\section{Introduction}\label{r:sec:introduction}

OpenAI's \emph{Finite Time Blowup for Navier--Stokes} proposes a
three-dimensional flow that starts from rest and develops unbounded
velocity under a smooth compactly supported force
\cite[Theorem~1.1]{openai-paper}. Its kinetic energy remains finite
because the large velocities occupy shrinking spatial regions.
A comparison based only on the size of the external force or the
kinetic energy cannot distinguish this concentration from a regular
flow. The relevant question is whether the proposed solution can
satisfy the stronger space--time bounds imposed by a regularity theorem
for its prescribed data.

Part~I relates the temporal derivatives of one unforced velocity field to
its spatial Sobolev energies. Part~II carries that construction into a new
forced history and retains every prescribed source row. Viscosity couples the
energy at order \(1/2\) to dissipation at order \(3/2\), while temporal
differentiation produces higher spatial orders together with differentiated
nonlinear and source terms. Both the explicit force work and the velocity's
nonlinear response remain in the comparison below.

The forced theorem proved in Part~II gives the full-interval solution used
here.  By Corollary~\ref{r:thm:supplied-forced}, finite collections of its mixed Sobolev norms bound
the velocity and all the terms in its momentum equation. The proposed
exterior has divergent positive contributions to precisely these
norms. Uniqueness identifies a proposed classical solution with the
regular solution for its data, so the bounds apply before any
regularity at the candidate's terminal time is assumed.

The source estimates distinguish contributions controlled by the
initial data and prescribed force from those that require a velocity
norm. This distinction yields quantitative sufficient conditions for
continuation, proved in the appendices together with the general-order
source expansions and product estimates.

\subsection{Force, pressure, and material acceleration}

We work on \(\R^3\), at fixed viscosity \(\nu>0\), with unit density
and no material boundary. Let \(u_0\) be divergence free and belong to
\(H^\infty=\bigcap_{m\ge0}H^m\). Fix \(0<T<\infty\). The force
\(f(x,t)\), measured per unit mass, is smooth through \([0,T]\), and
all its mixed derivatives decay rapidly in space, uniformly on this
interval. A subscript \(\sigma\) denotes the divergence-free subspace.
The whole-space Navier--Stokes system is
\cite[equations~(1)--(3)]{Fefferman2006}:
\begin{equation}
 u_t+(u\cdot\nabla)u+\nabla p-\nu\Delta u=f,
 \qquad \nabla\cdot u=0,\qquad u(0)=u_0.
 \label{r:eq:ns}
\end{equation}

At a fixed location, \(u_t\) records how the velocity changes with time.
A moving parcel also encounters spatial changes in velocity. Along
\(\dot X=u(X,t)\), its acceleration is
\begin{equation}
 \begin{aligned}
 \frac{\dd}{\dd t}u(X(t),t)&=D_tu(X(t),t),
 &D_t&=\partial_t+u\cdot\nabla,\\
 D_tu&=-\nabla p+\nu\Delta u+f.
 \end{aligned}
 \label{r:eq:parcel-law}
\end{equation}
Thus the prescribed force contributes to the acceleration alongside
pressure and viscous stress. Pressure depends on the surrounding velocity
through incompressibility. Viscosity depends on its spatial curvature.
These two responses must be retained when asking what a finite external
force can cause.

Eulerian and material derivatives describe different derivatives of
the velocity field along fixed points and moving trajectories. With \(V_n=\partial_t^nu\), acceleration is
\(D_tu=V_1+(V_0\cdot\nabla)V_0\). Every higher parcel derivative is
also a finite combination of the velocity's mixed derivatives.
At all finite orders the two derivative families determine each other
through product and commutator identities.
Appendix~\ref{r:app:endpoint} proves equality of the two complete Sobolev
intersections. Finite rectangles of that intersection bound both
the fixed-location derivatives and the parcel responses.

Differentiating the force along a parcel already exhibits this coupling:
\[
 D_tf=\partial_tf+(u\cdot\nabla)f.
\]
The spatial variation of a smooth prescribed force is sampled at a rate
set by the velocity. Higher material derivatives also differentiate that
velocity. The parcel's acceleration, jerk, and subsequent derivatives
are calculated within the mixed jet by exact product and commutator
identities. This relation also appears in the
Euler formulation \cite[Section~1; Proposition~6.1]{birnie-euler}.

Incompressibility determines the pressure contribution explicitly.
Let \(\Pp\) denote the orthogonal Leray projection, with Fourier symbol
\(I-\xi\otimes\xi/|\xi|^2\) for \(\xi\ne0\). Its value at zero
is immaterial to an \(L^2\)-based norm. Write \(\mathbb Q=I-\Pp\) for
the gradient projection. Then
\begin{equation}
 \begin{aligned}
 u_t-\nu\Delta u+\Pp\nabla\cdot(u\otimes u)&=\Pp f,\\
 \nabla p&=(I-\Pp)\bigl[f-(u\cdot\nabla)u\bigr].
 \end{aligned}
 \label{r:eq:projected}
\end{equation}
Pressure is determined up to a function of time; its gradient is the
physical quantity in momentum. A conservative force \(f=\nabla\phi\)
can be absorbed by setting \(q=p-\phi\). The velocity then obeys the
unforced equation exactly. In particular, rest remains a solution, with
\(p=\phi\). Boundary-driven flows require additional boundary conditions and are
outside the whole-space setting of this paper.


\subsection{The forced theorem and the comparison principle}

Part~II has already proved global smoothness and uniqueness for every datum
and force in the class used below. Its theorem gives the following
closed-interval consequence in the notation of this part.

\begin{corollary}[Finite-horizon forced intersection]
\label{r:thm:supplied-forced}
Let \(\nu>0\), \(u_0\in H^\infty_\sigma(\R^3)\), and
\(f\in C^\infty([0,\infty);\mathcal S(\R^3;\R^3))\). For every finite
\(T>0\), the unique global pair \((U,P)\) supplied by
Theorem~\ref{f:thm:ier-global-smoothness} satisfies
\begin{equation}
 U\in\bigcap_{r,m\ge0}H^r(0,T;H^m(\R^3)).
 \label{r:eq:direct-target}
\end{equation}
The mixed derivatives satisfy the differentiated momentum and
incompressibility equations.
\end{corollary}

\begin{proof}
Part~II gives the global pressure-normalized solution. On \([0,T]\), its
persistence argument places each spatial row in \(C([0,T];H^m)\), and the
complete forced recurrence produces every temporal row. Hence, for every
finite \(r,m\),
\[
 \sum_{n=0}^{r}\norm{\partial_t^nU}_{L^2(0,T;H^m)}^2
 \le T\sum_{n=0}^{r}
 \sup_{0\le t\le T}\norm{\partial_t^nU(t)}_{H^m}^2<\infty.
\]
Taking the intersection over the finite indices gives
\eqref{r:eq:direct-target}; the differentiated momentum and incompressibility
rows are restrictions of the complete forced mixed jet proved in Part~II.
\end{proof}

To see the comparison principle, let \(v\) be another classical solution
for these data on \([0,T)\), with finite spatial Sobolev norms on each
strict subinterval. Uniqueness gives \(v=U\) there. The adjacent norms
\(U\in L^2(0,T;H^3)\) and \(U_t\in L^2(0,T;H^1)\) imply
\begin{equation}
 \sup_{0\le t<T}\norm{v(t)}_\infty^2
 \le \frac{\norm{u_0}_{H^2}^2+
       \norm U_{L^2H^3}^2+\norm{U_t}_{L^2H^1}^2}{8\pi}<\infty.
 \label{r:eq:intro-comparison}
\end{equation}
The trace and embedding calculations are proved in
Section~\ref{r:sec:velocity}; the uniqueness argument is given in
Section~\ref{r:sec:exclusion}.  Before that argument is applied to the
candidate of \cite{openai-paper}, Section~\ref{r:sec:force-extension} proves
that its full localized residual belongs to the prescribed force class of
Part~II.  Once that source payment has been made, the estimate excludes
unbounded velocity. The
critical-height and material calculations below identify further
norms in which the proposed concentration fails the forced theorem.

\section{Forced mixed derivatives and Sobolev energies}\label{r:sec:forced-rows}

On each compact interval of local smooth existence, define
\(V_n=\partial_t^nu\) at fixed spatial coordinates. All spatial norms use
Euclidean vector and Frobenius tensor norms. Throughout the derivative formulas, \(f\) is the prescribed force in \eqref{r:eq:ns}.

Put \(F_n=\partial_t^nf\), \(Q_n=\partial_t^np\), and
\begin{equation}
 B_n=\sum_{a=0}^n\binom na(V_a\cdot\nabla)V_{n-a}.
 \label{r:eq:Bn}
\end{equation}
The product rule applied to \eqref{r:eq:ns} gives the forced recurrence already
derived in Proposition~\ref{f:prop:formal-mixed-recurrences}:
\begin{equation}
 V_{n+1}+B_n+\nabla Q_n-\nu\Delta V_n=F_n.
 \label{r:eq:row}
\end{equation}
Applying \(\partial_x^\alpha\) to transport gives the full mixed expansion
\begin{equation}
 \partial_x^\alpha B_n
 =\sum_{a=0}^n\sum_{\beta\le\alpha}
 \binom na\binom\alpha\beta
 (\partial_x^\beta V_a\cdot\nabla)
 \partial_x^{\alpha-\beta}V_{n-a}.
 \label{r:eq:mixed}
\end{equation}
Every force coordinate in this equation is the prescribed derivative
\(\partial_x^\alpha F_n\).

Taking divergence of \eqref{r:eq:row} gives the pressure Poisson equation.
In the Sobolev class used below, every nonpressure term of that row belongs
to \(H^m\), so \(\nabla Q_n\) does too. Applying the orthogonal
projection \(I-\Pp\) to the full row then fixes its gradient:
\begin{equation}
 -\Delta Q_n=\nabla\cdot B_n-\nabla\cdot F_n,
 \qquad \nabla Q_n=(I-\Pp)(F_n-B_n).
 \label{r:eq:pressure}
\end{equation}
The gradient part of the force contributes to pressure, while its solenoidal
part remains in the projected velocity equation. Formula \eqref{r:eq:pressure}
will bound pressure through its gradient, the quantity present in momentum.


The input regularity already gives a compatible source intersection:
\begin{equation}
 f\in\bigcap_{r,m\ge0}H^r(0,T;H^m(\R^3)).
 \label{r:eq:force-intersection}
\end{equation}
Indeed, each mixed derivative has finite spatial Sobolev norm uniformly on the
compact time interval, and integration in time preserves finiteness. The
initial velocity jet is generated recursively by \eqref{r:eq:row}, in the
form proved in Lemma~\ref{f:lem:formal-initial-jet}:
\begin{equation}
 a_0=u_0,\qquad
 a_{n+1}=\nu\Delta a_n
 -\Pp\sum_{b=0}^n\binom nb(a_b\cdot\nabla)a_{n-b}
 +\Pp F_n(0),\qquad a_n=V_n(0).
 \label{r:eq:initial}
\end{equation}
The product bound proved below shows by induction that \(a_n\in H^m\)
for every finite \(n,m\), when \(u_0\in H^\infty_\sigma\).

\subsection{Viscous coupling of Sobolev orders}

Set \(\Lambda=(-\Delta)^{1/2}\). Its power \(\Lambda^s\) weights a
Fourier component of wave number \(\xi\) by \(|\xi|^s\). Higher
Sobolev energies thus measure progressively finer spatial variation;
only height zero is the ordinary kinetic energy. For a temporal row
\(r\) and real height \(s\ge0\), define
\[
 E_{r,s}=\tfrac12\norm{\Lambda^sV_r}_2^2,\qquad
 \mathcal P_{r,s}=-\operatorname{Re}\ip{\Lambda^sV_r}{\Lambda^sB_r},
 \qquad
 W_{r,s}=\operatorname{Re}\ip{\Lambda^sV_r}{\Lambda^sF_r}.
\]
Pairing \eqref{r:eq:row} with \(\Lambda^{2s}V_r\) eliminates the pressure
gradient by incompressibility. The viscous pairing is
\(-\nu\norm{\Lambda^{s+1}V_r}_2^2\), giving the exact row balance
\begin{equation}
 \partial_tE_{r,s}=-2\nu E_{r,s+1}+\mathcal P_{r,s}+W_{r,s}.
 \label{r:eq:height}
\end{equation}
Differentiating repeatedly exposes the higher spatial heights:
\begin{equation}
 \partial_t^kE_{r,s}=(-2\nu)^kE_{r,s+k}
 +\sum_{j=0}^{k-1}(-2\nu)^j\partial_t^{k-1-j}
       (\mathcal P_{r,s+j}+W_{r,s+j}),\qquad k\ge1.
 \label{r:eq:completed}
\end{equation}
For \(k=1\) this is \eqref{r:eq:height}. Differentiating the formula for
\(k\), substituting \eqref{r:eq:height} at height \(s+k\), and collecting
the added term proves the formula for \(k+1\). Force leaves the viscous
coefficient of the exposed energy unchanged. At \(r=0,s=1/2\), the
successive spatial heights are \(1/2,3/2,5/2,\ldots\).
Here \(\partial_t^kE_{r,s}\) differentiates a scalar energy. It is not
\(E_{r+k,s}\), the energy of a higher velocity derivative. For example,
\[
 \partial_t^2E_{r,s}
 =\norm{\Lambda^sV_{r+1}}_2^2
 +\operatorname{Re}\ip{\Lambda^sV_r}{\Lambda^sV_{r+2}}.
\]
The mixed velocity rows organize both factors in this expression, and all
factors in the differentiated nonlinear and force terms.

The work derivatives retain both the response and source jets:
\begin{equation}
 \partial_t^\ell W_{r,s}
 =\sum_{b=0}^\ell\binom\ell b
 \operatorname{Re}\ip{\Lambda^sV_{r+b}}{\Lambda^sF_{r+\ell-b}}.
 \label{r:eq:work-jet}
\end{equation}
The nonlinear current has the full product rule
\begin{equation}
 \partial_t^\ell\mathcal P_{r,s}
 =-\sum_{a=0}^r\binom ra
 \sum_{b+c+d=\ell}\frac{\ell!}{b!c!d!}
 \operatorname{Re}\ip{\Lambda^sV_{r+b}}
 {\Lambda^s\bigl((V_{a+c}\cdot\nabla)V_{r-a+d}\bigr)}.
 \label{r:eq:production-jet}
\end{equation}
Spatial differentiation uses \eqref{r:eq:mixed} and the multi-index product
rule. Every source coordinate remains a derivative of the prescribed field
\(f\); the velocity factor in \eqref{r:eq:work-jet} remains part of the
response. The source intersection \eqref{r:eq:force-intersection} controls each source factor. Bounds on the velocity factors require coordinates of the velocity history.

These identities hold on every \([0,S]\) within smooth existence.
If the velocity belongs to the full-interval intersection, the continuous
limits in Appendix~\ref{r:app:endpoint} extend them to \(T\). All differentiated
nonlinear and force terms are retained before taking that limit.




\subsection{Forcing corrections at the critical heights}

At the critical starting height, put \(H_r=E_{r,1/2}\). Rearranging the
forced height identity gives the completed coordinate
\begin{equation}
 \begin{aligned}
 C^f_{r,k}&:=\frac{\partial_t^kH_r}{(-2\nu)^k}\\
 &\quad-\sum_{j=0}^{k-1}(-2\nu)^{j-k}\partial_t^{k-1-j}
 (\mathcal P_{r,j+1/2}+W_{r,j+1/2}),\\
 C^f_{r,k}&=E_{r,k+1/2}.
 \end{aligned}
 \label{r:eq:forced-completion}
\end{equation}
The factor \((-2\nu)^k\) is unchanged by forcing and is nonzero at every
finite order. Each next temporal differentiation thus exposes the spatial
height \(k+1/2\), with its complete nonlinear and source contributions.
The identity is exact on every interval of smooth existence.


At the first two exposed heights, the source correction is bounded
using kinetic energy and source regularity. Put \(W_s=W_{0,s}\) and
\(h_s=\Lambda^{2s}\Pp f\). Moving derivatives onto \(h_s\) writes the
work as \(W_s=\ip u{h_s}\). The energy inequality bounds
\(\norm u_{L^\infty L^2}\); differentiating this pairing moves the
transport and diffusion derivatives onto the smooth prescribed field.
Consequently \(W_s\) and \(W_s'\) have finite bounds
\(\mathsf B_s\) and \(\mathsf L_s\), defined in
\eqref{r:eq:source-work-constants} and proved in
Appendix~\ref{r:sec:source-primitives}.
Evaluate the completed-height formulas of
Proposition~\ref{f:prop:completed-height-recurrence} on the forced velocity
itself, and denote these expressions by
\begin{equation}
 \begin{aligned}
 \widehat C_1[u]
   &:=-\frac{H'-\mathcal P_{1/2}}{2\nu},\\
 \widehat C_2[u]
   &:=\frac{H''-\partial_t\mathcal P_{1/2}
                   +2\nu\mathcal P_{3/2}}{4\nu^2},
 \qquad H=E_{0,1/2}.
 \end{aligned}
 \label{r:eq:original-form-completions}
\end{equation}
Here \(\mathcal P_s=\mathcal P_{0,s}\) and \(W_s=W_{0,s}\).
Substituting the forced balance, with both time derivatives expanded,
gives
\begin{equation}
 \begin{aligned}
 \widehat C_1[u]
   &=E_{0,3/2}-\frac{W_{1/2}}{2\nu},\\
 \widehat C_2[u]
   &=E_{0,5/2}
       +\frac{W_{1/2}'-2\nu W_{3/2}}{4\nu^2}.
 \end{aligned}
 \label{r:eq:low-completion-source-shift}
\end{equation}
The estimates \eqref{r:eq:source-work-constants} control every source
expression on the right. With
\[
 \delta_1:=\frac{\mathsf B_{1/2}}{2\nu},
 \qquad
 \delta_2:=\frac{\mathsf L_{1/2}}{4\nu^2}
              +\frac{\mathsf B_{3/2}}{2\nu},
\]
we obtain, for \(j=1,2\),
\begin{equation}
 \norm{C^f_{0,j}-\widehat C_j[u]}_{L^\infty(0,T)}
 \le\delta_j<\infty.
 \label{r:eq:low-completion-uniform-difference}
\end{equation}
These constants depend on the prescribed force and the kinetic-energy
bound; no smallness condition is imposed. On a finite interval a bounded
additive correction belongs to every \(L^p\), \(1\le p\le\infty\).
Consequently, the exact membership relation is
\begin{equation}
 \begin{aligned}
 C^f_{0,j}\in L^p(0,T)
   &\quad\Longleftrightarrow\quad
     \widehat C_j[u]\in L^p(0,T),\\
 \norm{C^f_{0,j}}_{L^p(0,T)}
   &\le\norm{\widehat C_j[u]}_{L^p(0,T)}
            +T^{1/p}\delta_j,
 \end{aligned}
 \label{r:eq:low-completion-space-preservation}
\end{equation}
where \(T^{1/\infty}=1\). Thus the explicit force correction preserves
these completed-coordinate spaces, including the unweighted time integral
at the critical viscous height \(3/2\). Every expression here is evaluated
on the forced history. The notation \(\widehat C_j[u]\) retains the earlier
completion formula; it does not replace its argument by an unforced
solution.

\section{Finite rectangles and uniform Sobolev bounds}\label{r:sec:rectangles}

Assume that the forced solution belongs to the class
\begin{equation*}
 u\in\mathscr I_T:=\bigcap_{r,m\ge0}H^r(0,T;H^m(\R^3)).
 \tag{I}\label{r:eq:I}
\end{equation*}
For integers \(N,M\ge0\), collect two adjacent space--time norms in
\begin{equation}
 R_{N,M}:=\sum_{n=0}^N
 \left(\norm{V_n}_{L^2H^{M+1}}^2+
 \norm{V_{n+1}}_{L^2H^{M-1}}^2\right)<\infty.
\label{r:eq:rectangle}
\end{equation}
This is the velocity part of the finite rectangle in
Definition~\ref{f:def:finite-rectangle}; its organization by restriction is
given by Definition~\ref{f:def:compatible-intersection} and
Theorem~\ref{f:thm:compatible-inverse-intersection}.
Time norms in this paper are over \((0,T)\). Different finite selections
agree under restriction because their entries are derivatives of \(u\).
There is no bound common to all \(N,M\) in the hypothesis.
For \(M=0\), \(H^{M-1}=H^{-1}\) is the usual dual Sobolev space.

We first convert these integral bounds into uniform bounds on a row. For
\(v\in L^2H^{m+1}\) with \(v_t\in L^2H^{m-1}\), the Hilbert-pair
identity is
\[
 \frac{\dd}{\dd t}\norm v_{H^m}^2
 =2\operatorname{Re}\ip{(1-\Delta)^{(m+1)/2}v}
 {(1-\Delta)^{(m-1)/2}v_t}.
\]
For classical rows it follows by differentiation and Plancherel. The identity
extends to the displayed Hilbert spaces by spatial Fourier truncation and
time smoothing. To justify the continuous representative, apply the resulting
estimate to the difference of two truncated fields: their spatial
\(L^2H^{m+1}\) tails and temporal \(L^2H^{m-1}\) tails tend to zero.
The truncations are consequently Cauchy in \(C([0,T];H^m)\). Integrate
the identity in either direction from a time whose squared norm is at most
its average. Cauchy--Schwarz gives
\begin{equation}
 \sup_{0\le t\le T}\norm v_{H^m}^2
 \le T^{-1}\norm v_{L^2H^m}^2
 +2\norm v_{L^2H^{m+1}}\norm{v_t}_{L^2H^{m-1}}.
 \label{r:eq:trace}
\end{equation}
Define the finite row constants
\begin{equation}
 A_{n,m}:=\sqrt{(1+T^{-1})R_{n,m}},
 \qquad L_{n,m}:=\sup_{0\le t\le T}\norm{F_n(t)}_{H^m}.
 \label{r:eq:constants}
\end{equation}
The constants \(A_{n,m}\) are distinguished from the initial fields
\(a_n\) in \eqref{r:eq:initial}.
Using \(2ab\le a^2+b^2\) in \eqref{r:eq:trace} yields
\begin{equation}
 \sup_{0\le t\le T}\norm{V_n(t)}_{H^m}\le A_{n,m}.
 \label{r:eq:uniform-row}
\end{equation}
The estimate includes the endpoint through the continuous trace.
Its constant depends on the indicated finite rectangle.

\section{Critical production and velocity bounds}\label{r:sec:critical}

The two heights \(1/2\) and \(3/2\) enter together because viscosity
pairs the critical velocity energy with one additional spatial derivative.
We now extract both bounds from one finite rectangle and substitute them
into the nonlinear current. Select \(R=R_{0,1}\), and write \(D=\norm{u_0}_{H^1}^2\).
Integrating the trace identity from time zero gives
\[
 \sup_t\norm u_{H^1}^2\le D+R,
 \qquad \norm u_{L^2H^2}^2\le R.
\]
Since positive homogeneous derivative norms are bounded by these
inhomogeneous norms,
\begin{equation}
 \sup_t\norm u_{\dot H^{1/2}}\le\sqrt{D+R},
 \qquad \norm u_{L^2\dot H^{3/2}}^2\le R.
 \label{r:eq:critical-input}
\end{equation}
The sharp fractional Sobolev inequality gives
\[
 \norm v_3\le C_3\norm{\Lambda^{1/2}v}_2,\qquad
 C_3=2^{-1/6}\pi^{-1/3}
\]
\cite{Lieb1983,frank-lieb}.
In dimension three at order \(1/2\), the squared form has coefficient
\(2^{1/3}\pi^{2/3}\) on \(\|v\|_3^2\). The constant also applies
to Euclidean vectors and Frobenius tensors: the positive scalar kernel
of \(\Lambda^{-1/2}\) gives
\(|\Lambda^{-1/2}F|\le\Lambda^{-1/2}|F|\), followed by the scalar
inequality. This introduces no factor for the number of components.
H\"older applied to
\(\mathcal P_{1/2}=-\ip{\Lambda u}{(u\cdot\nabla)u}\)
and this embedding on its three factors gives
\[
 |\mathcal P_{1/2}|
 \le C_3^3\norm u_{\dot H^{1/2}}
 \norm u_{\dot H^{3/2}}^2.
\]
Substitution of \eqref{r:eq:critical-input} yields
\begin{equation}
 \int_0^T|\mathcal P_{1/2}|\dd t
 \le \frac{\sqrt{D+R}\,R}{\sqrt2\pi}<\infty.
 \label{r:eq:critical-production}
\end{equation}
For rest data on \((0,1)\), this reads
\[
 \sup_t E_{1/2}(t)\le R/2,\qquad
 \int_0^1 E_{3/2}(t)\dd t\le R/2,\qquad
 \int_0^1|\mathcal P_{1/2}|\dd t\le\frac{R^{3/2}}{\sqrt2\pi}.
\]
The first factor in production is uniformly bounded, and its two higher
spatial factors have a finite time integral. No comparison between the
size of this bound and viscosity is needed.

The forcing contribution is
\(W_{1/2}=\ip{\Lambda^{-1/2}f}{\Lambda^{3/2}u}\), so
\begin{equation}
 \int_0^T|W_{1/2}|\dd t
 \le\norm f_{L^2\dot H^{-1/2}}\sqrt R<\infty.
 \label{r:eq:critical-work}
\end{equation}
For completeness, the source's negative homogeneous norm is finite at low
frequencies as well as high frequencies. Splitting the Fourier integral at
\(|\xi|=1\), using \(\norm{\widehat f}_\infty\le
(2\pi)^{-3/2}\norm f_1\), gives
\[
 \norm{\Lambda^{-1/2}f}_2^2
 \le\frac1{4\pi^2}\norm f_1^2+\norm f_2^2.
\]
Both terms are uniformly finite under the source assumptions.
The dual pairing places the velocity at order \(3/2\), the order
generated by viscous dissipation from the critical energy at \(1/2\).
Thus forcing enters the critical pair used in Part~II:
source regularity controls one factor, and the rectangle controls
the velocity factor. The higher differentiated contributions remain
those in \eqref{r:eq:work-jet} and \eqref{r:eq:production-jet}.
Substituting the two integrable currents into the critical balance gives
\[
 \frac12\norm{\Lambda^{1/2}u(t)}_2^2
 +\nu\int_0^t\norm{\Lambda^{3/2}u}_2^2\dd s
 =\frac12\norm{\Lambda^{1/2}u_0}_2^2
 +\int_0^t(\mathcal P_{1/2}+W_{1/2})\dd s.
\]

Write \(F_-=\norm f_{L^2\dot H^{-1/2}}\). Substituting
\eqref{r:eq:critical-production} and \eqref{r:eq:critical-work} into this
identity bounds the complete forced critical balance:
\begin{equation}
 \begin{aligned}
 E_{1/2}(t)+2\nu\int_0^t E_{3/2}(s)\dd s
 &\le E_{1/2}(0)\\
 &\quad+\frac{\sqrt{D+R}\,R}{\sqrt2\pi}
       +F_-\sqrt R.
 \end{aligned}
 \label{r:eq:forced-critical-bound}
\end{equation}
For rest data, \(D=E_{1/2}(0)=0\), so the right side is
\(R^{3/2}/(\sqrt2\pi)+F_-\sqrt R\).
Here \(R\) bounds the two critical norms independently of the
production estimate. The calculation requires no absorption of
nonlinear production or source work into viscous dissipation.

\subsection{A uniform velocity bound}\label{r:sec:velocity}

The intersection also contains \(R_{0,2}<\infty\). This finite selection
already controls the speed of the forced fluid, before any estimate of
the individual momentum terms is needed. Put
\[
 a=\norm u_{L^2H^3},\qquad
 b=\norm{u_t}_{L^2H^1},\qquad R_{0,2}=a^2+b^2.
\]
Integrating the Hilbert-pair identity at \(m=2\) from the prescribed
initial velocity gives
\begin{equation}
 \begin{aligned}
 \norm{u(t)}_{H^2}^2
 &=\norm{u_0}_{H^2}^2\\
 &\quad+2\operatorname{Re}\int_0^t
 \ip{(1-\Delta)^{3/2}u}
 {(1-\Delta)^{1/2}u_t}\dd s,\\
 \sup_{0\le t\le T}\norm{u(t)}_{H^2}^2
 &\le\norm{u_0}_{H^2}^2+2ab\\
 &\le\norm{u_0}_{H^2}^2+R_{0,2}.
 \end{aligned}
 \label{r:eq:velocity-H2}
\end{equation}
Use the unitary Fourier transform and
\(\norm v_{H^s}^2=\int\langle\xi\rangle^{2s}|\widehat v(\xi)|^2\dd\xi\).
Since
\[
 \int_{\R^3}(1+|\xi|^2)^{-2}\dd\xi=\pi^2,
\]
Fourier inversion and Cauchy--Schwarz give
\begin{equation}
 \norm{\partial_x^\alpha v}_\infty
 \le c_\infty\norm v_{H^{|\alpha|+2}},
 \qquad c_\infty=\frac1{\sqrt{8\pi}}.
 \label{r:eq:embedding}
\end{equation}
This holds for vector fields with their Euclidean norm.

Applying \eqref{r:eq:embedding} to the velocity yields
\begin{equation}
 \sup_{0\le t\le T}\norm{u(t)}_\infty^2
 \le\frac{\norm{u_0}_{H^2}^2+R_{0,2}}{8\pi}.
 \label{r:eq:velocity-bound}
\end{equation}
In particular, for a fluid initially at rest,
\[
 \sup_{0\le t\le T}\norm{u(t)}_\infty^2
 \le\frac{R_{0,2}}{8\pi}<\infty.
\]
The endpoint belongs to the continuous \(H^2\) trace established in the intersection trace calculation. Thus the intersection supplies a uniform velocity bound
through the entire finite interval. The additional estimates below
resolve how acceleration, transport, pressure, and viscosity behave
individually within that velocity history.


\subsection{Higher spatial energies from the critical coordinate}
\label{r:sec:higher-spatial-source}

The critical coordinate also controls the growth of every higher spatial
energy. The homogeneous commutator inequality proved in
Lemma~\ref{f:lem:app-critical-commutator} is
\begin{equation}
 |\mathcal P_m|
 \le K_m\norm{\Lambda^{3/2}u}_2
          \norm{\Lambda^{m+1}u}_2\norm{\Lambda^mu}_2,
 \qquad m\ge1.
 \label{r:eq:actual-spatial-production}
\end{equation}
Here \(m\) is an integer and \(K_m\) is finite for each fixed order.
The inequality applies to the complete spatially differentiated energy
row, with its multinomial weights. It holds for a divergence-free field
independently of the force: the pressure gradient remains orthogonal to
the energy row, and the convection expression is unchanged.

Write
\[
 X_m=\norm{\Lambda^mu}_2^2,\qquad
 Y_m=\norm{\Lambda^{m+1}u}_2^2,\qquad
 a=\norm{\Lambda^{3/2}u}_2,\qquad
 b_m=\norm{\Lambda^{m-1}\Pp f}_2.
\]
The forced balance is
\[
 \frac12X_m'+\nu Y_m=\mathcal P_m+W_m,
 \qquad
 W_m=\ip{\Lambda^{m+1}u}{\Lambda^{m-1}\Pp f}.
\]
Young's inequality allocates a quarter of the dissipation to each pairing:
\begin{align}
 |\mathcal P_m|
 &\le\frac{\nu}{4}Y_m+\frac{K_m^2}{\nu}a^2X_m,
 \label{r:eq:actual-nonlinear-absorption}\\
 |W_m|&\le\frac{\nu}{4}Y_m+\frac1\nu b_m^2.
 \label{r:eq:actual-source-absorption}
\end{align}
Their substitution gives
\begin{equation}
 X_m'+\nu Y_m
 \le\frac{2K_m^2}{\nu}a^2X_m+\frac2\nu b_m^2.
 \label{r:eq:actual-forced-spatial-row}
\end{equation}
The current energy enters linearly. Its coefficient has a finite time
integral, since \eqref{r:eq:critical-input} gives
\(\int_0^T a(t)^2\dd t\le R\).

Set
\[
 A_m(t)=\frac{2K_m^2}{\nu}\int_0^t a(s)^2\dd s,\qquad
 \Phi_m=\int_0^T b_m(t)^2\dd t.
\]
Multiplying \eqref{r:eq:actual-forced-spatial-row} by \(e^{-A_m(t)}\)
and integrating gives
\[
 e^{-A_m(t)}X_m(t)+\nu\int_0^t e^{-A_m(s)}Y_m(s)\dd s
 \le X_m(0)+\frac2\nu\int_0^t e^{-A_m(s)}b_m(s)^2\dd s.
\]
Since \(A_m\) is nondecreasing and \(A_m(T)\le2K_m^2R/\nu\),
\begin{equation}
 \begin{aligned}
 X_m(t)+\nu\int_0^tY_m(s)\dd s
 &\le C_m^f(R),\\
 C_m^f(R)
 &:=\exp\!\left(\frac{2K_m^2R}{\nu}\right)
 \left[\norm{\Lambda^mu_0}_2^2+\frac2\nu\Phi_m\right].
 \end{aligned}
 \label{r:eq:actual-forced-spatial-bound}
\end{equation}
The source assumptions give \(\Phi_m<\infty\), with no restriction on its
size. Taking the supremum of the energy and the terminal limit of the
dissipation separately bounds each by \(C_m^f(R)\). The ordinary energy
bound supplies the low frequencies. Higher temporal rows remain derivatives
of this velocity and are generated by \eqref{r:eq:row} from these
spatial bounds and the prescribed source jets.

The factor controlling growth is the time integral of
\(\norm{\Lambda^{3/2}u}_2^2\). Once this critical quantity is finite,
the source contributes the finite additive term \(\Phi_m\), and
Gronwall's inequality bounds every fixed higher spatial order.

\section{Bounds on each momentum term}
\label{r:sec:individual}

A bounded sum in the momentum equation does not bound its summands:
large inertial, pressure, and viscous terms can cancel. The mixed
Sobolev estimates give stronger information. By selecting finitely
many additional derivatives of \(u\), we can bound each term before
forming the sum.

\subsection{Spatial product bounds}

For \(s\ge1\),
\[
 \langle\xi\rangle^s\le
 2^{s-1}\left(\langle\eta\rangle^s+
 \langle\xi-\eta\rangle^s\right).
\]
Apply this to the convolution formula for \(\widehat{a\otimes b}\),
use Young's convolution inequality, and use
\(\norm{\widehat a}_1\le\pi\norm a_{H^2}\). The result is
\begin{equation}
 \norm{a\otimes b}_{H^s}
 \le 2^{s-1}c_\infty
 \left(\norm a_{H^s}\norm b_{H^2}+
 \norm a_{H^2}\norm b_{H^s}\right).
 \label{r:eq:product}
\end{equation}
For tensors use the Frobenius norm. Divergence is bounded from \(H^{m+1}\)
to \(H^m\) with constant one. If \(\nabla\cdot a=0\), transport is
the divergence of the appropriate tensor product. Thus
\begin{equation}
 \norm{(a\cdot\nabla)b}_{H^m}
 \le C_m\norm a_{H^{m+2}}\norm b_{H^{m+2}},
 \qquad C_m=2^{m+1}c_\infty,\quad m\ge0.
 \label{r:eq:transport-product}
\end{equation}

\subsection{The complete momentum row}

For each temporal row and Sobolev height, define
\begin{equation}
 S_{n,m}:=C_m\sum_{a=0}^n\binom na
 A_{a,m+2}A_{n-a,m+2}.
 \label{r:eq:S}
\end{equation}
The finite sum retains both transport factors in every binomial summand.

\begin{theorem}[Individual mixed-term bounds]\label{r:thm:individual}
Assume \eqref{r:eq:ns}, the source regularity above, and \eqref{r:eq:I}.
For every \(n,m\ge0\), each term in the differentiated momentum equation
satisfies
\begin{align}
 \sup_t\norm{V_{n+1}(t)}_{H^m}&\le A_{n+1,m},\label{r:eq:time-bound}\\
 \sup_t\norm{B_n(t)}_{H^m}&\le S_{n,m},\label{r:eq:transport-bound}\\
 \sup_t\norm{\nabla Q_n(t)}_{H^m}&\le L_{n,m}+S_{n,m},\label{r:eq:pressure-bound}\\
 \sup_t\norm{\nu\Delta V_n(t)}_{H^m}&\le\nu A_{n,m+2}.\label{r:eq:visc-bound}
\end{align}
For \(|\alpha|=k\), their mixed spatial derivatives satisfy
\begin{align}
 \norm{\partial_x^\alpha V_{n+1}}_{L^\infty_{t,x}}
 &\le c_\infty A_{n+1,k+2},\label{r:eq:time-point}\\
 \norm{\partial_x^\alpha B_n}_{L^\infty_{t,x}}
 &\le c_\infty S_{n,k+2},\label{r:eq:transport-point}\\
 \norm{\partial_x^\alpha\nabla Q_n}_{L^\infty_{t,x}}
 &\le c_\infty(L_{n,k+2}+S_{n,k+2}),\label{r:eq:pressure-point}\\
 \norm{\nu\partial_x^\alpha\Delta V_n}_{L^\infty_{t,x}}
 &\le\nu c_\infty A_{n,k+4}.\label{r:eq:visc-point}
\end{align}
Every constant on the right is finite and comes from finitely many velocity
intersection coordinates and prescribed force coordinates.
\end{theorem}

\begin{proof}
The first bound is \eqref{r:eq:uniform-row} for temporal row \(n+1\).
Apply \eqref{r:eq:transport-product} to each summand of \eqref{r:eq:Bn} and
then apply \eqref{r:eq:uniform-row}; this gives \eqref{r:eq:transport-bound}.
The Fourier multiplier \(I-\Pp\) has operator norm at most one in every
\(H^m\). Its application to \eqref{r:eq:pressure} proves
\eqref{r:eq:pressure-bound}, including the forcing contribution.
The multiplier bound \(|\xi|^2\le\langle\xi\rangle^2\) proves
\eqref{r:eq:visc-bound}. Applying \eqref{r:eq:embedding} to these four
fields gives \eqref{r:eq:time-point}--\eqref{r:eq:visc-point}.
\end{proof}

For the undifferentiated momentum equation, the calculation is particularly
concrete. The four individual bounds are
\begin{align*}
 \norm{u_t}_{L^\infty_{t,x}}&\le c_\infty A_{1,2},\\
 \norm{(u\cdot\nabla)u}_{L^\infty_{t,x}}
 &\le c_\infty C_2 A_{0,4}^2=\frac{A_{0,4}^2}{\pi},\\
 \norm{\nabla p}_{L^\infty_{t,x}}
 &\le c_\infty L_{0,2}+\frac{A_{0,4}^2}{\pi},\\
 \norm{\nu\Delta u}_{L^\infty_{t,x}}
 &\le\nu c_\infty A_{0,4}.
\end{align*}
Here \(C_2=8c_\infty\). These estimates bound each summand before the
summands are added. The signs in \eqref{r:eq:ns} play no role in their
finiteness.


\subsection{Pressure completion and material acceleration}
\label{r:sec:orthogonal-acceleration}

Pressure reconstruction separates the material acceleration into
orthogonal solenoidal and gradient components. Let
\(\mathbb Q=I-\Pp\), let \(B_0=(u\cdot\nabla)u\), and write
\(a=D_tu\). Pressure reconstruction in \eqref{r:eq:projected} gives
\begin{equation}
 \begin{aligned}
 \nabla p&=\mathbb Q(f-B_0),\\
 a&=\nu\Delta u+\Pp f+\mathbb Q B_0,\\
 \Pp a&=\nu\Delta u+\Pp f,\qquad
 \mathbb Q a=\mathbb Q B_0.
 \end{aligned}
 \label{r:eq:acceleration-projections}
\end{equation}
The pressure response removes the gradient component of the applied
force and fixes the gradient component of material acceleration through
the velocity. The viscous contribution is divergence free. In every
Fourier Sobolev inner product, these two components are orthogonal:
\begin{equation}
 \norm a_{H^s}^2
 =\norm{\nu\Delta u+\Pp f}_{H^s}^2
  +\norm{\mathbb Q B_0}_{H^s}^2.
 \label{r:eq:acceleration-orthogonal}
\end{equation}
This identity also holds for homogeneous norms wherever the displayed
quantities are finite. For \(f=0\), it reads
\[
 \norm{D_tu}_{\dot H^s}^2
 =\nu^2\norm u_{\dot H^{s+2}}^2
  +\norm{\nabla p}_{\dot H^s}^2.
\]
Pressure cannot cancel the viscous contribution in this global norm.
With forcing retained, cancellation inside the divergence-free component
is governed by the complete sum \(\nu\Delta u+\Pp f\).

The critical rectangle gives a quantitative bound for these components.
Put \(h=\norm u_{\dot H^{1/2}}\), \(d=\norm u_{\dot H^{3/2}}\), and
\(c_*=(\sqrt2\pi)^{-1}\). The Sobolev inequality used in
the critical production calculation and its dual give
\[
 \norm{B_0}_{\dot H^{-1/2}}
 \le C_3\norm u_3\norm{\nabla u}_3
 \le c_*hd.
\]
The projector \(\mathbb Q\) has norm one, so this also bounds
\(\mathbb Q B_0\). There is a stronger spatial estimate for that
gradient component. Expanding its divergence and using
\(\nabla\cdot u=0\) puts one derivative on each factor:
\[
 \operatorname{div}B_0=\partial_i u_j\,\partial_j u_i.
\]
The Fourier symbol of \(\mathbb Q\) and the dual Sobolev inequality
now imply
\begin{equation}
 \begin{aligned}
 \norm{\mathbb Q B_0}_{\dot H^{1/2}}
 &=\norm{\operatorname{div}B_0}_{\dot H^{-1/2}}\\
 &\le C_3\norm{\operatorname{tr}[(\nabla u)^2]}_{3/2}\\
 &\le C_3\norm{\nabla u}_3^2
 \le c_*d^2.
 \end{aligned}
 \label{r:eq:pressure-critical-gain}
\end{equation}
This estimate controls the nonlinear pressure gradient
\(\nabla p-\mathbb Q f=-\mathbb Q B_0\). It retains the prescribed
source part of pressure separately.

Use the already selected \(R=R_{0,1}\),
\(D=\norm{u_0}_{H^1}^2\), and set
\(F_-^\sigma=\norm{\Pp f}_{L^2\dot H^{-1/2}}\le F_-\).
Substituting \eqref{r:eq:critical-input} into
\eqref{r:eq:acceleration-orthogonal} and
\eqref{r:eq:pressure-critical-gain} yields
\begin{equation}
 \begin{aligned}
 \norm{D_tu}_{L^2\dot H^{-1/2}}^2
 &\le\bigl(\nu\sqrt R+F_-^\sigma\bigr)^2
       +\frac{(D+R)R}{2\pi^2},\\
 \norm{\nabla p-\mathbb Q f}_{L^1\dot H^{1/2}}
 &\le\frac{R}{\sqrt2\pi}.
 \end{aligned}
 \label{r:eq:critical-material-bounds}
\end{equation}
These are finite consequences of the two critical coordinates.
The first bound uses orthogonality to add the squared component norms;
the second uses incompressibility to control the pressure contribution
at a higher spatial order. The pointwise acceleration bound below uses
a stronger finite rectangle.

Ordinary temporal differentiation preserves the two projections:
\[
 \partial_t^na
 =\nu\Delta V_n+\Pp F_n+\mathbb Q B_n.
\]
The orthogonal identity follows at each Eulerian temporal row. Repeated
material differentiation also differentiates the transport operator;
the commutators retained in Appendix~\ref{r:app:endpoint} give those higher
parcel derivatives.


\subsection{Equivalent critical continuation conditions}
\label{r:sec:terminal-characterization}

The projected acceleration and critical dissipation are linked by
an exact identity. Combined with the spatial energy estimate, it
characterizes continuation through \(T\) in terms of several
equivalent quantities.
Define
\[
 \mathcal B_T=\int_0^T\norm u_{\dot H^{3/2}}^2\dd t.
\]
Projection of the material momentum law gives the exact identity
\begin{equation}
 \mathcal B_T
 =\nu^{-2}\norm{\Pp(D_tu-f)}_{L^2(0,T;\dot H^{-1/2})}^2.
 \label{r:eq:projected-response-identity}
\end{equation}
The identity follows from \(\Pp(D_tu-f)=\nu\Delta u\) and Plancherel.
It holds on each strict interval and passes to the nonnegative terminal
integrals, including the value \(+\infty\).

For the stated smooth datum and source, the following assertions about
a classical forced history on \([0,T)\) are equivalent:
\begin{equation}
 \begin{gathered}
 u\in\mathscr I_T;\qquad
 \mathcal B_T<\infty;\qquad
 \Pp D_tu\in L^2(0,T;\dot H^{-1/2});\\
 \widehat C_1[u]\in L^1(0,T);\qquad
 \sup_{t<T}\int_0^t\mathcal P_{1/2}(s)\dd s<\infty.
 \end{gathered}
 \label{r:eq:terminal-equivalent-statements}
\end{equation}
Here \(\widehat C_1\) is the original completion formula evaluated on
this forced velocity, as in \eqref{r:eq:original-form-completions}.
The source has a finite \(L^2\dot H^{-1/2}\) norm, so subtraction of
\(\Pp f\) preserves the projected-acceleration condition.
Equation \eqref{r:eq:low-completion-source-shift} identifies the
completed-coordinate condition with \(\mathcal B_T<\infty\).

For the remaining implications, first suppose \(\mathcal B_T<\infty\).
At spatial height one, differentiation and incompressibility give
\[
 |\mathcal P_1|
 \le\norm{\nabla u}_3\norm{\nabla u}_2\norm{\nabla u}_6
 \le K_1\norm u_{\dot H^{3/2}}\norm{\nabla u}_2\norm{\Delta u}_2.
\]
Young's inequality in the forced energy balance gives
\[
 \frac{\dd}{\dd t}\norm{\nabla u}_2^2+\nu\norm{\Delta u}_2^2
 \le\frac{2K_1^2}{\nu}\norm u_{\dot H^{3/2}}^2
               \norm{\nabla u}_2^2
       +\frac2\nu\norm{\Pp f}_2^2.
\]
Thus Gronwall bounds the \(H^1\) norm uniformly through \(T\).
Local classical existence in \(H^1\), with the prescribed smooth source,
has a lifespan bounded below by its \(H^1\) bound and a source bound.
Restarting sufficiently near \(T\) extends the history past \(T\).
Uniqueness identifies that extension on its overlap.
Higher spatial differentiation and the momentum recurrence then supply
every fixed spatial and temporal Sobolev norm on the closed horizon.
Equivalently, the explicit higher-energy estimates of
Section~\ref{r:sec:higher-spatial-source} perform this reconstruction.
This proves \(u\in\mathscr I_T\).
The reverse implication follows by selecting the finite spatial coordinate.

Finally, the work \(W_{1/2}=\ip u{\Lambda\Pp f}\) is uniformly bounded
by kinetic energy and a prescribed source norm. The exact critical balance
is
\[
 E_{1/2}(t)+\nu\int_0^t\norm u_{\dot H^{3/2}}^2\dd s
 =E_{1/2}(0)+\int_0^t\mathcal P_{1/2}\dd s
                  +\int_0^tW_{1/2}\dd s.
\]
A uniform upper bound for the signed production integral controls both
nonnegative terms on the left, including \(\mathcal B_T\).
Conversely, membership in \(\mathscr I_T\) gives the absolute production
bound \eqref{r:eq:critical-production}. This proves all the equivalences.

These equivalences characterize full-interval regularity of a given forced
solution. For the global solution of Part~II,
Corollary~\ref{r:thm:supplied-forced} supplies all these conditions and every
finite-rectangle bound.

% Source: direct-forced-construction-20260910.md, section 8, equations (37)--(43).
% Pressure normalization: section 2, equation (5), of that source.
\subsection{From spatial-height integrals to the forced mixed intersection}
\label{r:sec:forced-domain-propagation}

Let \((u,p)\) be the classical forced history on every strict interval
\([0,S]\), \(S<T<\infty\), with fixed \(\nu>0\),
\(u_0\in H^\infty_\sigma(\R^3)\), and the prescribed smooth force whose
mixed derivatives decay rapidly in space uniformly through \([0,T]\).
Assume explicitly that every zeroth-row completed height has its
whole-interval integral:
\begin{equation}
 \int_0^T C^f_{0,k}(t)\dd t
 =\frac12\int_0^T\norm{\Lambda^{k+1/2}u(t)}_2^2\dd t<\infty,
 \qquad k=0,1,2,\ldots .
 \label{r:eq:forced-spatial-domain-hypothesis}
\end{equation}
The forced momentum equation carries this spatial information to the
temporal rows. The first step is to recover the low frequencies omitted
by the homogeneous heights. Put
\[
 K_T=\norm{u_0}_2+\int_0^T\norm{f(t)}_2\dd t,
 \qquad \mathsf S_m=\norm{u}_{L^2(0,T;H^m)}.
\]
The kinetic bound gives \(\sup_{t<T}\norm{u(t)}_2\le K_T\).
For each integer \(m\ge0\), split the Fourier integral at \(|\xi|=1\).
Below that radius use the kinetic bound; above it use
\((1+|\xi|^2)^m\le2^m|\xi|^{2m+1}\). This gives
\begin{equation}
 \mathsf S_m^2
 \le2^mTK_T^2+2^{m+1}\int_0^T C^f_{0,m}(t)\dd t<\infty.
 \label{r:eq:forced-domain-spatial-integrals}
\end{equation}

These time integrals bound the total change in each spatial norm.
Indeed, the projected equation is
\[
 u_t=\nu\Delta u-\Pp[(u\cdot\nabla)u]+\Pp f.
\]
The transport estimate \eqref{r:eq:transport-product} places its nonlinear
term in \(L^1(0,T;H^m)\), with norm at most
\(C_m\mathsf S_{m+2}^2\). Its viscous term has \(L^1H^m\) norm at most
\(\nu\sqrt T\,\mathsf S_{m+2}\). Integrating the equation from the
initial datum now yields
\begin{equation}
 \begin{aligned}
 \sup_{t<T}\norm{u(t)}_{H^m}
 &\le\mathsf U_{0,m},\\
 \mathsf U_{0,m}
 &:=\norm{u_0}_{H^m}+\nu\sqrt T\,\mathsf S_{m+2}
       +C_m\mathsf S_{m+2}^2
       +\norm{\Pp f}_{L^1(0,T;H^m)}<\infty.
 \end{aligned}
 \label{r:eq:forced-domain-base-row}
\end{equation}
The integral of this \(L^1H^m\) right side supplies the continuous
representative, so the bound does not assume a terminal trace.

Every finite spatial order is now available in the base row.
Differentiating the forced equation retains every transport product:
\begin{equation}
 \begin{aligned}
 V_{n+1}&=\nu\Delta V_n-\Pp B_n+\Pp F_n,\\
 B_n&=\sum_{a=0}^n\binom na(V_a\cdot\nabla)V_{n-a},
 \qquad V_n=\partial_t^nu,\quad F_n=\partial_t^nf.
 \end{aligned}
 \label{r:eq:forced-domain-row-recurrence}
\end{equation}
Suppose that rows \(0,\ldots,n\) have uniform bounds at every finite
spatial order. Applying \eqref{r:eq:transport-product} term by term in
this finite sum gives the next constants:
\begin{equation}
 \begin{aligned}
 \mathsf U_{n+1,m}:={}&\nu\mathsf U_{n,m+2}\\
 &+C_m\sum_{a=0}^n\binom na
       \mathsf U_{a,m+2}\mathsf U_{n-a,m+2}
       +\norm{\Pp F_n}_{L^\infty(0,T;H^m)}.
 \end{aligned}
 \label{r:eq:forced-domain-row-bounds}
\end{equation}
The prescribed force makes the last term finite. Induction on \(n\)
proves
\(\sup_{t<T}\norm{V_n(t)}_{H^m}\le\mathsf U_{n,m}<\infty\)
for every fixed pair \(n,m\). At each step only finitely many higher
spatial orders occur, and each of them was already available in the
earlier rows. Integrating these uniform bounds gives
\begin{equation}
 \sum_{n=0}^r\norm{V_n}_{L^2(0,T;H^m)}^2
 \le T\sum_{n=0}^r\mathsf U_{n,m}^2<\infty,
 \qquad
 u\in\mathscr I_T=\bigcap_{r,m\ge0}H^r(0,T;H^m).
 \label{r:eq:forced-domain-mixed-intersection}
\end{equation}
Conversely, choosing spatial order \(k+1\) in this intersection gives
\eqref{r:eq:forced-spatial-domain-hypothesis}. The full spatial-height
domain and the mixed intersection are equivalent for this forced
history.

Pressure completes the momentum rows without changing the prescribed
force. Write \(Q_n=\partial_t^np\), and use the decaying pressure
normalization
\begin{equation}
 \begin{aligned}
 \phi_n&=-(-\Delta)^{-1}\nabla\cdot F_n,\\
 Q_n&=R_iR_j\sum_{a=0}^n\binom na
                     V_{a,i}V_{n-a,j}+\phi_n,\\
 \nabla Q_n&=(I-\Pp)(F_n-B_n),
 \end{aligned}
 \label{r:eq:forced-domain-pressure}
\end{equation}
where \(R_i\) are the Riesz transforms and repeated spatial indices
are summed. The force potential is finite in every \(H^M\).
At low frequencies, \(\widehat F_n\) is uniformly bounded and its
inverse-derivative multiplier is controlled by
\(\int_{|\xi|\le1}|\xi|^{-2}\dd\xi=4\pi\).
At high frequencies, the prescribed Sobolev norms suffice. The smooth
time dependence of the source gives
\(\phi_n\in C([0,T];H^M)\) for every fixed \(n,M\).

The Riesz transforms are bounded on \(H^M\), and
\(H^{M+1}\cdot H^{M+1}\subset H^M\) in three dimensions.
For every integer \(M\ge0\), these product bounds and
Cauchy--Schwarz in time give
\begin{equation}
 \begin{aligned}
 &\norm{Q_n}_{L^1H^M}
       +\norm{\nabla Q_n}_{L^1H^{M-1}}\\
 &\quad\le c_M\sum_{a=0}^n\binom na
       \norm{V_a}_{L^2H^{M+1}}\norm{V_{n-a}}_{L^2H^{M+1}}\\
 &\qquad\quad+\norm{\phi_n}_{L^1H^M}
       +\norm{\nabla\phi_n}_{L^1H^{M-1}}<\infty.
 \end{aligned}
 \label{r:eq:forced-domain-pressure-integrals}
\end{equation}
Here all time norms are over \((0,T)\). In particular, \(M=0\)
uses \(H^{-1}\) for the gradient norm.
The velocity part of each finite rectangle satisfies
\begin{equation}
 R_{N,M}
 \le T\sum_{n=0}^N
 \left(\mathsf U_{n,M+1}^2+
       \mathsf U_{n+1,\max\{M-1,0\}}^2\right)<\infty.
 \label{r:eq:forced-domain-rectangles}
\end{equation}
Each prescribed \(F_n\) belongs to \(L^2(0,T;H^{M-1})\).
Together with \(Q_0,\ldots,Q_N\) in the spaces just proved and
\(F_0,\ldots,F_N\), this retains the full forced rectangle
\((V_0,\ldots,V_{N+1};Q_0,\ldots,Q_N;F_0,\ldots,F_N)\).
Every retained row obeys
\[
 V_{n+1}+B_n+\nabla Q_n-\nu\Delta V_n=F_n,
 \qquad \nabla\cdot V_n=0.
\]
Restriction to a smaller rectangle preserves all derivatives, pressure
terms, and force coordinates because these entries come from the
given history.

The next temporal row also controls approach to the endpoint.
For \(s,t<T\), integration of \(\partial_tV_n=V_{n+1}\) gives
\begin{equation}
 \norm{V_n(t)-V_n(s)}_{H^m}
 \le |t-s|\,\mathsf U_{n+1,m}.
 \label{r:eq:forced-domain-endpoint}
\end{equation}
Each row has an \(H^m\) limit at \(T\). If a trace is taken first in
a higher Sobolev space, its restriction is this limit by uniqueness
of the \(H^m\) limit. The traces assemble into one \(H^\infty\) jet,
with \(u_*=V_0(T)\in H^\infty_\sigma\).
Product continuity and \eqref{r:eq:forced-domain-pressure} give the
pressure traces and preserve the complete terminal equation:
\[
 V_{n+1}(T)+B_n(T)+\nabla Q_n(T)-\nu\Delta V_n(T)=F_n(T).
\]
Choosing \(m>k+3/2\) in \eqref{r:eq:forced-domain-endpoint} gives uniform
convergence of every spatial derivative of order at most \(k\).
The finite products in the pressure formula give the corresponding
limits for pressure and its gradient.

If the force is prescribed smoothly beyond \(T\), local forced
existence restarts the equation from \(u_*\).
The outgoing initial jet is fixed by
\eqref{r:eq:forced-domain-row-recurrence} and
\eqref{r:eq:forced-domain-pressure}, so it agrees with the incoming
jet. Local uniqueness joins these histories into a smooth solution
through \(T\).


\subsection{Material derivatives of incompressibility}

Material acceleration is generally not divergence free, even when
the velocity is. Write \(U_{ij}=\partial_j u_i\).
Differentiating \(\nabla\cdot u=0\) along a trajectory gives
\[
 0=D_t(\nabla\cdot u)
   =\nabla\cdot(D_tu)-\operatorname{tr}(U^2).
\]
The pressure equation consequently agrees with
\[
 \nabla\cdot(D_tu)=\operatorname{tr}(U^2)
 =-\Delta p+\nabla\cdot f.
\]
Every ordinary mixed derivative of \(\nabla\cdot u\) vanishes. Every
material derivative must retain the transport commutator. With
\(A_k=D_t^ku\) and \(g_k=\nabla\cdot A_k\), this reads
\[
 g_0=0,\qquad
 g_{k+1}=D_tg_k+\partial_i u_j\,\partial_j(A_k)_i.
\]
Thus the acceleration, jerk, and higher parcel responses belong to one
constraint hierarchy. Appendix~\ref{r:app:endpoint} derives its first
three rows and the complete conversion to mixed Eulerian derivatives.

For the deformation gradient \(F=\partial_\ell X\), the volume ratio
\(j=\det F\) obeys \(j_t=0\) and
\[
 j_{tt}=\bigl[\nabla\cdot(D_tu)-\operatorname{tr}(U^2)\bigr](X,t)j=0.
\]
The term \(\operatorname{tr}(U^2)\) accounts for the changing
deformation of the parcel. Volume preservation is consistent with a
nonzero divergence of acceleration. Quantitative bounds on the
acceleration require Sobolev control of the velocity, as in the
following corollary.

\begin{corollary}[Momentum and acceleration bounds from \(R_{1,2}\)]
For a locally smooth solution of \eqref{r:eq:ns} with the stated force,
the single condition \(R_{1,2}<\infty\) suffices for uniform bounds on
all four undifferentiated momentum terms. Set
\[
 R:=R_{1,2},\qquad A:=\sqrt{(1+T^{-1})R},
 \qquad F:=L_{0,2},\qquad
 G:=\sup_t\norm{\Pp f(t)}_{H^2}.
\]
Then
\begin{equation}
 \begin{aligned}
 \norm{u_t}_{L^\infty_{t,x}}&\le c_\infty A,\\
 \norm{(u\cdot\nabla)u}_{L^\infty_{t,x}}&\le A^2/(8\pi),\\
 \norm{\nabla p}_{L^\infty_{t,x}}&\le A^2/(2\pi)+c_\infty F,\\
 \norm{\nu\Delta u}_{L^\infty_{t,x}}&\le A^2/(2\pi)+c_\infty(A+G).
 \end{aligned}
 \label{r:eq:one-rectangle}
\end{equation}
The total material acceleration obeys
\begin{equation}
 \norm{D_tu}_{L^\infty_{t,x}}\le c_\infty A+\frac{A^2}{8\pi}.
 \label{r:eq:material-one-rectangle}
\end{equation}
Along each material trajectory,
\[
 |u(X(b,t),t)|
 \le |u_0(b)|+t\left(c_\infty A+\frac{A^2}{8\pi}\right).
\]
Also \(R_{0,1}\le R\), so \eqref{r:eq:critical-production} bounds the
integrated critical production by
\(\sqrt{D+R}\,R/(\sqrt2\pi)\), with \(D=\norm{u_0}_{H^1}^2\).
\end{corollary}

\begin{proof}
The rectangle contains \(u,u_t\in L^2(0,T;H^3)\). The time trace in
\(H^3\) gives
\[
 \sup_t\norm{u(t)}_{H^3}^2
 \le T^{-1}\norm u_{L^2H^3}^2
 +2\norm u_{L^2H^3}\norm{u_t}_{L^2H^3}\le A^2.
\]
The adjacent pair \(u_t\in L^2H^3\), \(u_{tt}\in L^2H^1\)
gives \(\sup_t\norm{u_t}_{H^2}\le A\) by \eqref{r:eq:trace}.
For \(B=(u\cdot\nabla)u\), pointwise embedding gives
\(\norm B_\infty\le c_\infty^2 A^2\).
The weighted Fourier convolution used in \eqref{r:eq:product}, applied
directly to this contraction, also gives
\[
 \norm B_{H^2}
 \le4c_\infty\norm u_{H^2}\norm{\nabla u}_{H^2}
 \le4c_\infty A^2.
\]
Euclidean and Frobenius norms bound the contraction with constant one.
Finally retain both projected identities,
\[
 \nabla p=(I-\Pp)(f-B),\qquad
 \nu\Delta u=u_t+\Pp B-\Pp f.
\]
Their \(H^2\) bounds and \eqref{r:eq:embedding} prove
\eqref{r:eq:one-rectangle}. Adding \(u_t+B\) proves
\eqref{r:eq:material-one-rectangle}. The bounded \(H^3\) velocity supplies
bounded speed and spatial gradient, hence trajectories throughout the
interval. Integration of their acceleration proves the last bound.
The critical estimate follows by restriction to \(R_{0,1}\).
\end{proof}

The finite condition \(R_{1,2}<\infty\) controls both the Eulerian
momentum terms and the acceleration along trajectories, without a
smallness restriction. Higher material derivatives are polynomials
in finitely many mixed Eulerian derivatives. Appendix~\ref{r:app:endpoint}
proves the two-way conversion between these derivative families,
gives an explicit rectangle for each requested order, and derives
their endpoint traces and convergence rates. General production,
pressure, and acceleration estimates are developed in
Appendix~\ref{r:sec:recovered-rectangle}.

\section{Exact cancellation in the prescribed exterior}
\label{r:sec:exact-cancellation}

The proposed construction combines a concentrating core, an exact
swirling exterior, and spatial cutoffs. We begin with the exterior,
where all correction fields vanish and the velocity and pressure have
explicit formulas \cite[p.~116, Step~4; p.~138, Lemma~A.6]{openai-paper}.
This region already gives positive contributions to the Sobolev norms
used above. After calculating its momentum balance and integer-order
growth, we obtain fractional lower bounds that remain valid for any
interior completion agreeing with this exterior.

Work at viscosity one. Write
\(\tau=1-t\), \(A=\tfrac12+h\), \(D=\tfrac12-h\), with
\(0<h<1/100\), and \(b=2^A\kappa>0\).
The amplitude \(\kappa\) is the constant denoted by \(c_\infty\) in
that construction, unrelated to our embedding constant. Define
\begin{equation}
 H(Z)=\frac1{\Gamma(1+h)}\int_0^\infty
 e^{-v}v^h(1+Zv)^{-h}\dd v,
 \qquad K(r,t)=b r^{-2A}H(4\tau/r^2).
 \label{r:eq:heat-integral}
\end{equation}
On its exterior the velocity and pressure are
\begin{equation}
 u=K(r,t)e_\theta,\qquad
 p=-\int_r^\infty K(\rho,t)^2\frac{\dd\rho}{\rho}.
 \label{r:eq:uncut-exterior}
\end{equation}
These fields are smooth for \(r>0,t<1\), and the pressure integral
converges since its integrand decays as \(\rho^{-3-4h}\).

\begin{proposition}[Exact exterior cancellation and failure of a finite rectangle]
\label{r:prop:exterior-cancellation}
The exterior \eqref{r:eq:uncut-exterior} has zero momentum residual.
On the shrinking annuli specified below, all four individual momentum
terms grow without bound. Any whole-space velocity agreeing with this
exterior on those annuli satisfies
\[
 R_{0,1}=\int_0^1
 \bigl(\norm u_{H^2}^2+\norm{u_t}_2^2\bigr)\dd t=\infty.
\]
The localized field prescribed in \cite{openai-paper} agrees with the
exterior on these annuli.
Its material acceleration also fails to belong to \(L^2_{t,x}\), although
the exterior trajectories preserve volume and satisfy every differentiated
material incompressibility identity.
\end{proposition}

\begin{proof}
Differentiation under the integral, justified by an integrable exponential
majorant, gives
\[
 H'(Z)=-\frac{h}{\Gamma(1+h)}\int_0^\infty
 e^{-v}v^{h+1}(1+Zv)^{-h-1}\dd v<0.
\]
Thus \(H(0)=1\), \(H'(0)=-h(1+h)\), and \(H\) is positive.
Integration by parts in \(v\), with vanishing boundary terms, gives
\[
 Z^2H''+[1+2(1+h)Z]H'+h(1+h)H=0.
\]
Substitution into the radial derivatives of \(K\) proves
\(K_t=(\partial_{rr}+r^{-1}\partial_r-r^{-2})K\).
The angular unit vector contributes the last term. Computing transport
and the pressure derivative now gives the four signed momentum terms:
\begin{equation}
 \begin{aligned}
 u_t&=K_t e_\theta,&
 (u\cdot\nabla)u&=-K^2r^{-1}e_r,\\
 \nabla p&=K^2r^{-1}e_r,&
 -\Delta u&=-K_t e_\theta.
 \end{aligned}
 \label{r:eq:exact-pairs}
\end{equation}
Their sum is zero on this open region, as is every fixed mixed derivative
of their sum.

We next locate an open region where the exact localized field has these
formulas. Its coordinates satisfy
\(\tau=q(1-\eta^2)\), \(z=q^D\eta\), and \(X=r^2/(2q)\)
\cite[(3.2)]{openai-paper}.
Fix \(0<d<1\) and choose \(c>0\) with
\(c^2(1-d^2)/2>X_{\rm ext}\). Set
\begin{equation}
 \mathcal E_\tau=
 \{c\sqrt\tau<r<2c\sqrt\tau,\quad |z|<d\tau^D\}.
 \label{r:eq:shrinking-annulus}
\end{equation}
Since \(q\ge\tau\), throughout this region
\[
 |\eta|\le d,\qquad
 \tau\le q\le\tau/(1-d^2),\qquad
 X\ge c^2(1-d^2)/2>X_{\rm ext}.
\]
All corrections and the vector potential vanish there
\cite[p.~116, Step~4]{openai-paper}. For sufficiently small \(\tau\),
\(q<q_*\), \(2c\sqrt\tau<r_0/2\),
\(d\tau^D<z_0/2\), and \(\tau<\tau_0/2\).
Both cutoffs in \cite[(10.3)--(10.5)]{openai-paper} equal one.
This proves agreement with \eqref{r:eq:uncut-exterior} throughout
\(\mathcal E_\tau\).

Take the Eulerian derivatives at fixed spatial position before evaluating
at \(r=y\sqrt\tau\). The positive functions
\[
 k(y)=b y^{-2A}H(4/y^2),\qquad
 j(y)=-4b y^{-2A-2}H'(4/y^2)
\]
satisfy
\begin{equation}
 \begin{aligned}
 |u|&=k(y)\tau^{-1/2-h},\\
 |u_t|=|\Delta u|&=j(y)\tau^{-3/2-h},\\
 |(u\cdot\nabla)u|=|\nabla p|
 &=k(y)^2y^{-1}\tau^{-3/2-2h}.
 \end{aligned}
 \label{r:eq:exterior-powers}
\end{equation}
Every coefficient is strictly positive. The individual terms diverge
while cancelling exactly as in \eqref{r:eq:exact-pairs}.

Finally integrate over the annulus, rather than infer a Sobolev norm from
one point. Cylindrical integration and \(r=y\sqrt\tau\) give
\begin{equation}
 \begin{aligned}
 \int_{\mathcal E_\tau}|u_t|^2\dd x
 &=\int_{\mathcal E_\tau}|\Delta u|^2\dd x\\
 &=C_{c,d}\tau^{-3/2-3h},\qquad
 C_{c,d}=4\pi d\int_c^{2c}j(y)^2y\dd y>0.
 \end{aligned}
 \label{r:eq:exterior-integral}
\end{equation}
The coefficient is finite and positive by continuity and strict positivity
on the fixed compact interval. Since
\(\norm u_{H^2}^2\ge\norm{\Delta u}_2^2\), for any sufficiently small
\(\varepsilon>0\) and \(0<a<\varepsilon\),
\begin{equation}
 \int_{1-\varepsilon}^{1-a}
 (\norm u_{H^2}^2+\norm{u_t}_2^2)\dd t
 \ge\frac{2C_{c,d}}{1/2+3h}
 \left(a^{-1/2-3h}-\varepsilon^{-1/2-3h}\right).
 \label{r:eq:rectangle-divergence}
\end{equation}
Letting \(a\downarrow0\) proves the claimed rectangle divergence.

Material differentiation retains the rotation of the cylindrical basis.
Writing \(a_{\rm mat}=D_tu\), we obtain
\begin{equation}
 \begin{aligned}
 a_{\rm mat}&=-K^2r^{-1}e_r+K_t e_\theta,\\
 D_ta_{\rm mat}
 &=-3KK_t r^{-1}e_r+(K_{tt}-K^3r^{-2})e_\theta .
 \end{aligned}
 \label{r:eq:exterior-material}
\end{equation}
The radial and angular acceleration components are orthogonal. Hence
\begin{equation}
 \begin{aligned}
 \int_{\mathcal E_\tau}|a_{\rm mat}|^2\dd x
 &=C_{c,d}\tau^{-3/2-3h}
   +B_{c,d}\tau^{-3/2-5h},\\
 B_{c,d}&=4\pi d\int_c^{2c}\frac{k(y)^4}{y}\dd y>0 .
 \end{aligned}
 \label{r:eq:exterior-material-integral}
\end{equation}
Both powers are nonintegrable at zero. This proves the failure of the
material \(L^2_{t,x}\) coordinate without identifying it with \(u_t\).

To verify volume preservation, follow a trajectory segment that stays
inside this open exterior. With initial coordinates
\((r_*,\theta_*,z_*)\), its exact motion is
\[
 r(t)=r_*,\qquad z(t)=z_*,\qquad
 \theta(t)=\theta_*+\Phi(r_*,t),\qquad
 \Phi(r,t)=\int_{t_0}^t\frac{K(r,s)}r\dd s .
\]
Its cylindrical coordinate Jacobian is one; the unchanged radius makes
its Cartesian volume Jacobian one as well. Every derivative of the volume
of a transported parcel vanishes. In particular, with
\(U_{ij}=\partial_j u_i\), direct differentiation gives
\[
 \nabla\cdot a_{\rm mat}
 =-\frac{2KK_r}{r}
 =\operatorname{tr}(U^2).
\]
This is precisely the acceleration constraint derived in
Appendix~\ref{r:app:endpoint}. Higher orders follow by differentiating the
identically constant volume Jacobian, or by its material divergence
recurrence. The growing acceleration norm violates none of these
identities.
\end{proof}

The region's kinetic energy instead equals
\[
 \int_{\mathcal E_\tau}|u|^2\dd x
 =4\pi d\left(\int_c^{2c}k(y)^2y\dd y\right)
 \tau^{1/2-3h}\longrightarrow0.
\]
The shrinking volume permits large speeds and derivatives without
requiring large kinetic energy in that region.

The cancellation has a direct mechanical meaning:
\(D_tu=\Delta u-\nabla p\) in this exterior.
The inward acceleration is supplied by pressure, and the change in swirl
speed is supplied by viscous diffusion. A zero external force does not
make either contribution zero. Moving them to one side of the momentum
law produces the two cancelling pairs in \eqref{r:eq:exact-pairs}.

The annulus \(\mathcal E_\tau\) is not a transported parcel: its radii and
height shrink while exterior parcels retain their own radii and heights.
Its decreasing volume and growing norms sample different particle labels.
Diffusion can increase
velocity at a point by transporting momentum from surrounding fluid;
its dissipation of total energy is an integral statement.

This exterior occupies a region with a moving inner boundary. It is not
a global unforced history on \(\R^3\), so applying the unforced
theorem to this restriction would change that theorem's hypotheses.
The whole candidate includes an interior and cutoff forcing. The
exterior calculation depends only on the formulas in the displayed
region; the interior residual estimate enters the source calculation
below.

\subsection{Kinetic energy, strain, and critical concentration}

The exterior's acceleration cost can be compared with its strain and
kinetic costs on the explicit annulus \eqref{r:eq:shrinking-annulus}.
For the azimuthal velocity \(K(r,t)e_\theta\), the symmetric velocity
gradient \(S=(\nabla u+\nabla u^{\mathsf T})/2\) has squared norm
\[
 |S|^2=\tfrac12(K_r-K/r)^2.
\]
Use \(K=\tau^{-A}k(r/\sqrt\tau)\) and cylindrical volume measure
\(r\,dr\,d\theta\,dz\). Direct substitution gives
\begin{equation}
\begin{aligned}
 \int_{\mathcal E_\tau}|S|^2\,dx
 &=C_S\tau^{-1/2-3h},\\
 C_S&=2\pi d\int_c^{2c}(k'(y)-k(y)/y)^2y\,dy>0,\\
 \int_{\mathcal E_\tau}|u|^2\,dx
 &=4\pi d\int_c^{2c}|k(y)|^2y\,dy\;\tau^{1/2-3h},\\
 \int_{\mathcal E_\tau}|u|^3\,dx
 &=4\pi d\int_c^{2c}|k(y)|^3y\,dy\;\tau^{-4h}.
\end{aligned}
\label{r:eq:exterior-strain-critical-costs}
\end{equation}
Here \(C_S>0\) because the analytic, decaying heat profile cannot obey
\(k'=k/y\) throughout an open annulus. The first time integral is finite
since \(h<1/6\), and the kinetic contribution tends to zero. The
\(L^3\) contribution grows without bound. This explains how finite
kinetic energy and finite viscous action coexist with critical
concentration in the displayed region. Their bounds do not substitute
for the critical coordinates controlled by the full intersection.

\subsection{Positive lower bounds at fractional Sobolev orders}

The exterior also determines lower bounds at the precise fractional heights used by the intersection argument. These bounds use pairs of points lying entirely in that region. No choice of the unknown interior can cancel their contribution.

\begin{theorem}[Critical mixed-row lower bounds]\label{r:joined-thm:criticalexterior}
Suppose a smooth whole-space velocity agrees with \eqref{r:eq:uncut-exterior} on the prescribed exterior for all sufficiently late times. For every fixed pair of integers $n,m\ge0$, there are $c_{n,m}>0$ and $\delta_{n,m}>0$ such that
\begin{equation}\label{r:joined-eq:criticalexterior}
\norm{\partial_t^n u(1-\tau)}_{\dot H^{m+1/2}}^2
\ge c_{n,m}\tau^{-2n-m-3h}
\qquad(0<\tau<\delta_{n,m}).
\end{equation}
The constants are allowed to depend on the two derivative orders. The bound also holds, with changed constants, after subtracting any history smooth through time one whose finite mixed Sobolev norms are bounded. In particular,
\begin{equation}\label{r:joined-eq:criticalexteriorbase}
\begin{aligned}
\norm{u(1-\tau)}_{\dot H^{1/2}}^2&\ge c_0\tau^{-3h},\\
\norm{u(1-\tau)}_{\dot H^{3/2}}^2&\ge c_1\tau^{-1-3h}.
\end{aligned}
\end{equation}
Thus the critical dissipation integral diverges at the proposed singular time.
\end{theorem}

\begin{proof}
For a scalar function $F$ on $\R^3$, the positive fractional seminorm formula is
\begin{equation}\label{r:joined-eq:gagliardopairs}
\norm F_{\dot H^{1/2}}^2
=\frac1{2\pi^2}
\iint_{\R^3\times\R^3}
\frac{|F(x)-F(y)|^2}{|x-y|^4}\,dx\,dy .
\end{equation}
It follows by Plancherel and
$\int_{\R^3}(1-\cos(\xi\cdot z))|z|^{-4}dz=\pi^2|\xi|$.
In particular any selected set of pairs supplies a lower bound.

Differentiate the heat profile in \eqref{r:eq:uncut-exterior}. On the exterior,
\[
\partial_t^n u
=\tau^{-A-n}k_n(r/\sqrt\tau)e_\theta,
\]
where
\begin{equation*}
\begin{aligned}
k_n(\rho)
={}&\frac{b4^n(h)_n}{\Gamma(1+h)}
\rho^{-2A-2n}\\
&\times\int_0^\infty e^{-v}v^{h+n}
(1+4v/\rho^2)^{-h-n}\,dv .
\end{aligned}
\end{equation*}
Here $(a)_n$ denotes the rising factorial, with $(a)_0=1$. For each fixed $m$,
\[
\begin{aligned}
k_n^{(m)}(\rho)
&=(-1)^m d_{n,m}\rho^{-2A-2n-m}
\bigl(1+O_{n,m}(\rho^{-2})\bigr),\\
d_{n,m}&=b4^n(h)_n(1+h)_n(2A+2n)_m>0.
\end{aligned}
\]
The expansion, including its derivatives, follows from the exponentially integrable profile formula. Choose a fixed sufficiently large $c$. At transverse points $(c,0)$ and $(2c,0)$, the second component of
$\partial_1^m[k_n(r)e_\theta]$ equals $k_n^{(m)}(c)$ and $k_n^{(m)}(2c)$. Their difference is nonzero. Two small fixed transverse patches retain a strictly positive gap.

Scale both patches by $\sqrt\tau$. Choose the first axial coordinate in an interval of length proportional to $\tau^D$, and the second within a distance proportional to $\sqrt\tau$ of it. Since $D<1/2$, both axial coordinates remain in the exterior slab for small $\tau$. The transverse patches can be placed in the exact exterior by taking $c$ large, and the physical localization cutoffs equal one there for sufficiently small $\tau$. The selected pair set has measure bounded below by a positive constant times
$\tau^{5/2+D}$. Its distances are at most a fixed multiple of $\sqrt\tau$. For $F=\partial_1^m\partial_t^n u_2$, the squared difference on those pairs is at least
$c'\tau^{-2A-2n-m}$.
Formula \eqref{r:joined-eq:gagliardopairs} now gives
\[
\norm F_{\dot H^{1/2}}^2
\ge c''\tau^{5/2+D-2-2A-2n-m}
=c''\tau^{-2n-m-3h}.
\]
The Fourier multiplier inequality
$|\xi_1|^{2m}|\xi|\le|\xi|^{2m+1}$
proves \eqref{r:joined-eq:criticalexterior}. Finally, the triangle inequality shows that subtracting a history with bounded mixed Sobolev norms preserves each divergent lower bound for small enough $\tau$.
\end{proof}

The first pressure-complete nonlinear prefix must also grow. Define
\[
P_{1/2}(t)
=-\ip{\Lambda^{1/2}u}
{\Lambda^{1/2}\PP((u\cdot\nabla)u)}.
\]
For $t_0<1$ fixed sufficiently late, the exact critical balance gives
\begin{equation*}
\begin{aligned}
\int_{t_0}^{t}P_{1/2}(s)\,ds
={}&E_{1/2}(t)-E_{1/2}(t_0)\\
&+\nu\int_{t_0}^{t}\norm u_{\dot H^{3/2}}^2\,ds
-\int_{t_0}^{t}\ip{u}{\Lambda\PP f}\,ds .
\end{aligned}
\end{equation*}
The last integral is bounded by the kinetic budget and the smooth source. Inserting \eqref{r:joined-eq:criticalexteriorbase} yields constants $c_*>0$ and $C_*<\infty$ such that
\begin{equation}\label{r:joined-eq:exteriorcriticalprefix}
\int_{t_0}^{1-\tau}P_{1/2}(s)\,ds
\ge c_*\tau^{-3h}-C_* .
\end{equation}
Every completion satisfying the stated forced equation must obey this
lower bound. Its signed critical production consequently has an
unbounded time integral, contradicting the finite production bound
for a solution in the full intersection. The exterior calculation
establishes the lower bound independently of that membership premise.


\subsection{Divergence of the positive time convolutions}

The mixed-row lower bound also supplies the integrated divergence
required in Section~\ref{r:sec:all-row-primitives}. Write
\(H_r=E_{r,1/2}=\tfrac12\|\partial_t^r u\|_{\dot H^{1/2}}^2\).
For a constant \(d_r>0\), Theorem~\ref{r:joined-thm:criticalexterior} gives
\[
 H_r(1-\tau)\ge d_r\tau^{-2r-3h}
\]
when \(\tau\) is sufficiently small. Let \(I^r\) denote the
\(r\)-fold integral from time zero, defined in
\eqref{r:eq:source-primitive-def}. For \(r\ge1\), its convolution
kernel is positive. Restrict its lag
\(\ell=(1-\tau)-a\) to \([\tau,2\tau]\). On that interval,
\(1-a=\tau+\ell\le3\tau\), the kernel is at least
\(\tau^{r-1}/(r-1)!\), and the interval has length \(\tau\). Thus
\begin{equation}
 I^rH_r(1-\tau)
 \ge\frac{d_r3^{-2r-3h}}{(r-1)!}\tau^{-r-3h}
 \longrightarrow\infty.
\label{r:eq:positive-row-convolution}
\end{equation}
The case \(r=0\) uses the original pointwise lower bound. The source
primitive and initial polynomial in the odd-row identity stay bounded.
Consequently every odd completed height requires unbounded upper
nonlinear production for a whole-space solution with this exterior.
The conclusion uses the positive convolution explicitly; it does not
infer integrated divergence from pointwise growth alone.


\subsection{A material trajectory in the vortex core}
\label{r:sec:axis-trajectory}

The exterior annuli sample different particles as they shrink. Inside the
core, the stated formulas permit a stronger calculation: one particle
approaches the origin while its velocity and every material derivative
increase without bound. We can recover that trajectory directly from
the velocity on the axis.

The leading construction fixes \(U(0,\eta)=4\eta+j_0\), with
\(0<j_0\le0.05\)
\cite[p.~37; (B.1), (B.12)]{openai-paper}.
Every positive-order axial profile has zero axis value
\cite[p.~47, Lemma~5.1]{openai-paper}.
The cutoffs depend on \(q(z,t)\); their additional curl terms are radial
and leave this axial value unchanged. The wave and mean corrections
vanish near the axis at each \(t<1\)
\cite[(5.45); pp.~114--115]{openai-paper}.
Thus the candidate's exact axis trace, where the final cutoffs equal one,
is
\begin{equation}
 u(0,0,z,t)=q^{-A}(4\eta+j_0)e_z,\qquad
 z=q^D\eta,\quad 1-t=q(1-\eta^2).
 \label{r:eq:axis-trace}
\end{equation}

Define \(H_c(\eta)=D\eta+(1-\eta^2)(4\eta+j_0)\).
Its values at \(-j_0/4\) and zero have opposite signs, and
\(H_c'>0\) between them. Let \(\eta_c\) be its unique zero in that
interval, and put
\[
 d_c=1-\eta_c^2,\qquad U_c=4\eta_c+j_0
       =-\frac{D\eta_c}{d_c}>0.
\]
For sufficiently late \(t_0<1\), the curve
\begin{equation}
 x_c(t)=\bigl(0,0,\eta_c[(1-t)/d_c]^D\bigr),
 \qquad t_0\le t<1,
 \label{r:eq:axis-path}
\end{equation}
lies inside the cutoff plateaux and local domain. Differentiating its
axial component gives
\[
 \dot z_c=-D\eta_c d_c^{-D}(1-t)^{-A}
         =d_c^A U_c(1-t)^{-A}
         =u_z(x_c(t),t).
\]
The curve is a material trajectory, rather than a choice of spatial
sampling points. Along it, with \(C_c=d_c^AU_c>0\),
\begin{equation}
 (D_t^n u)(x_c(t),t)
 =C_c(A)_n(1-t)^{-A-n}e_z,\qquad n\ge0,
 \label{r:eq:axis-material-tower}
\end{equation}
where \((A)_0=1\) and
\((A)_n=A(A+1)\cdots(A+n-1)\). Acceleration, jerk, and every
fixed higher material derivative diverge along this particle. Yet its
position has a finite limit, since \(A<1\).

This growth remains compatible with local volume preservation. Ordinary
axial differentiation of \eqref{r:eq:axis-trace} yields
\[
 \partial_z u_z(x_c(t),t)=\frac{\beta_c}{1-t},\qquad
 \beta_c=\frac{d_c[-2A\eta_c U_c+4d_c]}
                   {1-2h\eta_c^2}>0.
\]
Axis symmetry and incompressibility give transverse contraction rates
\(-\beta_c/[2(1-t)]\), with an additional transverse rotation.
The linearized flow from \(t_0\) has axial stretch
\([(1-t_0)/(1-t)]^{\beta_c}\) and transverse length factors
\([(1-t)/(1-t_0)]^{\beta_c/2}\). Rotation changes neither factor's
determinant. Their volume product is exactly one at every \(t<1\).
Unbounded stretch and vanishing transverse widths coexist with constant
material volume; differentiating that constant produces zero at every
order.

For the source along this trajectory, assume the core residual estimate
\cite[(9.20)]{openai-paper}: every finite Eulerian derivative of \(f\)
is \(O(q^N)\) for every \(N\). Along \eqref{r:eq:axis-path}, each fixed
derivative of \(x_c\) has only a finite power singularity. The chain
rule gives
\[
 \frac{\dd^k}{\dd t^k}f(x_c(t),t)=O((1-t)^N)
 \quad\text{for every fixed }k,N\ge0.
\]
Under this residual estimate, every fixed material derivative of the
force approaches zero along the particle, while
\eqref{r:eq:axis-material-tower} diverges. The momentum law permits this
distinction because pressure and viscous stresses also contribute to
the acceleration.

\section{Admission of the candidate force to the forced theorem}
\label{r:sec:force-extension}

Part~II applies after the momentum residual of the proposed field has been
placed in its prescribed force class. The cutoff calculation below keeps the
full localization error, proves uniform bounds for every fixed mixed
derivative, and produces compatible compactly supported terminal jets.

\subsection{The localization residual}

Multiplication by a spatial cutoff changes the fluid's velocity gradients.
The resulting force must include those changes, including the correction
needed to preserve incompressibility. OpenAI's local representation is
\(v=\operatorname{curl}\mathcal A+B e_\theta\), with \(B\) independent
of the angle. Its localization is exactly
\begin{equation}
 u=cv+w,\qquad w=\nabla c\times\mathcal A,\qquad p=c\pi,
 \qquad c=\chi_x\chi_t.
 \label{r:eq:localization-fields}
\end{equation}
Here \(\pi\) is the local pressure, and \(c\) is axisymmetric.
This is \cite[(10.4)]{openai-paper}: the curl acts on \(c\mathcal A\),
while the swirl is multiplied directly by \(c\). In particular,
\(\nabla\cdot u=0\). Introducing a vector potential for the exterior
swirl and localizing that additional potential would change the field.

Write \(\mathcal R_\nu(v,\pi)=v_t+(v\cdot\nabla)v-\nu\Delta v+\nabla\pi\).
Expanding \eqref{r:eq:localization-fields} gives the entire localization error:
\begin{equation}
\begin{aligned}
 f-c\mathcal R_\nu(v,\pi)
 ={}&(c_t-\nu\Delta c)v-2\nu\sum_{i=1}^3(\partial_i c)\partial_i v
       +\pi\nabla c\\
 &+(c^2-c)(v\cdot\nabla)v+c(v\cdot\nabla c)v\\
 &+w_t-\nu\Delta w+c(v\cdot\nabla)w\\
 &+c(w\cdot\nabla)v+(w\cdot\nabla)w.
\end{aligned}
\label{r:eq:localization-residual}
\end{equation}
The apparent additional product \((w\cdot\nabla c)v\) vanishes because
\((\nabla c\times\mathcal A)\cdot\nabla c=0\).
Every remaining term belongs to the prescribed force. Where
both cutoffs are constant at one, \(w=0\) and the error vanishes.
Where they vanish on a neighborhood, the localized fields vanish too.
The extra force is confined to the transitions between these regions.

In the heat exterior, near time one, \(\mathcal A=0\) and
\(v=K e_\theta\). The formula becomes a concrete force field:
\begin{equation}
\begin{aligned}
 f={}&\left[\chi_x(1-\chi_x)\frac{K^2}{r}
                 +\pi\partial_r\chi_x\right]e_r
        +\pi\partial_z\chi_x e_z\\
 &-\nu\left[2(\partial_r\chi_x)K_r
       +K\left(\partial_{rr}\chi_x+r^{-1}\partial_r\chi_x
                         +\partial_{zz}\chi_x\right)\right]e_\theta.
\end{aligned}
\label{r:eq:exterior-cutoff-force-full}
\end{equation}
The radial force accounts for the changed centrifugal balance and the
pressure cutoff. The axial term comes from the pressure cutoff alone.
The angular force accounts for diffusion across the tapered swirl.
All three vanish where \(\chi_x\) is constant at one. The bounds below
control every finite derivative on their transition support.

For a space--time region \(\Sigma\), define the finite mixed-jet
seminorm
\[
 \mathcal J_m(F;\Sigma)
 =\sum_{|\gamma|\le m}\frac{\norm{\partial^\gamma F}_{L^\infty(\Sigma)}}{\gamma!},
 \qquad \gamma\in\mathbb N_0^4.
\]
The four entries of \(\gamma\) count three spatial derivatives and one
temporal derivative. The multinomial product rule gives
\(\mathcal J_m(FG)\le\mathcal J_m(F)\mathcal J_m(G)\), also for dot
and cross products. Each coordinate derivative obeys
\(\mathcal J_m(\partial_i F)\le(m+1)\mathcal J_{m+1}(F)\).

On a transition region put
\[
 \begin{aligned}
 C_m&=\mathcal J_m(c),& V_m&=\mathcal J_m(v),\\
 A_m&=\mathcal J_m(\mathcal A),& P_m&=\mathcal J_m(\pi).
 \end{aligned}
 \qquad
 U_m=C_mV_m+3(m+1)C_{m+1}A_m.
\]
Then \(\mathcal J_m(u)\le U_m\) and
\(\mathcal J_m(p)\le C_mP_m\). Differentiating the full momentum
residual gives the explicit force-jet bound
\begin{equation}
\begin{aligned}
 \mathcal J_m(f)\le{}&(m+1)U_{m+1}
       +3(m+1)U_mU_{m+1}\\
 &+3\nu(m+1)(m+2)U_{m+2}
       +3(m+1)C_{m+1}P_{m+1}.
\end{aligned}
\label{r:eq:localized-jet-bound}
\end{equation}
This bounds every fixed temporal and spatial source derivative once the
local field jets on the transition have been bounded. It requires no
constant shared by infinitely many orders.

\subsection{Uniform derivative bounds on the cutoff transitions}

The relevant bounds follow from the actual support geometry. The cutoff
in \cite[p.~118]{openai-paper} equals one for
\(r\le r_0/2\), \(|z|\le z_0/2\), and has compact support inside
\(r<r_0\), \(|z|<z_0\). Every spatial cutoff error in
\eqref{r:eq:localization-residual}, including \(c^2-c\) near time one,
is supported where
\[
 r\ge r_0/2\quad\hbox{or}\quad |z|\ge z_0/2.
\]
The construction's coordinates satisfy
\[
 \tau=q(1-\eta^2),\quad z=q^D\eta,\quad
 X=\frac{r^2}{2q},\quad |\eta|\le1,
 \qquad D=\frac12-h>0.
\]
On the axial transition, \(q\ge |z|^{1/D}\) yields
\[
 q\ge q_z:=(z_0/2)^{1/D}>0.
\]
This remains true as \(r\) tends to zero. On the radial transition,
any correction inside \(X\le X_{\rm ext}\) satisfies
\[
 q=\frac{r^2}{2X}\ge q_r:=\frac{r_0^2}{8X_{\rm ext}}>0.
\]
Thus every interior correction meeting a spatial transition has
\(q\ge\min(q_z,q_r)>0\). The cutoff choices also ensure
\(q<q_*/2\), by \cite[(10.3)]{openai-paper}. This is precisely the
closed scale range covered by Theorem~3.1(ii) of that paper.

The one-sided interior bound has a concrete finite-order meaning.
On \(q\ge\delta>0\), only finitely many summands in
\cite[(9.21)]{openai-paper} survive, since their cutoff scales tend to
infinity. The finite representatives, fixed phase labels, and closed
parameter bounds are those specified in Proposition~9.9, Step~3.
Coordinate conversion has denominator
\(1-2h\eta^2\ge1-2h>0\). At the axis, the potential is represented
in Cartesian coordinates as \(a(r^2,z,t)(-x_2,x_1,0)\), as in
\cite[(10.1)]{openai-paper}. These are the stated interior inputs to
the finite numbers \(V_m,A_m,P_m\) in \eqref{r:eq:localized-jet-bound}.

If \(q\) approaches zero on the remaining radial transition, the field
is instead in \(X\ge X_{\rm ext}\). There
\(\mathcal A=0\), \(v=K e_\theta\), and
\(\pi=-\int_r^\infty K^2\dd\rho/\rho\).
The heat-profile formula gives
\[
 |\partial_\tau^jK|\le C_jr^{-1-2h-2j},\qquad
 |\partial_\tau^j\pi|
 \le\frac{C_j}{2+4h+2j}r^{-2-4h-2j}.
\]
Their further spatial derivatives are bounded at \(r\ge r_0/2\).
Temporal cutoff derivatives occur only where
\(\tau_0/2\le\tau\le\tau_0\), which gives \(q\ge\tau_0/2\).
Every transition thus supplies finite jets for
\eqref{r:eq:localized-jet-bound}.

Near the proposed singular point both cutoffs equal one. The source is
then the local residual itself. Theorem~3.1(iii) supplies its all-orders
flatness on \(X\le X_{\rm ext}\); beyond that region the residual is
exactly zero. Combining this input with the transition bounds proves
\begin{equation}
 \sup_{x\in\R^3,\,0\le t<1}|\partial_x^\alpha\partial_t^j f(x,t)|<\infty
 \qquad\hbox{for every finite }\alpha,j.
 \label{r:eq:localized-source-jets}
\end{equation}
One additional time derivative in this bound makes each preceding jet
uniformly Cauchy as \(t\uparrow1\). Uniform spatial derivative limits
give compatible smooth terminal fields with fixed compact support.
Let \(K\subset\mathbb R^3\) be a fixed compact set containing these supports.
The source's finite Sobolev coordinates can now be bounded explicitly.
Put
\[
 M_{N,M}=\max_{0\le j\le N,\ |\alpha|\le M}
 \sup_{x,t<1}|\partial_x^\alpha\partial_t^j f(x,t)|.
\]
Expanding the integer Fourier weight \((1+|\xi|^2)^M\), whose
nonnegative multinomial coefficients sum to \(4^M\), gives
\begin{equation}
 \sum_{j=0}^N\norm{\partial_t^jf}_{L^2(0,1;H^M)}^2
 \le (N+1)4^M|K|M_{N,M}^2<\infty.
 \label{r:eq:localized-source-rectangle}
\end{equation}
Thus the actual localized source belongs to every finite mixed Sobolev
coordinate under the stated local inputs. The next construction extends
its terminal fields through time one.

Localization thus preserves source smoothness under the interior
hypotheses of \cite[Theorem~3.1]{openai-paper}. The local residual
estimate \cite[(3.4)]{openai-paper}, based on the finite-stage estimate
\cite[(9.18)]{openai-paper}, is an input to this conclusion; the
calculation above bounds the additional terms produced by the cutoffs.

\subsection{Continuing the terminal force jets}


Write \(F_j\in C_c^\infty(\R^3;\R^3)\), with
\(\operatorname{supp}F_j\subset K\), for the terminal field obtained from
\(\partial_t^j f(\cdot,t)\) as \(t\uparrow1\). The source vanishes near
\(t=0\), and all its left terminal jets have support in one fixed compact
set \(K\). These jets can be continued through time one without imposing a
bound common to all derivative orders.

Choose \(\chi_0\in C_c^\infty(\R)\), equal to one near zero and zero for
arguments at least one. For \(\sigma=t-1\ge0\), define
\begin{equation}
 \widetilde f_{\rm OA}(x,1+\sigma)=\sum_{j=0}^\infty
 \chi_0(b_j\sigma)\frac{\sigma^j}{j!}F_j(x).
 \label{r:eq:force-borel}
\end{equation}
For \(0\le t<1\), set \(\widetilde f_{\rm OA}(x,t)=f(x,t)\), the full
localized momentum residual. Formula \eqref{r:eq:force-borel} is the
extension used in \cite[(10.11)]{openai-paper}. Its full product rule is
\[
 \partial_\sigma^m
 \left(\chi_0(b_j\sigma)\frac{\sigma^j}{j!}\right)
 =\sum_{\ell=\max(0,m-j)}^m
 \binom m\ell b_j^\ell\chi_0^{(\ell)}(b_j\sigma)
 \frac{\sigma^{j-m+\ell}}{(j-m+\ell)!}.
\]
On a summand's support, \(\sigma\le b_j^{-1}\), so every product
\(b_j^\ell\sigma^{j-m+\ell}\) is at most \(b_j^{m-j}\). Writing
\[
 C_{j,m}=\sum_{\ell=\max(0,m-j)}^m
 \binom m\ell\frac{\norm{\chi_0^{(\ell)}}_\infty}{(j-m+\ell)!},
\]
we obtain
\begin{equation}
 \left\|\partial_x^\alpha\partial_\sigma^m
 \left[\chi_0(b_j\sigma)\frac{\sigma^j}{j!}F_j\right]\right\|_\infty
 \le C_{j,m}b_j^{m-j}\norm{F_j}_{C^{|\alpha|}}.
 \label{r:eq:force-borel-bound}
\end{equation}
For each \(j\ge1\), choose \(b_j>b_{j-1}\), with \(b_0=1\), so that
this bound is at most \(2^{-j}\) for all
\(|\alpha|+m\le\lfloor j/2\rfloor\). There are finitely many constraints at
that step, and \(m-j<0\) in each one. Such a choice exists for every finite
collection of norms of \(F_j\).

For fixed \(\alpha,m\), the differentiated tail beginning at
\(j\ge\max\{1,2(|\alpha|+m)\}\) is uniformly summable. At
\(\sigma=0\), the \(m\)-th temporal derivative of the series receives a
contribution only from \(j=m\), and that contribution is \(F_m\). All left
and right jets agree. Each term retains support in \(K\) and vanishes for
\(\sigma\ge1\). Uniform differentiated convergence gives smoothness at the
joining time and at the outer support boundary. Hence
\begin{equation}
 \widetilde f_{\rm OA}\in
 C_c^\infty(\R^3\times(0,\infty);\R^3)
 \subset C^\infty([0,\infty);\mathcal S(\R^3;\R^3)),
 \label{r:eq:oa-force-class}
\end{equation}
and its restriction to \(t<1\) is exactly the localized momentum residual.
For the quantitative integrals below, write
\begin{equation}
 f_{\rm OA}:=\widetilde f_{\rm OA}\big|_{[0,1)}=f.
 \label{r:eq:oa-force-restriction}
\end{equation}
Its support lies in \(K\times[0,2]\), and it retains the prescribed initial
interval of rest. A different smooth continuation beyond time one has the
same restriction on every \([0,S]\), \(S<1\), so strict-subinterval
uniqueness gives the same comparison below.


\section{Strict-subinterval uniqueness and refutation}
\label{r:sec:exclusion}

A regular solution constructed from the prescribed data constrains every
classical candidate for those data. This follows from uniqueness before
any terminal bound on the candidate is assumed.

\begin{theorem}[Exclusion of the proposed finite-time blowup]
\label{r:thm:openai-exclusion}
No classical pair with the zero datum, positive viscosity, prescribed force
\(\widetilde f_{\rm OA}\), and exterior asserted in
\cite[Theorem~1.1]{openai-paper} can solve that forced initial-value problem
on \([0,1)\).
\end{theorem}

\begin{proof}
Apply Corollary~\ref{r:thm:supplied-forced} to
\((0,\nu,\widetilde f_{\rm OA})\), and let \(U,P\) be its regular solution on
\([0,1]\). Let \(v,p\) be a proposed classical
solution on \([0,1)\), with identical initial datum, prescribed force,
viscosity, and spatial domain. For the compactly supported candidate in
OpenAI's Theorem~1.1, every spatial Sobolev norm is finite on each strict
interval \([0,S]\), \(S<1\). Put \(z=v-U\), \(\pi=p-P\).
Subtracting their momentum equations gives
\begin{equation}
 \begin{aligned}
 z_t+(v\cdot\nabla)z+(z\cdot\nabla)U+\nabla\pi
 &=\nu\Delta z,\\
 \nabla\cdot z=0,\qquad z(0)&=0.
 \end{aligned}
 \label{r:eq:candidate-identification}
\end{equation}
The prescribed force cancels because both histories solve that one forced
problem. Pairing with \(z\) eliminates the pressure and the divergence-free
transport term. With \(L(t)=\norm{\nabla U(t)}_\infty\),
\[
 \frac12\frac{\dd}{\dd t}\norm z_2^2+\nu\norm{\nabla z}_2^2
 =-\int z_i z_j\partial_jU_i\dd x
 \le L(t)\norm z_2^2.
\]
The intersection for \(U\) supplies \(L\in L^1(0,1)\). Gronwall's
inequality on \([0,S]\) yields
\[
 \norm{z(t)}_2^2
 \le\norm{z(0)}_2^2\exp\left(2\int_0^t L(s)\dd s\right)=0.
\]
Since \(S<1\) was arbitrary, \(v=U\) throughout \([0,1)\). Every estimate for \(U\), including its mixed and material
derivative estimates, consequently applies to \(v\). No terminal regularity of \(v\) was used to
identify it with \(U\).

For rest data, the finite rectangle \(R_{0,2}[U]\) and
\eqref{r:eq:velocity-bound} give
\[
 \begin{aligned}
 \sup_{0\le t<1}\norm{v(t)}_\infty^2
 &=\sup_{0\le t<1}\norm{U(t)}_\infty^2\\
 &\le\frac{R_{0,2}[U]}{8\pi}<\infty.
 \end{aligned}
\]
This bound excludes the velocity growth in \cite[Theorem~1.1]{openai-paper}.
Section~\ref{r:sec:orthogonal-acceleration} also bounds the acceleration,
transport term, pressure gradient, and viscous term individually.
Each bound follows from finite mixed norms of \(U\), so it applies
irrespective of cancellation among the momentum terms.

The prescribed exterior already lies outside a lower-order rectangle.
Equation~\eqref{r:eq:rectangle-divergence} gives a positive contribution
whose time integral is infinite, so \(R_{0,1}=\infty\).
Every velocity in \(\mathscr I_1\) instead has \(R_{0,1}<\infty\).
No velocity containing that exterior belongs to the full-interval
compatible intersection.

Material coordinates preserve the conclusion.
Equation~\eqref{r:eq:exterior-material-integral} shows a divergent
material-acceleration integral, while the exact axis trajectory in
Section~\ref{r:sec:axis-trajectory} has unbounded velocity and every fixed
higher material derivative. Incompressibility is retained through the
volume Jacobian and its differentiated identities. The resulting norm blowup
is consequently present in both the Eulerian and material descriptions of the
proposed flow.

Equation~\eqref{r:eq:oa-force-class} gives the compactly supported smooth
continuation of the candidate's residual. The initial datum is zero and the
viscosity is positive, so Corollary~\ref{r:thm:supplied-forced} gives the
regular solution \(U\) used above. In particular,
\[
 \sup_{t<1}\norm{v(t)}_\infty^2\le R_{0,2}[U]/(8\pi)<\infty
\]
whereas the candidate has unbounded velocity. Its divergent critical
coordinate in Theorem~\ref{r:joined-thm:criticalexterior} gives a second
contradiction. Both comparisons use uniqueness on strict subintervals;
the terminal bounds follow from the global forced theorem proved in Part~II.
\end{proof}

\subsection{Quantitative source criteria and conservative forcing}

The appendices provide two quantitative source criteria for
continuation. They also reduce conservative forcing to the unforced
problem proved in Part~I. Each result
provides a comparison solution under its own hypotheses.
For the rest datum and unit terminal time, the two quantitative
hypotheses are, respectively,
\begin{equation}
 \begin{aligned}
 G_{\rm OA}&:=\int_0^1\norm{\Pp f_{\rm OA}(t)}_{H^3}\dd t
             <\frac{\pi\sqrt\nu}{8},\\
 B_{\rm OA}&:=\int_0^1\norm{\Pp f_{\rm OA}(t)}_{\dot H^{-1/2}}^2\dd t
             <\frac{4\pi^2\nu^3}{3}.
 \end{aligned}
 \label{r:eq:candidate-admission}
\end{equation}
Either inequality would independently give the displayed finite-horizon
estimates.
The explicit source fails both tests. Appendix~\ref{r:app:critical-source}
checks its actual residual directly and proves
\begin{equation}
 \begin{aligned}
 G_{\rm OA}&\ge\frac{\sqrt{\pi\nu}}2
                >\frac{\pi\sqrt\nu}{8},\\
 B_{\rm OA}&>\frac{4\pi^2\nu^3}{3}.
 \end{aligned}
 \label{r:eq:necessary-admission-budgets}
\end{equation}
The first bound follows from the \(H^3\) energy balance and the
exterior's growing spatial norm. The second uses the critical energy
balance up to its first threshold crossing, followed by the nonzero
source in the fixed cutoff region. Neither calculation assumes terminal
regularity of the velocity or global smooth continuation of the source.
Thus the prescribed force lies outside both quantitative source
classes. Equation~\eqref{r:eq:actual-pressure-obstruction} also excludes
the conservative-force reduction. Failure of these sufficient conditions does
not affect Corollary~\ref{r:thm:supplied-forced}, which follows from the
unrestricted forced theorem of Part~II.
The first-rectangle criterion in Section~\ref{r:sec:first-rectangle}
also fails for this candidate: its conclusion would give the finite
rectangle contradicted by \eqref{r:eq:rectangle-divergence}.
In each case the unbounded candidate norms prevent application
of the corresponding sufficient condition.


Theorem~\ref{r:thm:openai-exclusion} completes the formal refutation.  The
appendices that follow retain the two proof-lane coordinate catalogues and the
auxiliary source, endpoint, material, and quantitative calculations used in
Part~III.
\appendix

\section*{Appendix map}
\addcontentsline{toc}{section}{Appendix map}
The first catalogue records the canonical coordinates of the unforced proof;
the second records the distinct canonical coordinates of the forced proof.
The remaining appendices belong to the refutation and supply its source-work,
terminal-domain, parcel, and quantitative estimates.  Their placement after
the refutation does not move any appendix calculation upstream of the theorem
that uses it.

\UnforcedQuantitativeCatalogue
\ForcedQuantitativeCatalogue

\section{Source work and its iterated time integrals}
\label{r:sec:source-primitives}

The source work can be controlled further before invoking any higher
velocity rectangle. For the base temporal row, write \(W_s=W_{0,s}\)
and put \(h_s=\Lambda^{2s}\Pp f\). Moving the spatial derivatives onto
the prescribed field gives \(W_s=\ip{u}{h_s}\). This pairing measures
how the fluid responds to a smooth spatial weight. Kinetic energy alone bounds this pairing; no higher spatial norm
of the velocity enters the estimate.

The energy balance implies
\[
 \norm{u(t)}_2\le K_T,
 \qquad K_T=\norm{u_0}_2+\int_0^T\norm{f(t)}_2\dd t.
\]
Let \(S_s=\operatorname{Sym}\nabla h_s\). Substitution of the
momentum equation into the derivative of \(\ip{u}{h_s}\) gives
\begin{equation}
 W_s'=\ip{u}{(\partial_t+\nu\Delta)h_s}
 +\int_{\R^3}u_i u_j(S_s)_{ij}\dd x+\ip{f}{h_s}.
 \label{r:eq:source-first}
\end{equation}
Incompressibility removes the pressure pairing; integration by parts
moves the transport and viscous derivatives onto \(h_s\).
Using the matrix operator norm for \(S_s\), define
\begin{equation}
 \begin{aligned}
 \mathsf B_s&:=K_T\norm{h_s}_{L^\infty_tL^2_x},\\
 \mathsf L_s&:=K_T\norm{(\partial_t+\nu\Delta)h_s}_{L^\infty_tL^2_x}
 +K_T^2\norm{S_s}_{L^\infty_{t,x}}\\
 &\quad+\norm{f}_{L^\infty_tL^2_x}\norm{h_s}_{L^\infty_tL^2_x}.
 \end{aligned}
 \label{r:eq:source-work-constants}
\end{equation}
Every fixed constant is finite under the source assumptions, and
\(\norm{W_s}_\infty\le\mathsf B_s\),
\(\norm{W_s'}_\infty\le\mathsf L_s\). The projection need not preserve
rapid spatial decrease; its boundedness on Sobolev spaces supplies the
norms used here.

These two bounds control an integrated form of the complete source
correction at every spatial height. Define
\begin{equation}
 \begin{aligned}
 \mathcal S_{s,k}&:=\sum_{j=0}^{k-1}(-2\nu)^j
 \partial_t^{k-1-j}W_{s+j},\qquad k\ge2,\\
 I^0&:=\mathrm{Id},\qquad
 (I^a g)(t):=\frac1{(a-1)!}\int_0^t(t-v)^{a-1}g(v)\dd v,
 \quad a\ge1.
 \end{aligned}
 \label{r:eq:source-primitive-def}
\end{equation}

\begin{proposition}[Integrated source correction]
\label{r:prop:source-primitive}
For each fixed \(s\ge0\) and integer \(k\ge2\),
\(I^{k-2}\mathcal S_{s,k}\) is bounded on \((0,T)\) by the datum and
source. If \(u_0=0\) and the force vanishes near \(t=0\), then
\begin{equation}
 \begin{aligned}
 I^{k-2}\mathcal S_{s,k}
 &=W_s'-2\nu W_{s+1}
 +\sum_{j=2}^{k-1}(-2\nu)^jI^{j-1}W_{s+j},\\
 \norm{I^{k-2}\mathcal S_{s,k}}_\infty
 &\le\mathsf L_s+
 \sum_{j=1}^{k-1}\frac{(2\nu)^jT^{j-1}}{(j-1)!}\mathsf B_{s+j}.
 \end{aligned}
 \label{r:eq:source-primitive-rest}
\end{equation}
No source amplitude restriction or higher velocity estimate is used.
\end{proposition}

\begin{proof}
Write \(D=\partial_t\) and \(m=k-2\). Repeated integration gives
\begin{equation}
 I^mD^dW=
 \begin{cases}
 D^{d-m}W-\displaystyle\sum_{\ell=0}^{m-1}
 W^{(d-m+\ell)}(0)\dfrac{t^\ell}{\ell!},&d\ge m,\\[5pt]
 I^{m-d}W-\displaystyle\sum_{\ell=0}^{d-1}
 W^{(\ell)}(0)\dfrac{t^{m-d+\ell}}{(m-d+\ell)!},&d<m.
 \end{cases}
 \label{r:eq:primitive-initial-terms}
\end{equation}
Empty sums vanish. In summand \(j\) of \(\mathcal S_{s,k}\), set
\(d=k-1-j=m+1-j\). The cases \(j=0,1\) leave \(W_s'\) and
\(-2\nu W_{s+1}\); the other cases leave the integrals in
\eqref{r:eq:source-primitive-rest}. Each initial correction is a polynomial
whose coefficients are determined by
\[
 W_s^{(\ell)}(0)=\sum_{b=0}^{\ell}\binom\ell b
 \ip{a_b}{\partial_t^{\ell-b}h_s(0)}.
\]
They are finite by \eqref{r:eq:initial}. Equation
\eqref{r:eq:primitive-initial-terms} specifies every correction for general
datum. The inequality \(\norm{I^aW}_\infty\le T^a\norm{W}_\infty/a!\)
proves the bound. For rest datum and forcing that vanishes initially,
local uniqueness gives a zero initial velocity history. All initial
corrections vanish, yielding \eqref{r:eq:source-primitive-rest}.
\end{proof}

The first instance bounds \(W_s'-2\nu W_{s+1}\) directly. The next
bounds the time integral of
\(W_s''-2\nu W_{s+1}'+4\nu^2W_{s+2}\). At higher orders the bound
uses exactly \(k-2\) integrations. It controls the signed source
contribution in that coordinate, rather than the absolute integral of
each differentiated term. This specifies which part of the completed
height calculation is bounded before using the velocity intersection.

\subsection{Source absorption in every temporal row}
\label{r:sec:all-row-source-absorption}

The kinetic estimate above bounds the pairing with \(u\). At a higher
temporal row, source work pairs the prescribed \(F_r\) with the response
\(V_r\). Viscous dissipation controls one additional spatial derivative
of \(V_r\). Pairing that derivative with one lower source derivative
allows part of the dissipation to absorb the source work.
Since \(V_r\) is divergence free,
\begin{equation}
 \begin{aligned}
 W_{r,s}
 &=\operatorname{Re}\ip{\Lambda^sV_r}{\Lambda^s\Pp F_r}\\
 &=\operatorname{Re}\ip{\Lambda^{s+1}V_r}
                         {\Lambda^{s-1}\Pp F_r}.
 \end{aligned}
 \label{r:eq:all-row-source-duality}
\end{equation}
For the prescribed rapidly decreasing force, the second factor is finite
at every \(s\ge0\). If \(s<1\), its Fourier weight has a singularity
at zero, but \(\widehat F_r\) is bounded there and the three-dimensional
radial integral is \(\int_0^1\rho^{2s}\dd\rho<\infty\).
The smoothness of the force controls the high frequencies. All these
bounds are uniform on the prescribed time interval.

Young's inequality now allocates one quarter of the viscous dissipation
to source work:
\begin{equation}
 \begin{aligned}
 |W_{r,s}|
 &\le\frac\nu4\norm{\Lambda^{s+1}V_r}_2^2
       +\frac1\nu\norm{\Lambda^{s-1}\Pp F_r}_2^2\\
 &=\frac\nu2E_{r,s+1}
       +\frac1\nu\norm{\Lambda^{s-1}\Pp F_r}_2^2.
 \end{aligned}
 \label{r:eq:all-row-source-absorption}
\end{equation}
Substitution into the complete height identity gives
\begin{equation}
 \partial_tE_{r,s}+\frac{3\nu}{2}E_{r,s+1}
 \le\mathcal P_{r,s}
       +\frac1\nu\norm{\Lambda^{s-1}\Pp F_r}_2^2.
 \label{r:eq:all-row-source-retained-production}
\end{equation}
The nonlinear production retains its sign and its full binomial sum.
The displayed source norm has a finite time integral determined by the
prescribed force. An absolute integral of \(W_{r,s}\) also requires
control of the dissipation appearing in
\eqref{r:eq:all-row-source-absorption}.

The argument can be written for the broader source class
\(\bigcap_{r,m}H^r_tH^m_x\), without a spatial-decay assumption.
Use \(J=(1-\Delta)^{1/2}\) and set
\[
 A_{r,s}=\tfrac12\norm{J^sV_r}_2^2,
 \qquad D_{r,s}=\norm{\nabla J^sV_r}_2^2,
 \qquad
 \mathcal P^J_{r,s}=-\operatorname{Re}\ip{J^sV_r}{J^sB_r}.
\]
Pressure again pairs to zero. The identity
\(\norm{J^{s+1}V_r}_2^2=D_{r,s}+2A_{r,s}\)
and Young's inequality yield
\begin{equation}
 A'_{r,s}+\frac\nu2D_{r,s}
 \le\mathcal P^J_{r,s}+\nu A_{r,s}
       +\frac1{2\nu}\norm{J^{s-1}\Pp F_r}_2^2.
 \label{r:eq:inhomogeneous-source-absorption}
\end{equation}
For \(s<1\), the multiplier defining \(J^{s-1}\) is bounded at
zero. Multiplying by \(e^{-\nu t}\) and integrating controls the
source contribution on every strict window with constants independent
of the time cutoff. The resulting inequality still contains the
signed integral of \(\mathcal P^J_{r,s}\).

\subsection{The source primitive at every temporal row}
\label{r:sec:all-row-primitives}

For \(h_s=\Lambda^{2s}\Pp f\), the kinetic bound also controls
the full direct source current at every finite temporal row.
Set

\[
M_{s,k}(t)=\operatorname{Re}\langle u(t),\partial_t^kh_s(t)\rangle.
\]

The exact binomial transfer is

\[
W_{r,s}
=\sum_{b=0}^r(-1)^b\binom rb
\partial_t^{r-b}M_{s,r+b}.
\]

To verify all indices and signs, expand the derivative on the right. The coefficient of
$\langle V_a,\partial_t^{2r-a}h_s\rangle$ is

\[
\sum_{b=0}^{r-a}(-1)^b\binom rb\binom{r-b}{a}
=\binom ra\sum_{b=0}^{r-a}(-1)^b\binom{r-a}{b}.
\]

It is zero for $a<r$ and one for $a=r$. The surviving term is
$\langle V_r,\partial_t^rh_s\rangle=W_{r,s}$.

The momentum calculation in \eqref{r:eq:source-first} also applies to
every prescribed test field \(g_{s,k}=\partial_t^kh_s\). It gives
\begin{equation}
\begin{aligned}
 M_{s,k}'={}&\operatorname{Re}\ip{u}{(\partial_t+\nu\Delta)g_{s,k}}\\
 &+\int_{\R^3}u_i u_j(\operatorname{Sym}\nabla g_{s,k})_{ij}\dd x
 +\operatorname{Re}\ip{f}{g_{s,k}}.
 \label{r:eq:all-row-test-derivative}
\end{aligned}
\end{equation}
Thus both the test pairing and its first time derivative are bounded
by kinetic energy and prescribed source norms:
\begin{equation}
\begin{aligned}
 \norm{M_{s,k}}_\infty&\le\mathsf B_{s,k}
 :=K_T\norm{g_{s,k}}_{L^\infty_tL^2_x},\\
 \norm{M_{s,k}'}_\infty&\le\mathsf L_{s,k}
 :=K_T\norm{(\partial_t+\nu\Delta)g_{s,k}}_{L^\infty_tL^2_x}\\
 &\quad+K_T^2\norm{\operatorname{Sym}\nabla g_{s,k}}_{L^\infty_{t,x},\mathrm{op}}
 +\norm{f}_{L^\infty_tL^2_x}\norm{g_{s,k}}_{L^\infty_tL^2_x}.
 \label{r:eq:all-row-test-bounds}
\end{aligned}
\end{equation}
The additional derivative bound removes two integrations from the
source primitive at every temporal row. Write the complete source
correction as
\[
 \mathcal F_{r,s,n}:=\sum_{j=0}^{n-1}(-2\nu)^j
 \partial_t^{n-1-j}W_{r,s+j},\qquad n\ge1.
\]

\begin{proposition}[Source correction with two fewer integrations]
\label{r:prop:all-row-sharp-primitive}
For every fixed \(r\ge0\), \(n\ge1\), and \(s\ge0\), the complete
source correction satisfies
\begin{equation}
 \norm{I^{\max(n+r-2,0)}\mathcal F_{r,s,n}}_{L^\infty(0,T)}<\infty.
 \label{r:eq:all-row-sharp-primitive}
\end{equation}
The bound uses \(\nu,T\), the kinetic bound, prescribed force norms,
and finitely many initial jets. No restriction on the force amplitude
or higher velocity norm is required.
\end{proposition}

\begin{proof}
Substitution of the binomial transfer into the complete source
correction gives
\begin{equation}
 \mathcal F_{r,s,n}
 =\sum_{j=0}^{n-1}\sum_{b=0}^r(-2\nu)^j(-1)^b\binom rb
 \partial_t^{n+r-1-j-b}M_{s+j,r+b}.
 \label{r:eq:all-row-source-transfer}
\end{equation}
When \((r,n)=(0,1)\), this is \(M_{s,0}\), already bounded by
\eqref{r:eq:all-row-test-bounds}. Otherwise put \(q=n+r-2\). For
\(0\le v\le q+1\), repeated integration gives
\begin{equation}
 I^q\partial_t^{q+1-v}M=\mathcal J_{q,v}[M],\qquad
 \mathcal J_{q,v}[M]=
 \begin{cases}
 M'-\displaystyle\sum_{\ell=0}^{q-1}
 M^{(\ell+1)}(0)\dfrac{t^\ell}{\ell!},&v=0,\\[7pt]
 M-\displaystyle\sum_{\ell=0}^{q-1}
 M^{(\ell)}(0)\dfrac{t^\ell}{\ell!},&v=1,\\[7pt]
 I^{v-1}M-\displaystyle\sum_{\ell=0}^{q-v}
 M^{(\ell)}(0)\dfrac{t^{v-1+\ell}}{(v-1+\ell)!},&v\ge2.
 \end{cases}
 \label{r:eq:all-row-sharp-integration}
\end{equation}
Empty sums vanish. In summand \((j,b)\) of
\eqref{r:eq:all-row-source-transfer}, take \(v=j+b\). The first two
branches are bounded by \(\mathsf L_{s+j,r+b}\) and
\(\mathsf B_{s+j,r+b}\), together with their initial polynomials.
For the third branch,
\[
 \norm{I^{v-1}M_{s,k}}_\infty
 \le\frac{T^{v-1}}{(v-1)!}\mathsf B_{s,k}.
\]
Every initial coefficient is finite: the product rule writes it as
\[
 M_{s,k}^{(\ell)}(0)
 =\sum_{c=0}^{\ell}\binom\ell c
 \operatorname{Re}\ip{V_c(0)}{g_{s,k+\ell-c}(0)},
\]
and \eqref{r:eq:initial} determines each \(V_c(0)\) from the smooth
datum and prescribed force jet. There are finitely many summands for
each fixed \((r,s,n)\), proving \eqref{r:eq:all-row-sharp-primitive}.
\end{proof}

For the nonnegative integrated height identity below, it remains useful
to write the source correction after \(n+r\) integrations explicitly.
Apply $I^{n+r}$ to \(\mathcal F_{r,s,n}\). For $0\le j<n$ and $0\le b\le r$, put $q=j+b+1$. The integration formula \eqref{r:eq:primitive-initial-terms}, with total integration order $n+r$, gives

\[
\begin{aligned}
&I^{n+r}
\partial_t^{n-1-j+r-b}M_{s+j,r+b}(t)\\
&\quad=I^qM_{s+j,r+b}(t)
-\sum_{\ell=0}^{n+r-2-j-b}
M_{s+j,r+b}^{(\ell)}(0)
\frac{t^{q+\ell}}{(q+\ell)!}.
\end{aligned}
\]

An upper index of \(-1\) means an empty sum. Multiplying the preceding identity by
$(-2\nu)^j(-1)^b\binom rb$ and summing gives the entire source current, with no terminal derivative.

\Needspace{10\baselineskip}
For an explicit bound, set

\[
S_{s,k}=\sup_{0\le t\le T}\|\partial_t^kh_s(t)\|_2.
\]

Then $|M_{s,k}|\le K_TS_{s,k}$ and
\[
\sup_{t<T}|I^{n+r}\mathcal F_{r,s,n}(t)|\le B_{r,s,n},
\]
where

\[
\begin{aligned}
B_{r,s,n}:=
\sum_{j=0}^{n-1}\sum_{b=0}^r
(2\nu)^j\binom rb
\Bigg[
&K_TS_{s+j,r+b}\frac{T^{j+b+1}}{(j+b+1)!}\\
&+\sum_{\ell=0}^{n+r-2-j-b}
|M_{s+j,r+b}^{(\ell)}(0)|
\frac{T^{j+b+1+\ell}}{(j+b+1+\ell)!}
\Bigg]<\infty .
\end{aligned}
\]

Every initial jet is finite by the forced recurrence and smooth initial data. This establishes finite integrated direct forcing in every temporal row and at every spatial height, separately for each finite $r,s,n$.

Write
\(\mathcal N_{r,s,n}=\sum_{j=0}^{n-1}(-2\nu)^j
\partial_t^{n-1-j}\mathcal P_{r,s+j}\).
For odd \(n\), integrating \eqref{r:eq:completed} $n+r$ times gives

\[
\begin{aligned}
I^rE_{r,s}
+(2\nu)^nI^{n+r}E_{r,s+n}
={}&I^rT_{r,s,n}
+I^{n+r}\mathcal N_{r,s,n}\\
&+I^{n+r}\mathcal F_{r,s,n},
\end{aligned}
\]

where
$T_{r,s,n}(t)=\sum_{\ell=0}^{n-1}E_{r,s}^{(\ell)}(0)t^\ell/\ell!$.
Both left-hand terms are nonnegative. Thus divergence of $I^rE_{r,s}$ requires an unbounded upper value of $I^{n+r}\mathcal N_{r,s,n}$ at every odd $n$. For $r>0$, this conclusion requires divergence of the indicated time integral; pointwise divergence of $E_{r,s}$ alone is insufficient.

\subsection{Source work in the first two temporal rows}
\label{r:app:recovered-source-work}
The source work has a useful bound before any higher velocity rectangle is
read off. This includes mixed spatial derivatives and the first temporal
row. Keep the actual forced velocity, set \(V_r=\partial_t^ru\) and
\(F_r=\partial_t^rf\), and write
\[
 U_*:=\|u_0\|_2+\int_0^T\|f(t)\|_2\dd t.
\]
The kinetic identity gives \(\sup_{t<T}\|u(t)\|_2\le U_*\).
For a spatial multi-index \(\alpha\) and a real height \(s\ge0\), introduce
the nonnegative self-adjoint multiplier
\[
 M_{\alpha,s}=(-1)^{|\alpha|}\partial_x^{2\alpha}\Lambda^{2s},
 \qquad
 \widehat{M_{\alpha,s}v}(\xi)
   =\xi^{2\alpha}|\xi|^{2s}\widehat v(\xi).
\]
Spatial self-adjointness gives the exact source pairing
\begin{equation}
 W_{r,\alpha,s}
 :=\ip{\Lambda^s\partial^\alpha V_r}
          {\Lambda^s\partial^\alpha F_r}
 =\ip{V_r}{M_{\alpha,s}F_r}.
 \label{r:eq:recovered-mixed-source-work}
\end{equation}
All pairings in this section mean their real parts. At \(r=0\), every
spatial derivative can thus be carried by the prescribed source:
\begin{equation}
 \int_0^T|W_{0,\alpha,s}|\dd t
 \le U_*\int_0^T\|M_{\alpha,s}f\|_2\dd t<\infty.
 \label{r:eq:recovered-source-row-zero}
\end{equation}
The constant is finite at each chosen order. This estimate consumes no
viscous dissipation and requires no higher velocity norm.

For \(r=1\), let \(h=M_{\alpha,s}\Pp F_1\). This is a prescribed smooth
solenoidal field. Pairing the original momentum equation with \(h\),
moving spatial derivatives to \(h\), and using incompressibility gives
\begin{equation}
 W_{1,\alpha,s}
 =\nu\ip{u}{\Delta h}
  +\int_{\mathbb R^3}u_i u_j(\operatorname{Sym}\nabla h)_{ij}\dd x
  +\ip{f}{h}.
 \label{r:eq:recovered-source-row-one-identity}
\end{equation}
Pressure contributes zero because \(\nabla\cdot h=0\). The transport term
has the displayed positive sign after integration by parts; contraction
with \(u\otimes u\) removes the antisymmetric part of \(\nabla h\).
Consequently,
\begin{align}
 \int_0^T|W_{1,\alpha,s}|\dd t
 \le{}&\nu U_*\int_0^T\|\Delta h\|_2\dd t\notag\\
 &+U_*^2\int_0^T
       \|\operatorname{Sym}\nabla h\|_{L^\infty_x,\mathrm{op}}\dd t
   +\int_0^T\|f\|_2\|h\|_2\dd t<\infty.
 \label{r:eq:recovered-source-row-one-bound}
\end{align}
Every quantity on the right is fixed by the datum, viscosity, and source.
Sobolev boundedness of \(\Pp\), followed by an adequately high spatial
embedding, controls these norms. No spatial support or rapid-decay
property of \(\Pp f\) is required. The estimate imposes no restriction on
force amplitude.

The higher temporal rows have an exact finite reduction. For \(r\ge1\),
put
\[
 B_{r,\alpha,s}
 =\sum_{j=0}^{r-1}(-1)^j
   \ip{V_{r-1-j}}{M_{\alpha,s}F_{r+j}}.
\]
Differentiating this sum cancels adjacent terms and leaves
\[
 B_{r,\alpha,s}'
 =W_{r,\alpha,s}
  +(-1)^{r-1}\ip{u}{M_{\alpha,s}F_{2r}}.
\]
Hence
\begin{equation}
 \int_0^tW_{r,\alpha,s}\dd v
 =B_{r,\alpha,s}(t)-B_{r,\alpha,s}(0)
  +(-1)^r\int_0^t\ip{u}{M_{\alpha,s}F_{2r}}\dd v.
 \label{r:eq:recovered-source-temporal-telescope}
\end{equation}
The terminal pairings involve temporal velocity rows strictly below
\(r\). If \(K_j=\sup_{t<T}\|V_j(t)\|_2\) for \(j<r\), the explicit bound
is
\begin{align}
 \sup_{t<T}\left|\int_0^tW_{r,\alpha,s}\dd v\right|
 \le{}&2\sum_{j=0}^{r-1}K_{r-1-j}
           \sup_{t\le T}\|M_{\alpha,s}F_{r+j}(t)\|_2\notag\\
 &+U_*\int_0^T\|M_{\alpha,s}F_{2r}\|_2\dd t.
 \label{r:eq:recovered-source-temporal-bound}
\end{align}
At \(r=1\), only the kinetic bound is needed. At larger \(r\), the stated
lower temporal norms enter explicitly; the datum-generated intersection
supplies them in the principal argument. Formula
\eqref{r:eq:recovered-source-temporal-bound} controls a signed integral,
while \eqref{r:eq:recovered-source-row-one-bound} controls its absolute
integral at the first temporal row.

For completeness, with
\(E_{r,\alpha,s}=\|\Lambda^s\partial^\alpha V_r\|_2^2/2\),
the corresponding energy identity can be written
\[
 \partial_t(E_{r,\alpha,s}-B_{r,\alpha,s})
 =-\nu\|\Lambda^{s+1}\partial^\alpha V_r\|_2^2
   +\mathcal P_{r,\alpha,s}
   +(-1)^r\ip{u}{M_{\alpha,s}F_{2r}}.
\]
Here \(\mathcal P_{r,\alpha,s}\) retains the complete differentiated
nonlinear row and its pressure completion. This identity changes the
recorded scalar by an explicitly bounded pairing when the lower rows are
bounded. It leaves the underlying velocity and its equation intact.

\subsection{Source insertions in the velocity rows}
Write
\[
g=\mathbb Pf,\qquad g_j=\partial_t^jg,\qquad
L=\nu\Delta,\qquad
B(v,w)=\mathbb P((v\cdot\nabla)w),
\]
\[
Q(v,w)=B(v,w)+B(w,v),\qquad
\mathcal K(u)=Lu-B(u,u),\qquad
\mathcal A_u h=Lh-Q(u,h).
\]
Here \(\mathcal K\) is the spatial differential expression in the
projected momentum equation, and \(\mathcal A_u\) is its derivative
with respect to the velocity argument. The Leray projection is
bounded in every Sobolev norm used below.

The momentum equation is \(V_1=\mathcal K(u)+g\). To differentiate it correctly, use
\[
\partial_t\mathcal K(u)=\mathcal A_uV_1,\qquad
\partial_t(\mathcal A_u h)=\mathcal A_u\partial_t h-Q(V_1,h).
\]
In particular, the coefficients of \(\mathcal A_u\) also change with time. Carrying this term gives
\[
\begin{aligned}
V_2&=\mathcal A_u\mathcal K(u)+S_2,\\
S_2&=\mathcal A_u g+g_1
=Lg-Q(u,g)+g_1,
\end{aligned}
\]
and then
\[
\begin{aligned}
V_3&=\mathcal A_u^2\mathcal K(u)
     -Q(\mathcal K(u),\mathcal K(u))+S_3,\\
S_3&=\mathcal A_u^2g+\mathcal A_u g_1
     -2Q(\mathcal K(u),g)-Q(g,g)+g_2.
\end{aligned}
\]
These are exact expansions of the full binomial rows. The coefficient two in the formula for \(S_3\) comes from differentiating both the velocity-dependent operator and the velocity argument. The quantities \(S_2,S_3\) collect source insertions into the velocity jet. They contain the prescribed source together with the velocity on which it acts.

The kinetic-energy inequality already bounds the first insertion
through the terminal time. Put
\[
U=\|u_0\|_2+\|f\|_{L^1_tL^2_x},\qquad
D^2=\nu^{-1}\left(\tfrac12\|u_0\|_2^2
                      +U\|f\|_{L^1_tL^2_x}\right).
\]
The actual forced history satisfies
\(\sup_{t<T}\|u(t)\|_2\le U\) and
\(\|\nabla u\|_{L^2_tL^2_x}\le D\).
Since \(B\) is followed by an \(L^2\)-contractive projection, the formula for \(S_2\) gives
\[
\begin{aligned}
\|S_2\|_{L^2_tL^2_x}
\le{}&\nu\|\Delta g\|_{L^2_tL^2_x}
      +\|g_1\|_{L^2_tL^2_x}\\
&+\|g\|_{L^\infty_{t,x}}D
 +\|\nabla g\|_{L^\infty_{t,x}}\sqrt{T}\,U
<\infty.
\end{aligned}
\]
Thus this first force--velocity insertion has a terminal bound independent of the critical velocity rectangle.

The next insertion admits an \(L^2_tH^{-3}_x\) estimate from
the kinetic bounds. The weaker spatial norm accommodates the
additional derivatives in this row. Define
\[
Z=\|u\|_{L^2_tH^1_x}\le(TU^2+D^2)^{1/2},\qquad
G=\max_{0\le j\le2}\sup_{t\le T}\|g_j(t)\|_{H^6}.
\]
The divergence form and Sobolev duality give
\[
\begin{aligned}
\|Q(u,g)\|_{L^2L^2}&\le C GZ,\\
\|B(u,u)\|_{L^2H^{-3/2}}&\le C UZ,\\
\|Q(u,Q(u,g))\|_{L^2H^{-3}}&\le C G UZ,\\
\|Q(Lu,g)\|_{L^2H^{-2}}&\le C\nu GZ,\\
\|Q(B(u,u),g)\|_{L^2H^{-5/2}}&\le C G UZ.
\end{aligned}
\]
For the second line use \(u\otimes u\in L^2_tL^{3/2}_x\), obtained from \(u\in L^\infty_tL^2_x\cap L^2_tL^6_x\), and the dual of \(H^{1/2}\hookrightarrow L^3\). For the third line, \(u\otimes Q(u,g)\) belongs to \(L^2_tL^1_x\); pairing its divergence against \(H^3\) is bounded because \(H^2\hookrightarrow L^\infty\). Smooth multiplication by \(g\) and its derivatives gives the final two lines. Expanding \(\mathcal A_u^2g\) in the formula for \(S_3\) now gives
\[
\|S_3\|_{L^2_tH^{-3}_x}
\le C\left[
\sqrt{T}\bigl((1+\nu+\nu^2)G+G^2\bigr)
+G(1+\nu+U)Z\right]<\infty.
\]
The constant \(C\) depends on fixed Sobolev product and multiplier constants. All velocity dependence is contained in the kinetic quantities \(U\) and \(Z\). The two estimates control \(S_2\) and \(S_3\) at the stated spatial heights.

The all-order version is also explicit. Define spatial jet polynomials by
\[
\mathcal J_0(v)=v,\qquad
\mathcal J_{n+1}(v)=D\mathcal J_n(v)[\mathcal K(v)].
\]
Here \(D\mathcal J_n(v)[h]\) means the derivative of the spatial polynomial with respect to \(v\) in direction \(h\). Along the actual forced history,
\[
V_n=\mathcal J_n(u)+S_n,\qquad S_0=0,
\]
\[
S_{n+1}=D\mathcal J_n(u)[g]+\partial_t S_n.
\]
This follows by the chain rule and \(u_t=\mathcal K(u)+g\). Every \(S_n\) is a finite differential expression with at least one source factor. The recurrence retains all velocity factors as well. The bounds above apply at the specified spatial heights.


\subsection{Time-integrated completion and terminal integrability}

To apply Proposition~\ref{r:prop:source-primitive} to the
completed-height identity \eqref{r:eq:completed}, take
\(r=0\), \(k\ge2\), \(m=k-2\), and \(H=E_{0,1/2}\). Define
\[
 \mathcal N_{s,k}:=\sum_{j=0}^{k-1}(-2\nu)^j
 \partial_t^{k-1-j}\mathcal P_{0,s+j}.
\]
Applying \(I^m\) to the full completion gives
\begin{equation}
 \begin{aligned}
 (-2\nu)^kI^m C^f_{0,k}
 &=H''-\sum_{\ell=0}^{m-1}H^{(\ell+2)}(0)\frac{t^\ell}{\ell!}\\
 &\quad-I^m\mathcal N_{1/2,k}-I^m\mathcal S_{1/2,k}.
 \end{aligned}
 \label{r:eq:integrated-completion}
\end{equation}
The source term on the last line has the proved finite bound. The energy
derivative and complete nonlinear term remain in the expression; a bound
on their combination has not been obtained from that source bound.

A full-interval Sobolev coordinate also specifies an integrability
requirement. For example, \(V_r\in L^2(0,T;\dot H^{k+1/2})\) means
\[
 \int_0^T C^f_{r,k}(t)\dd t
 =\frac12\int_0^T\norm{V_r(t)}_{\dot H^{k+1/2}}^2\dd t<\infty.
\]
This is the meaning of membership at that height. It is stronger than
the existence of every derivative on \([0,S]\) for each \(S<T\).
In particular, the full mixed product rule for \(W_{r,s}\) contains both
source and response factors, even when every source coordinate is finite.
For \(k\ge4\), boundedness of \(I^{k-2}C^f_{0,k}\) would itself give
only a weighted terminal integral,
\[
 \int_0^T(T-t)^{k-3}C^f_{0,k}(t)\dd t<\infty.
\]
The positive Volterra kernel gives this statement by monotone convergence;
its vanishing terminal weight does not give the unweighted integral above.
At \(k=2,3\), control of the full right side of
\eqref{r:eq:integrated-completion} would suffice. These distinctions specify which terminal integrals follow from
the integrated source estimate.

\section{Stokes lifting and differentiated source interactions}\label{r:joined-sec:unrestricted}

The prescribed force determines a linear Stokes response with
finite norms at every fixed order. Separating this response from the
velocity makes the additional nonlinear interactions explicit.
Let $z$ solve
\begin{equation}\label{r:joined-eq:stokeslift}
z_t-\nu\Delta z=\PP f,\qquad z(0)=0,
\qquad z(t)=\int_0^t e^{\nu(t-s)\Delta}\PP f(s)\,ds.
\end{equation}
The heat semigroup is a contraction in every $H^q$. For $Z_n=\partial_t^nz$,
\begin{equation}\label{r:joined-eq:liftjets}
\begin{aligned}
Z_{n+1}&=\nu\Delta Z_n+\PP F_n,\\
Z_n(0)&=\sum_{j=0}^{n-1}(\nu\Delta)^{n-1-j}\PP F_j(0)
\quad(n\ge1),\\
\sup_{t\le T}\norm{Z_n(t)}_{H^q}
&\le\norm{Z_n(0)}_{H^q}+\int_0^T\norm{F_n(s)}_{H^q}\,ds.
\end{aligned}
\end{equation}
Every norm in this estimate depends on finitely many derivatives of
the prescribed force. Set \(v=u-z\). This decomposition separates the
known Stokes response from the remainder of the forced solution.

With $B(a,b)=\PP[(a\cdot\nabla)b]$, the equation for $v$ is
\begin{equation}\label{r:joined-eq:liftequation}
v_t-\nu\Delta v+B(v,v)+B(z,v)+B(v,z)+B(z,z)=0.
\end{equation}
Put $W_n=\partial_t^nv$, so $V_n=W_n+Z_n$. Every temporal row becomes
\begin{equation}\label{r:joined-eq:liftrow}
\begin{aligned}
W_{n+1}-\nu\Delta W_n
+\sum_{a=0}^n\binom na\bigl\{
&B(W_a,W_{n-a})+B(Z_a,W_{n-a})\\
&+B(W_a,Z_{n-a})+B(Z_a,Z_{n-a})\bigr\}=0.
\end{aligned}
\end{equation}
The binomial self-interaction is retained in full. The other three terms contain a prescribed lift factor. If $\nabla\phi=\QQ f$ and $\pi=p-\phi$, pressure is retained by
\begin{equation}\label{r:joined-eq:liftpressure}
-\Delta\partial_t^n\pi
=\sum_{a=0}^n\binom na
\partial_i(W_a+Z_a)_j\,
\partial_j(W_{n-a}+Z_{n-a})_i.
\end{equation}

\begin{theorem}[Finite-order estimate for the Stokes decomposition]\label{r:joined-thm:lift}
For a fixed finite $N$ and real $s\ge0$, define
\begin{equation}\label{r:joined-eq:liftenergy}
\begin{aligned}
\mathscr E&=\tfrac12\sum_{n=0}^N\norm{W_n}_{H^s}^2,
&\mathscr D&=\sum_{n=0}^N\norm{\nabla W_n}_{H^s}^2,\\
\mathscr Z&=\sum_{a=0}^N\binom Na\norm{Z_a}_{H^{s+2}},
&G_n&=\sum_{a=0}^n\binom na B(Z_a,Z_{n-a}),\\
\mathscr G^2&=\sum_{n=0}^N\norm{G_n}_{H^{s-1}}^2.
\end{aligned}
\end{equation}
Let the signed self-production be
\begin{equation}\label{r:joined-eq:selfproducer}
\mathscr P_{N,s}(W)
=-\sum_{n=0}^N\ip{W_n}{
 \sum_{a=0}^n\binom na B(W_a,W_{n-a})}_{H^s}.
\end{equation}
On every strict classical interval,
\begin{equation}\label{r:joined-eq:unrestrictedrowestimate}
\mathscr E'+\frac\nu2\mathscr D
\le\mathscr P_{N,s}(W)
+\left(\nu+\frac{8C_s^2}{\nu}\mathscr Z^2\right)\mathscr E
+\frac{\mathscr G^2}{\nu},
\end{equation}
where $C_s=2^{s/2}\pi(2\pi)^{-3/2}$. The two added coefficients have finite time integrals determined entirely by the prescribed force. No smallness condition is used.
\end{theorem}
\begin{proof}
The Fourier weight inequality
$\langle\xi\rangle^s\le2^{s/2}\langle\eta\rangle^s
\langle\xi-\eta\rangle^s$
and Cauchy--Schwarz give
\begin{equation*}
\norm{ab}_{H^s}\le C_s\norm a_{H^{s+2}}\norm b_{H^s},
\end{equation*}
using $\int\langle\eta\rangle^{-4}d\eta=\pi^2$.
Since both factors are divergence-free, the divergence form of $B$ gives
$\norm{B(a,b)}_{H^{s-1}}\le\norm{a\otimes b}_{H^s}$.
In either cross term, the factor placed at height $s+2$ is $Z$.

Let $L_n$ be the sum of the two cross interactions in \eqref{r:joined-eq:liftrow}. Discrete convolution and $\binom na\le\binom Na$ give
\begin{equation*}
\left(\sum_{n=0}^N\norm{L_n}_{H^{s-1}}^2\right)^{1/2}
\le2C_s\mathscr Z\sqrt{2\mathscr E}.
\end{equation*}
Pairing through heights $s-1$ and $s+1$, all added terms are bounded by
\begin{equation*}
\sqrt{\mathscr D+2\mathscr E}
\left(2C_s\mathscr Z\sqrt{2\mathscr E}+\mathscr G\right).
\end{equation*}
Young's inequality bounds this by
$\frac\nu2\mathscr D+\nu\mathscr E
+8C_s^2\mathscr Z^2\mathscr E/\nu+\mathscr G^2/\nu$.
The remaining pairing is exactly \eqref{r:joined-eq:selfproducer}, proving
\eqref{r:joined-eq:unrestrictedrowestimate}. Finally,
$\mathscr G\le C_s\mathscr Z(\sum_{n\le N}\norm{Z_n}_{H^s}^2)^{1/2}$,
so \eqref{r:joined-eq:liftjets} bounds every added coefficient.
\end{proof}

The added interactions have coefficients determined by the source,
and their contribution is at most linear in \(\mathscr E\).
At \(N=s=0\), incompressibility also cancels the self-production.
At higher orders, the signed nonlinear term in
\eqref{r:joined-eq:selfproducer} remains. The estimate controls the
source-dependent perturbation of the energy equation without giving
an independent bound for that nonlinear term.

The decomposition must also retain the original initial derivatives.
For \(a_n=V_n(0)\), the recurrence \eqref{r:eq:initial} reads
\begin{equation}\label{r:joined-eq:initialjet}
a_{n+1}=\nu\Delta a_n
-\PP\sum_{b=0}^{n}\binom nb(a_b\cdot\nabla)a_{n-b}
+\PP F_n(0).
\end{equation}
Thus the initial derivatives remain fixed by \(u_0\) and the
source derivatives at zero. Full-interval regularity for arbitrary
forcing is supplied by Corollary~\ref{r:thm:supplied-forced}, as a direct
consequence of the global forced theorem proved in Part~II.


\subsection{The force inside differentiated nonlinear production}

The explicit work term does not contain every occurrence of the source after time differentiation. The nonlinear production is evaluated on a forced velocity, so differentiating it inserts the force into each velocity slot. At the second critical-height row this entire insertion can also be controlled from the kinetic budget.

Write $E_s(u)=\frac12\norm u_{\dot H^s}^2$, $g=\PP f$,
$\mathcal A(u)=\nu\Delta u-B(u,u)$ and
$\mathcal P_s(u)=-\langle\Lambda^su,\Lambda^sB(u,u)\rangle$.
The vector field $\mathcal A$ is the unforced right side evaluated at the current velocity; the actual trajectory obeys $u_t=\mathcal A(u)+g$.
Let $\mathcal J_s=D\mathcal P_s(u)[g]$. Differentiating its three velocity slots gives
\begin{equation}\label{r:joined-eq:sourceinsertion}
\begin{aligned}
\mathcal J_s={}&-\langle\Lambda^sg,\Lambda^s((u\cdot\nabla)u)\rangle\\
&-\langle\Lambda^su,\Lambda^s((g\cdot\nabla)u)\rangle\\
&-\langle\Lambda^su,\Lambda^s((u\cdot\nabla)g)\rangle .
\end{aligned}
\end{equation}
The projection has been removed only inside global pairings with solenoidal fields.

Define $U=\sup_{t<T}\norm{u(t)}_2$,
$D=\norm{\nabla u}_{L^2(0,T;L^2)}$, and
\begin{equation*}
a_\rho(t)=(2\pi)^{-3/2}\int_{\R^3}|\xi|^\rho
|\widehat g(t,\xi)|\,d\xi,\qquad
h(t)=\norm{u(t)}_{\dot H^{1/2}}.
\end{equation*}
Kinetic energy bounds $U$ and $D$ from the prescribed data. Every $a_\rho$ used below is bounded by a sufficiently high source Sobolev norm.

\begin{proposition}[Complete critical source insertion]\label{r:joined-prop:sourceinsertion}
Without a source smallness condition,
\begin{equation}\label{r:joined-eq:insertionbound}
|\mathcal J_{1/2}|\le
2a_1h^2+a_{3/2}Uh+a_2U^2.
\end{equation}
In particular, with $a_\rho^*=\sup_{t\le T}a_\rho(t)$,
\begin{equation}\label{r:joined-eq:insertionL1}
\norm{\mathcal J_{1/2}}_{L^1(0,T)}
\le2a_1^*U\sqrt T\,D
+a_{3/2}^*U^{3/2}T^{3/4}D^{1/2}
+a_2^*TU^2.
\end{equation}
\end{proposition}
\begin{proof}
In the first pairing of \eqref{r:joined-eq:sourceinsertion}, move all spatial derivatives onto $g$; its absolute value is bounded by $a_2U^2$.
In the second pairing the undifferentiated transport cancels by $\diver g=0$. Its commutator obeys
\begin{equation*}
\norm{[\Lambda^{1/2},g\cdot\nabla]u}_2\le a_1h,
\end{equation*}
because
$\bigl||\xi|^{1/2}-|\eta|^{1/2}\bigr|\,|\eta|
\le|\xi-\eta|\,|\eta|^{1/2}$
in its Fourier kernel. Young's convolution inequality proves the bound.
For the final pairing, splitting the half-derivative Fourier weight gives
\begin{equation*}
\norm{\Lambda^{1/2}((u\cdot\nabla)g)}_2
\le a_1h+a_{3/2}U.
\end{equation*}
Summing gives \eqref{r:joined-eq:insertionbound}. Finally, Fourier Cauchy--Schwarz gives $h^2\le U\norm{\nabla u}_2$. Integrating this inequality and its square root proves \eqref{r:joined-eq:insertionL1}.
\end{proof}

The direct work terms can also be estimated using only kinetic energy and source norms. Set
$W_s=\langle u,\Lambda^{2s}g\rangle$ and $k_s=\Lambda^{2s}g$.
Moving derivatives onto the prescribed test field yields
\begin{equation}\label{r:joined-eq:workderivative}
\begin{aligned}
W_s'={}&\langle u,(\partial_t+\nu\Delta)k_s\rangle\\
&+\int_{\R^3}u_i u_j\,\partial_j(k_s)_i\,dx
+\langle g,k_s\rangle .
\end{aligned}
\end{equation}
Each term is uniformly bounded from $U$ and smooth source norms. The pressure pairing vanishes because $k_s$ is divergence-free.

Let $H(u)=\frac12\norm u_{\dot H^{1/2}}^2$ and write
$L_{\mathcal A}\Phi=D\Phi[\mathcal A]$ for differentiation of a scalar functional along the unforced vector field. Along the forced history the exact second-height equation is
\begin{equation}\label{r:joined-eq:generatorcomparison}
\begin{aligned}
H''&=L_{\mathcal A}^2H(u)+R_f,\\
R_f&=\mathcal J_{1/2}+W_{1/2}'-2\nu W_{3/2},
\qquad R_f\in L^1(0,T).
\end{aligned}
\end{equation}
Thus the source correction hidden inside nonlinear differentiation is controlled along with the explicit work terms. The remaining functional is
\begin{equation}\label{r:joined-eq:autonomousgenerator}
L_{\mathcal A}^2H(u)
=4\nu^2E_{5/2}(u)
+D\mathcal P_{1/2}(u)[\mathcal A(u)]
-2\nu\mathcal P_{3/2}(u).
\end{equation}
The leading viscous coefficient remains positive.
The remainder depends on the nonlinear evolution of \(u\).
Subtracting the double time integral of \(R_f\) corrects the scalar
energy identity, but the velocity still satisfies
\(u_t=\mathcal A(u)+g\). Thus neither this scalar correction nor the
Stokes decomposition identifies the forced solution with an unforced
one. The estimates establish bounded source contributions at the
stated orders; full-interval control of the nonlinear remainder is
a separate assertion.


\subsection{Graph domains of the Stokes linearization}
\label{r:sec:linear-domains}

The linearization around the smooth Stokes response has the
parabolic graph domain of the heat operator.
Write \(w=z\) for the prescribed Stokes response in
\eqref{r:joined-eq:stokeslift}, and define
\[
 K_wv=\Pp[(w\cdot\nabla)v+(v\cdot\nabla)w],
 \qquad L_w=\partial_t-\nu\Delta+K_w.
\]
For an integer \(m\ge0\), let
\[
 \mathcal E_m(0,S)=
 C([0,S];H^m_\sigma)\cap L^2(0,S;H^{m+1}_\sigma),
 \qquad v_t\in L^2(0,S;H^{m-1}_\sigma).
\]
The derivative condition is included in the definition of \(\mathcal E_m\).

\begin{proposition}[Linear graph domains]
For \(0<S\le T\), the initial-trace graph completions of \(L_w\) and
\(L_0=\partial_t-\nu\Delta\) both equal \(\mathcal E_m(0,S)\).
Their norms are equivalent with constants bounded independently of \(S\):
\begin{equation}
 \begin{aligned}
 &\norm v_{CH^m}+\norm v_{L^2H^{m+1}}
       +\norm{v_t}_{L^2H^{m-1}}\\
 &\qquad\le C_{m,\nu,T,w}
 \left(\norm{v(0)}_{H^m}+\norm{L_wv}_{L^2H^{m-1}}\right).
 \end{aligned}
 \label{r:eq:linear-domain-bound}
\end{equation}
Every coefficient in the constant is determined by finitely many prescribed
source norms, with no restriction on their size.
\end{proposition}

\begin{proof}
Product differentiation and the divergence form at \(m=0\) give
\[
 \norm{K_wv}_{H^{m-1}}
 \le C_m\norm w_{W^{m+1,\infty}}\norm v_{H^m}.
\]
Put \(r=L_wv\), and pair the equation at height \(m\), distributing
the duality between \(H^{m-1}\) and \(H^{m+1}\).
Since
\(\norm{\nabla v}_{H^m}^2=\norm v_{H^{m+1}}^2-\norm v_{H^m}^2\),
Young's inequality yields
\[
 \frac{\dd}{\dd t}\norm v_{H^m}^2+\nu\norm v_{H^{m+1}}^2
 \le C_{m,\nu,w}\norm v_{H^m}^2+C_\nu\norm r_{H^{m-1}}^2.
\]
Gronwall and integration give the first two norms in
\eqref{r:eq:linear-domain-bound}; the equation gives the third.
The exponential factor is bounded by its value at \(T\).
Conversely, each member of \(\mathcal E_m\) has \(L_wv\in L^2H^{m-1}\)
by the product bound. Approximation of its initial trace and residual,
followed by the difference estimate, gives the graph completion.
Existence for the approximating linear problems follows from the heat
integral equation on a finite partition, on which its smooth bounded
coefficients give contraction. The energy estimate joins those intervals.
\end{proof}

Temporal differentiation retains the earlier lift rows. If
\(W_n=\partial_t^nv\), \(Z_a=\partial_t^aw\), and
\(h_w=-\Pp[(w\cdot\nabla)w]\), the exact equation is
\[
 L_wW_n+\sum_{a=0}^n\binom na
       \Pp[(W_a\cdot\nabla)W_{n-a}]
 =\partial_t^nh_w-\sum_{a=1}^n\binom na K_{Z_a}W_{n-a}.
\]
For \(c_{a,m}=C_m\sup_{t\le T}\norm{Z_a}_{W^{m+1,\infty}}\),
\[
 \norm{\partial_t^n(K_wv)}_{L^2H^{m-1}}
 \le\sum_{a=0}^n\binom na c_{a,m}\norm{W_{n-a}}_{L^2H^m}.
\]
The adjacent pair has parabolic degree \(2n+m+1\); each velocity
factor on this right side has degree at most \(2n+m\).
At the base corner, kinetic energy already supplies the lower row.
These source contributions preserve the terminal linear domain.
The nonlinear self-interaction remains in the displayed equation, so
the domain estimate does not identify \(v\) with an unforced solution.

\section{Critical growth, source reserve, and chord density}
\label{r:sec:critical-reserve-density}

The source work at critical order is \(\ip{u}{\Lambda g}\), where
\(g=\Pp f\). Its spatial derivatives can all be placed on \(g\), so the
kinetic bound controls how much this term can still add before a prescribed
finite time. Keeping that remaining contribution with the critical height
gives a positive quantity whose growth is governed by the nonlinear
production of the velocity.

Let \(u\) be a classical history on \([0,T)\), with the force smooth through
\(T\) as in Section~\ref{r:sec:introduction}. All calculations first take place
on strict classical intervals. Write
\begin{equation}
 \begin{aligned}
 \mathsf K_T&=\norm{u_0}_2+\int_0^T\norm{g(t)}_2\dd t,\\
 H(t)&=\tfrac12\norm{\Lambda^{1/2}u(t)}_2^2,
 &D(t)&=\norm{\Lambda^{3/2}u(t)}_2^2.
 \end{aligned}
 \label{r:crd-eq:height}
\end{equation}
The kinetic work satisfies
\[
 \int_0^t\norm{g(s)}_2\norm{u(s)}_2\dd s
 \le\int_0^t\mathsf K_s\mathsf K_s'\dd s
 =\tfrac12\bigl(\mathsf K_t^2-\norm{u_0}_2^2\bigr).
\]
The norm inequality and the energy identity give
\begin{equation}
 \norm{u(t)}_2\le\mathsf K_t,\qquad
 \norm{u(t)}_2^2+2\nu\int_0^t\norm{\nabla u(s)}_2^2\dd s
 \le\mathsf K_t^2\le\mathsf K_T^2.
 \label{r:crd-eq:kinetic}
\end{equation}
At critical order, put \(P_H=\mathcal P_{0,1/2}\) and
\(W=\operatorname{Re}\ip{u}{\Lambda g}\). The exact balance is
\[
 H'+\nu D=P_H+W,\qquad
 |W(t)|\le b_f(t):=\mathsf K_T\norm{\Lambda g(t)}_2.
\]
The source assumptions give \(b_f\in L^1(0,T)\). Its remaining integral
decreases at the rate \(b_f\), so define
\begin{equation}
 \mathscr R_f(t)=\int_t^T b_f(s)\dd s,\qquad
 \Phi=H+\mathscr R_f,\qquad Z_f=b_f-W.
 \label{r:crd-eq:reserve}
\end{equation}
Then
\begin{equation}
 \Phi'+\nu D+Z_f=P_H,\qquad
 Z_f\ge0,\qquad \Phi\ge H\ge0.
 \label{r:crd-eq:reserve-balance}
\end{equation}
The known function \(\mathscr R_f\) is absolutely continuous and tends to
zero at \(T\). Neither statement requires a terminal value of \(H\).
The production in \eqref{r:crd-eq:reserve-balance} is still evaluated on the
forced velocity. The force and its pressure contribution remain in the
momentum equation.

\subsection{Where critical growth is distributed}

The critical norm measures velocity differences over every separation.
For a center \(z\), a chord \(y\), and \(0\le\theta\le1\), put
\[
 x_+=z+(1-\theta)y,\qquad x_-=z-\theta y,\qquad
 \delta u_\theta=u(x_+)-u(x_-),\qquad \widehat y=y/|y|.
\]
The positive tensor
\begin{equation}
 \mathsf T_u(z)=\frac1{\pi^2}\int_0^1\int_{\R^3}
 \frac{|\delta u_\theta(z,y)|^2}{|y|^4}
 \widehat y\otimes\widehat y\dd y\dd\theta,\qquad
 \tau=\tr\mathsf T_u
 \label{r:crd-eq:tensor}
\end{equation}
weights each chord orientation by the squared magnitude of its velocity difference.
The fractional Fourier identity and the complete transport pairing give
\begin{equation}
 \int_{\R^3}\tau\dd z=2\norm{\Lambda^{1/2}u}_2^2=4H,\qquad
 P_H=-\int_{\R^3}S:\mathsf T_u\dd z,
 \quad S=\tfrac12(\nabla u+\nabla u^T).
 \label{r:crd-eq:current}
\end{equation}
The transport identity follows by integration in the chord variable and
averaging the strain along that chord; the local derivation below retains
the resulting redistribution flux.

Where \(\tau>0\), define \(r=-S:\mathsf T_u/\tau\), and set \(r=0\)
where \(\tau=0\). Positivity of \(\mathsf T_u\) gives
\[
 |r|\le\norm S_{\mathrm{op}}\le |S|_F,\qquad
 P_H=\int\tau r\dd z.
\]
Thus the production samples strain with the nonlocal density \(\tau\).
For \(H>0\), use the probability measure and its positive sampled rate
\[
 \dd\mu_t=\frac{\tau(z,t)\dd z}{4H(t)},\qquad
 \overline r_+(t)=\int r_+\dd\mu_t.
\]
Set \(\overline r_+=0\) when \(H=0\). Finite energy then forces \(u=0\),
so \(\tau=P_H=D=W=0\); no probability measure is assigned to zero mass.
Equation~\eqref{r:crd-eq:reserve-balance} gives
\begin{equation}
 \Phi'\le P_H\le\int\tau r_+\dd z
 =4H\overline r_+\le4\Phi\overline r_+.
 \label{r:crd-eq:sampled-growth}
\end{equation}
On an interval \([a,b]\) where \(\Phi>0\), a doubling requires
\[
 4\int_a^b\overline r_+(t)\dd t
 \ge\log\frac{\Phi(b)}{\Phi(a)}\ge\log2.
\]
The rate itself has an ordinary-volume square bound. For each eigenvalue
\(\lambda_i\) of the trace-free strain \(S\), the other two sum to
\(-\lambda_i\), so
\[
 |S|_F^2\ge\lambda_i^2+\tfrac12\lambda_i^2
 \quad\text{for every }i,\qquad
 |r|^2\le\norm S_{\mathrm{op}}^2\le\tfrac23|S|_F^2.
\]
Incompressibility gives \(\norm S_2^2=\norm{\nabla u}_2^2/2\), hence
\begin{equation}
 \int_0^T\int_{\R^3}|r|^2\dd z\dd t
 \le\frac{\mathsf K_T^2}{6\nu}.
 \label{r:crd-eq:rate-budget}
\end{equation}
The integral through \(T\) is the monotone limit of strict-subinterval
integrals.

To compare this volume bound with the density that samples it, put
\begin{equation}
 \rho_f=\frac{\tau}{4\Phi},\qquad
 V_f^{-1}(t)=\int_{\{r>0\}}\rho_f(z,t)^2\dd z,\qquad
 G_f=\{t:\Phi'(t)>0\}.
 \label{r:crd-eq:volume}
\end{equation}
These definitions apply where \(\Phi>0\); set \(\rho_f=0\) when
\(\Phi=0\). Its mass is \(H/\Phi\le1\). If
\(V_{\mathrm{work}}^{-1}=\int_{\{r>0\}}(\tau/(4H))^2\dd z\) for
\(H>0\), then
\[
 V_f^{-1}=(H/\Phi)^2V_{\mathrm{work}}^{-1}.
\]
At \(H=0<\Phi\), equation~\eqref{r:crd-eq:reserve-balance} gives
\(\Phi'=-b_f\le0\), so such times add no positive growth.

On a positive interval, discard the negative part of the logarithmic
growth and apply Cauchy--Schwarz in space and then in time:
\begin{equation}
 \log\frac{\Phi(b)}{\Phi(a)}
 \le4
 \left(\int_{[a,b]\cap G_f}\int r_+^2\dd z\dd t\right)^{1/2}
 \left(\int_{[a,b]\cap G_f}V_f^{-1}\dd t\right)^{1/2}.
 \label{r:crd-eq:growth-volume}
\end{equation}
A finite inverse-volume integral on the growth set thus bounds \(\Phi\)
from any positive initial value, and bounds \(H\) with it.
For \(N\) disjoint positive doubling intervals \(I_j\), write
\[
 A_j=\int_{I_j\cap G_f}\int r_+^2\dd z\dd t,\qquad
 C_j=\int_{I_j\cap G_f}V_f^{-1}\dd t.
\]
Then \(A_jC_j\ge(\log2)^2/16\). If \(\mathsf K_T>0\), a second
Cauchy--Schwarz inequality and \eqref{r:crd-eq:rate-budget} give
\begin{equation}
 \sum_{j=1}^N C_j
 \ge\frac{(\log2)^2N^2}{16\sum_{j=1}^N A_j}
 \ge\frac{3\nu(\log2)^2}{8\mathsf K_T^2}N^2.
 \label{r:crd-eq:repeated-growth}
\end{equation}
Infinite actions are read in the extended sense. If \(\mathsf K_T=0\),
the velocity is zero and no positive doubling occurs. Repeated growth
requires increasing concentration of the density that samples strain.
The evolution of that density must retain the endpoints of each chord.

\subsection{The local density equation and its source}

Write \(p=p_u+p_f\), with
\[
 -\Delta p_u=\partial_i u_j\,\partial_j u_i,\qquad
 \nabla p_f=(I-\Pp)f.
\]
Combining \(p_f\) with the prescribed force gives
\(D_tu=\nu\Delta u-\nabla p_u+g\). Define
\[
 a_\pm=u(x_\pm)-u(z),\qquad
 D_z=\partial_t+u(z)\cdot\nabla_z,
\]
and use the chord integral
\[
 \mathcal K[F](z)=\frac1{\pi^2}\int_0^1\int_{\R^3}
             \frac{F(z,\theta,y)}{|y|^4}\dd y\dd\theta.
\]
A fixed chord translated by the center velocity does not move with either
endpoint. Substituting the momentum equation at both endpoints gives
\begin{equation}
 \begin{aligned}
 D_z\delta u_\theta&=\nu\Delta_z\delta u_\theta
                  +\mathcal R+\delta g_\theta,\\
 \mathcal R&=-\delta\nabla p_u
 -(a_+\cdot\nabla)u(x_+)+(a_-\cdot\nabla)u(x_-).
 \end{aligned}
 \label{r:crd-eq:chord-momentum}
\end{equation}
Since \(\tau=\mathcal K[|\delta u_\theta|^2]\), the product rule yields
\begin{equation}
 \begin{aligned}
 (D_z-\nu\Delta_z)\tau
   &=\mathcal I_\tau-2\nu Q_\tau+\mathcal F_\tau,\\
 \mathcal I_\tau&=2\mathcal K[\delta u_\theta\cdot\mathcal R],\\
 Q_\tau&=\mathcal K[|\delta\nabla u_\theta|^2],\qquad
 \mathcal F_\tau=2\mathcal K[\delta u_\theta\cdot\delta g_\theta].
 \end{aligned}
 \label{r:crd-eq:density}
\end{equation}
Pressure and endpoint transport remain together in \(\mathcal I_\tau\).
The direct force term has an absolute spatial bound. Cauchy--Schwarz
over all centers, chords, and center parameters gives
\begin{equation}
 \norm{\mathcal F_\tau(t)}_1
 \le4\norm{\Lambda^{1/2}u(t)}_2\norm{\Lambda^{1/2}g(t)}_2
 =4\sqrt{2H(t)}\,\norm{\Lambda^{1/2}g(t)}_2.
 \label{r:crd-eq:local-source}
\end{equation}

The time integral can be bounded before a supremum of \(H\) is known.
For \(0<s\le1\), Fourier interpolation gives
\[
 E_s^{1/s}\le
 2^{-1/s}\norm u_2^{2(1-s)/s}\norm{\nabla u}_2^2,
 \qquad E_s=\tfrac12\norm{\Lambda^su}_2^2.
\]
Using \eqref{r:crd-eq:kinetic},
\begin{equation}
 \int_0^T E_s^{1/s}\dd t
 \le2^{-1/s-1}\nu^{-1}\mathsf K_T^{2/s},\qquad
 \int_0^T H^2\dd t\le\frac{\mathsf K_T^4}{8\nu}.
 \label{r:crd-eq:low-height}
\end{equation}
When \(\mathsf K_T=0\), both sides vanish. Applying H\"older to
\eqref{r:crd-eq:local-source} now gives
\begin{equation}
 \int_0^T\norm{\mathcal F_\tau(t)}_1\dd t
 \le2^{7/4}\mathsf K_T\nu^{-1/4}
       \norm{\Lambda^{1/2}g}_{L^{4/3}_tL^2_x}.
 \label{r:crd-eq:local-source-budget}
\end{equation}
The manuscript's smooth source class supplies this norm as well as the
\(L^1_tH^1_x\) norm used in \eqref{r:crd-eq:reserve}.
The direct chord forcing has finite absolute mass through \(T\), without
a terminal critical-height bound.

The integrated density equation has the normalization
\begin{equation}
 \int Q_\tau\dd z=2D,\qquad
 \int\mathcal F_\tau\dd z=4W,\qquad
 \int\mathcal I_\tau\dd z=4P_H.
 \label{r:crd-eq:integrated-density}
\end{equation}
For each fixed chord, \(\delta u_\theta\) is divergence free in \(z\),
so its global pairing with \(\delta\nabla p_u\) vanishes.
Integrating \eqref{r:crd-eq:density} gives
\(4H'=4P_H-4\nu D+4W\).
Pressure can still redistribute density locally. To expose that
redistribution, its flux must be retained before spatial integration.

\subsection{Intrinsic production and redistribution}

Endpoint advection acts on the chord coordinates through
\[
 c_\theta\cdot\nabla_z+\delta u_\theta\cdot\nabla_y,\qquad
 c_\theta=\theta u(x_+)+(1-\theta)u(x_-).
\]
Both coordinate divergences vanish by incompressibility. The intrinsic
integrand in \eqref{r:crd-eq:density} can be written
\[
 \nabla_z\cdot\bigl[(u(z)-c_\theta)|\delta u_\theta|^2
               -2\delta p_\theta\,\delta u_\theta\bigr]
 -\nabla_y\cdot\bigl(\delta u_\theta|\delta u_\theta|^2\bigr),
\]
where \(\delta p_\theta=p_u(x_+)-p_u(x_-)\).
The first term already has a flux in the center coordinate:
\begin{equation}
 J_0=\mathcal K\bigl[(u(z)-c_\theta)|\delta u_\theta|^2
               -2\delta p_\theta\,\delta u_\theta\bigr].
 \label{r:crd-eq:first-flux}
\end{equation}
Integrating the second term against the radial kernel produces
\(-4|\delta u_\theta|^2(\delta u_\theta\cdot\widehat y)/|y|\).
The velocity difference is the integral of its gradient along the chord.
With \(s(z,y)=\widehat y\cdot S(z)\widehat y\),
\[
 \frac{\delta u_\theta\cdot\widehat y}{|y|}
 =\int_0^1s(z+(\alpha-\theta)y,y)\dd\alpha.
\]
Changing the center parameter preserves the endpoints:
\[
 \delta u_\alpha(z+(\alpha-\theta)y,y)
 =\delta u_\theta(z,y).
\]
Put \(F_\alpha(z,y)=|\delta u_\alpha(z,y)|^2s(z,y)\).
Its translation away from \(z\) has the exact divergence form
\[
 F_\alpha(z+h,y)-F_\alpha(z,y)
 =\nabla_z\cdot
 \left[h\int_0^1F_\alpha(z+\lambda h,y)\dd\lambda\right].
\]
With \(h_{\alpha\theta}=(\alpha-\theta)y\), define
\begin{equation}
 \begin{aligned}
 L(z)={}&\frac1{\pi^2}\int_{\R^3}\frac{\dd y}{|y|^4}
 \int_0^1\int_0^1 h_{\alpha\theta}
 \left[\int_0^1F_\alpha(z+\lambda h_{\alpha\theta},y)
                         \dd\lambda\right]\dd\alpha\dd\theta,\\
 J_{\mathrm{int}}={}&J_0-4L.
 \end{aligned}
 \label{r:crd-eq:full-flux}
\end{equation}
The untranslated part is \(S:\mathsf T_u\), giving
\begin{equation}
 \mathcal I_\tau=-4S:\mathsf T_u+\nabla_z\cdot J_{\mathrm{int}}
 =4\tau r+\nabla_z\cdot J_{\mathrm{int}}.
 \label{r:crd-eq:intrinsic-split}
\end{equation}
The complete local equation is
\begin{equation}
 \partial_t\tau+\nabla_z\cdot(u\tau-J_{\mathrm{int}})
 -\nu\Delta_z\tau
 =4\tau r-2\nu Q_\tau+\mathcal F_\tau.
 \label{r:crd-eq:local-conservation}
\end{equation}
The strain at the center controls the displayed production, while
\(J_{\mathrm{int}}\) retains pressure and the change between chord centers.

These operations can first be performed with
\(\varepsilon<|y|<R\). The boundary remainder in the chord integration
obeys
\[
 \norm{B_{\varepsilon,R}}_1
 \le C\bigl(\varepsilon\norm{\nabla u}_3^3
                 +R^{-2}\norm u_3^3\bigr)\longrightarrow0.
\]
No endpoint in \(\theta\) or \(\alpha\) is differentiated.
On strict smooth time windows the paired differences also give two
powers of \(|y|\) near zero, and the radial kernel is integrable at
infinity. These bounds justify the local identities in spatial
distributions.

The flux itself is absolutely integrable in space on such a window.
Write \(U=\norm u_2\), \(A=\norm{\nabla u}_2\), and \(B=\norm u_3\).
Fractional Sobolev embedding and the Fourier kernel identity yield
\[
 \int_{\R^3}\frac{\norm{\delta_yu}_3^2}{|y|^4}\dd y
 \le C\int_{\R^3}
 \frac{\norm{\Lambda^{1/2}\delta_yu}_2^2}{|y|^4}\dd y
 \le CA^2.
\]
Since \(\norm{u-c_\theta}_3\le2B\), the transport part of \(J_0\)
has \(L^1\) norm at most \(CBA^2\).

For \(L\), split the chords at length \(\ell>0\).
The estimate
\(\norm{\delta_yu}_4^2\le C|y|^{1/2}A^2\) gives
\(C\ell^{1/2}\norm S_2A^2\) for short chords.
For long chords, average the squared endpoint velocities in the center
parameter. For an endpoint fraction \(0<a<1\), changing variables to
\(h=ay\) gives the kernel
\[
 k_\ell(h)=|h|^{-3}\min(1,|h|/\ell),\qquad
 \norm{k_\ell}_{6/5}^{6/5}=\frac{40\pi}{3}\ell^{-3/5}.
\]
H\"older and Young convolution inequalities bound this part by
\(C\ell^{-1/2}\norm S_2B^2\).
Choosing \(\ell=(B/A)^2\) for nonzero \(u\) gives
\(\norm L_1\le C\norm S_2AB\le CBA^2\).
For the zero field there is no division.

The pressure part of \(J_0\) satisfies
\[
 \norm{J_{0,\mathrm{pressure}}}_1
 \le4\norm{\Lambda^{1/2}p_u}_2\norm{\Lambda^{1/2}u}_2.
\]
The pressure multiplier is bounded, and Fourier Cauchy--Schwarz for
each product \(u_i u_j\) uses
\[
 \int_{\R^3}\frac{\dd\eta}{|\eta|^2|\xi-\eta|^2}
 =C|\xi|^{-1}.
\]
Multiplying by \(|\xi|\), integrating in \(\xi\), and applying Tonelli
gives \(\norm{\Lambda^{1/2}p_u}_2\le CA^2\).
Finally \(B\le CU^{1/2}A^{1/2}\) and
\(\norm{\Lambda^{1/2}u}_2\le U^{1/2}A^{1/2}\), so
\begin{equation}
 \norm{J_{\mathrm{int}}}_1\le CU^{1/2}A^{5/2},\qquad
 \int_0^T\norm{J_{\mathrm{int}}}_1^{4/5}\dd t
 \le C\nu^{-1}\mathsf K_T^{12/5}.
 \label{r:crd-eq:flux-budget}
\end{equation}
The time exponent is \(4/5\). A terminal \(L^1_tL^1_z\) flux estimate
does not follow from this bound.

\subsection{Diffusion of the normalized density}

Diffusion gives a local constraint on how sharply the density varies.
Cauchy--Schwarz inside the chord integral gives
\[
 |\nabla\tau|^2\le4\tau Q_\tau,\qquad
 -2\nu Q_\tau\le-\frac\nu2\frac{|\nabla\tau|^2}{\tau}
 \quad\hbox{where }\tau>0.
\]
The intrinsic production in \eqref{r:crd-eq:density} remains alongside
this favorable term. Its time integral cannot be discarded when estimating
the diffusion of \(\rho_f\).

On an interval with \(\Phi>0\), put
\[
 \psi=\sqrt{\rho_f},\qquad
 m=\int\rho_f\dd z=\frac H\Phi\le1,\qquad
 \mathcal J=\int|\nabla\psi|^2\dd z.
\]
Differentiating the normalization gives
\begin{equation}
 (D_z-\nu\Delta_z)\rho_f
 =\frac{\mathcal I_\tau+\mathcal F_\tau}{4\Phi}
 -\frac{\nu Q_\tau}{2\Phi}
 -\frac{\Phi'}{\Phi}\rho_f.
 \label{r:crd-eq:normalized-density}
\end{equation}
Where \(\tau>0\), its square root satisfies
\begin{equation}
 \begin{aligned}
 (D_z-\nu\Delta_z)\psi
 ={}&\frac{\mathcal I_\tau+\mathcal F_\tau}{8\Phi\psi}
 -\frac{\nu}{\psi}
     \left(\frac{Q_\tau}{4\Phi}-|\nabla\psi|^2\right)
 -\frac{\Phi'}{2\Phi}\psi.
 \end{aligned}
 \label{r:crd-eq:square-root}
\end{equation}
The bracket is nonnegative. The derivative of \(\Phi\) retains its full
sign, and pressure and endpoint transport remain inside \(\mathcal I_\tau\).

The norm derivative inequality extends across zeros by Sobolev approximation,
and \eqref{r:crd-eq:integrated-density} gives
\(\mathcal J\le D/(2\Phi)\).
Let \(S_6\) be a constant in
\(\norm h_6\le S_6\norm{\nabla h}_2\). Interpolation gives
\begin{equation}
 V_f^{-1}\le\norm\psi_4^4
 \le\min\left\{
 S_6^3m^{1/2}\mathcal J^{3/2},
 \frac{S_6^2\norm\tau_{3/2}\mathcal J}{4\Phi}
 \right\}.
 \label{r:crd-eq:density-interpolation}
\end{equation}
The first estimate uses the \(L^2\) and \(L^6\) norms of \(\psi\);
the second uses its \(L^3\) and \(L^6\) norms.
For the chord density,
\(\norm\tau_{3/2}\le C_\tau\norm{\nabla u}_2^2\):
Minkowski bounds this norm by
\(C\int |y|^{-4}\norm{\delta_yu}_3^2\dd y\), which is controlled by
the fractional estimate used for \(J_0\).
Substituting these bounds gives
\begin{equation}
 V_f^{-1}\le\min\left\{
 \frac{S_6^3H^{1/2}D^{3/2}}{2^{3/2}\Phi^2},
 \frac{C_\tau S_6^2\norm{\nabla u}_2^2D}{8\Phi^2}
 \right\}.
 \label{r:crd-eq:inverse-volume}
\end{equation}
Neither time integral on the right is bounded by
\eqref{r:crd-eq:kinetic} and \eqref{r:crd-eq:local-source-budget}.
Even a finite integral of \(\mathcal J\) would give only
\[
 \int(V_f^{-1})^{2/3}\dd t\le S_6^2\int\mathcal J\dd t,
\]
which has a smaller time exponent than the action in
\eqref{r:crd-eq:growth-volume}.

For this route, a sufficient additional estimate is
\[
 \int_{G_f}V_f^{-1}(t)\dd t<\infty
\]
with a bound derived from \(u_0,\nu,f\). The source reserve and the absolute
chord-source estimate control the direct external contribution.
The required intrinsic estimate must control how strain samples the density
under the complete production and redistribution in
\eqref{r:crd-eq:local-conservation}. The rate bound
\eqref{r:crd-eq:rate-budget}, the flux bound \eqref{r:crd-eq:flux-budget},
and the instantaneous diffusion estimates above do not supply that
time-integrated control.


\section{A sufficient condition from the squared momentum equation}
\label{r:sec:first-rectangle}

The temporal and viscous terms can be separated directly in an integrated
identity. Write
\[
 \mathcal G=\Pp f-\Pp[(u\cdot\nabla)u],
 \qquad u_t-\nu\Delta u=\mathcal G.
\]
For a strict classical horizon \(S<T\), put
\[
 X(S)=\int_0^S(\norm{u_t}_2^2+\nu^2\norm{\Delta u}_2^2)\dd t,
 \qquad G_0=\norm{\nabla u_0}_2^2.
\]
The cross term in the squared equation is an endpoint energy:
\begin{equation}
 X(S)+\nu\norm{\nabla u(S)}_2^2
 =\nu G_0+\int_0^S\norm{\mathcal G}_2^2\dd t.
 \label{r:eq:first-rectangle-exact}
\end{equation}
Indeed, \(-2\ip{u_t}{\Delta u}=\partial_t\norm{\nabla u}_2^2\).
This identity retains the full projected nonlinear term.

Set
\[
 U=\norm{u_0}_2+\int_0^T\norm f_2\dd t,\qquad
 L=TU^2+U^2/\nu,\qquad
 \Phi=\int_0^T\norm{\Pp f}_2^2\dd t.
\]
The kinetic balance gives \(\sup_t\norm u_2\le U\) and
\(\nu\int_0^S\norm{\nabla u}_2^2\dd t\le U^2/2\).
Using the Bessel Sobolev norm,
\[
 \norm u_{H^2}^2=\norm u_2^2+2\norm{\nabla u}_2^2+\norm{\Delta u}_2^2.
\]
Consequently the first rectangle
\(R_{0,1}(S)=\int_0^S(\norm u_{H^2}^2+\norm{u_t}_2^2)\dd t\)
satisfies
\begin{equation}
 \begin{aligned}
 R_{0,1}(S)&\le L+\max(1,\nu^{-2})X(S),\\
 X(S)&\le\max(1,\nu^2)R_{0,1}(S).
 \end{aligned}
 \label{r:eq:first-rectangle-equivalence}
\end{equation}
A bound for the full \(\mathcal G\) integral thus supplies both
velocity coordinates without first assuming that rectangle.

The gradient trace identity and Young's inequality also give
\[
 \sup_{t\le S}\norm{\nabla u(t)}_2^2\le G_0+X(S)/\nu,
 \qquad
 \int_0^S\norm u_{H^2}^2\dd t\le L+X(S)/\nu^2.
\]
Fourier inversion gives
\(\norm u_\infty^2\le\norm u_{H^2}^2/(8\pi)\), proved explicitly in
\eqref{r:eq:embedding}. Since
\(\norm{(u\cdot\nabla)u}_2\le\norm u_\infty\norm{\nabla u}_2\),
substitution into \eqref{r:eq:first-rectangle-exact} yields
\begin{equation}
 X(S)\le\nu G_0+2\Phi+
 \frac{(G_0+X(S)/\nu)(L+X(S)/\nu^2)}{4\pi}.
 \label{r:eq:first-rectangle-polynomial}
\end{equation}
This proves a quantified criterion directly from the data. Define
\[
 a=\frac1{4\pi\nu^3},\qquad
 b=\frac{G_0/\nu^2+L/\nu}{4\pi},\qquad
 c=\nu G_0+2\Phi+\frac{G_0L}{4\pi}.
\]
If \(b<1\) and \((1-b)^2>4ac\), then
\begin{equation}
 \sup_{S<T}X(S)\le
 \frac{1-b-\sqrt{(1-b)^2-4ac}}{2a}.
 \label{r:eq:first-rectangle-root}
\end{equation}
To see this, rewrite \eqref{r:eq:first-rectangle-polynomial} as
\(aX^2+(b-1)X+c\ge0\). Its two roots enclose a forbidden interval.
The continuous function \(X\), starting from zero, cannot cross that interval.
For rest datum the criterion becomes
\[
 \frac{U^2(T+\nu^{-1})}{4\pi\nu}
 +\sqrt{\frac{2\Phi}{\pi\nu^3}}<1.
\]
The bound is uniform through \(T\). Without the stated root separation,
\eqref{r:eq:first-rectangle-polynomial} does not supply an upper bound.

% Recovered from direct-forced-intersection.md, lines 466--516, and
% navier-stokes-pressure-material-investigation.md, lines 141--212,
% 553--763, and 821--868.  The main paper uses B_n for full transport.
% This fragment preserves that notation; projected transport is \Pp B_n.
% All norms in time are over (0,T).  Insert after the basic rectangle,
% product, pressure-projection, and acceleration identities are available.

\section{General bounds for production, pressure, and acceleration}\label{r:sec:recovered-rectangle}
\label{r:sec:recovered-rectangle-bounds}

The selected rectangles control more than the critical energy and the
pointwise momentum terms.  Each temporal row has its own nonlinear
production and source work.  Retaining the full binomial row lets us bound
these quantities together, and then use the two pressure projections to
resolve how transport contributes to material acceleration.

\subsection{The production and work of all temporal rows}

Fix nonnegative integers \(N,m\).  The pressure work in row \(r\) vanishes
after pairing with the solenoidal field \(V_r\).  The product estimate
\eqref{r:eq:transport-product} gives
\begin{equation}
 |\mathcal P_{r,m}|
 \le C_m\norm{V_r}_{H^m}
 \sum_{a=0}^r\binom ra
 \norm{V_a}_{H^{m+2}}\norm{V_{r-a}}_{H^{m+2}},
 \qquad C_m=2^{m+1}c_\infty.
 \label{r:eq:all-row-production-point}
\end{equation}
The two factors in every summand are retained: temporal differentiation
does not preserve every cancellation of the zeroth row.

Write
\[
 D_{N,m}:=\sum_{r=0}^N\norm{V_r(0)}_{H^m}^2.
\]
Integrating the adjacent trace identity from time zero and using
\(2ab\le a^2+b^2\) gives
\begin{equation}
 \max_{r\le N}\norm{V_r}_{L^\infty H^m}
 \le\sqrt{D_{N,m}+R_{N,m}}.
 \label{r:eq:all-row-initial-trace}
\end{equation}
For \(x_a=\norm{V_a}_{L^2H^{m+2}}\), time Cauchy--Schwarz and the
binomial symmetry give
\[
 \sum_{a=0}^r\binom ra x_ax_{r-a}
 \le 2^r\sum_{a=0}^N x_a^2
 \le 2^rR_{N,m+1}.
\]
We can now integrate \eqref{r:eq:all-row-production-point} and sum the full nonlinear contribution through row \(N\):
\begin{equation}
 \sum_{r=0}^N\norm{\mathcal P_{r,m}}_{L^1}
 \le C_m(2^{N+1}-1)
 \sqrt{D_{N,m}+R_{N,m}}\,R_{N,m+1}.
 \label{r:eq:all-row-production-integral}
\end{equation}
This bounds the absolute production in every selected row, including
each signed prefix integral.

The source work admits a bound with one adjacent spatial derivative
placed on each side of the pairing.  Since
\[
 W_{r,m}
 =\operatorname{Re}\ip{\Lambda^{m+1}V_r}{\Lambda^{m-1}F_r},
\]
define
\[
 \mathcal F_{N,m-1}^{\rm work}
 :=\left(\sum_{r=0}^N
 \norm{\Lambda^{m-1}F_r}_{L^2L^2}^2\right)^{1/2}.
\]
Cauchy--Schwarz first in time and then over the finite rows yields
\begin{equation}
 \sum_{r=0}^N\norm{W_{r,m}}_{L^1}
 \le\mathcal F_{N,m-1}^{\rm work}\sqrt{R_{N,m}}.
 \label{r:eq:all-row-force-work}
\end{equation}
At \(m=0\), the source coordinate is \(\dot H^{-1}\).  The assumed
spatial decay supplies its low frequencies.  For each fixed temporal
order \(\ell\), differentiation of the pairing gives the further bound
\begin{equation}
 \begin{aligned}
 \norm{\partial_t^\ell W_{r,m}}_{L^1}
 &\le\sum_{a=0}^{\ell}\binom\ell a
 \norm{\Lambda^{m+1}V_{r+a}}_{L^2L^2}
 \norm{\Lambda^{m-1}F_{r+\ell-a}}_{L^2L^2}.
 \end{aligned}
 \label{r:eq:all-row-force-work-derivatives}
\end{equation}
Every factor is a finite coordinate of the velocity or the prescribed
source.  The constants may depend on \(N,m,\ell\); the calculation
requires no bound uniform over all derivative orders.

\subsection{Projected transport from products and adjacent rows}

The pressure projection also gives direct integral estimates for the
nonlinear field, before it is paired with velocity.  Put
\[
 X_{a,j}:=\norm{V_a}_{L^2H^j}.
\]
For every nonnegative integer \(M\), choose the finite product constant
\(C_M^{\rm alg}\) in
\[
 \norm{v\otimes w}_{H^M}
 \le C_M^{\rm alg}\norm v_{H^{M+1}}\norm w_{H^{M+1}}.
\]
For \(M=0\), this follows from \(H^1\subset L^4\) and H\"older;
at higher integer orders it follows by the finite Leibniz expansion
and the Sobolev product bounds.  Tensor divergence maps \(H^M\) to
\(H^{M-1}\) with norm at most one, and both \(\Pp\) and
\(\mathbb Q=I-\Pp\) are contractions in these Fourier Sobolev norms.
Applying them to the full transport row gives
\begin{equation}
 \begin{aligned}
 \norm{\Pp B_n}_{L^1H^{M-1}},\quad
 \norm{\mathbb Q B_n}_{L^1H^{M-1}}
 &\le C_M^{\rm alg}
 \sum_{a=0}^n\binom na X_{a,M+1}X_{n-a,M+1}\\
 &\le C_M^{\rm alg}2^nR_{N,M},\qquad n\le N.
 \end{aligned}
 \label{r:eq:projected-transport-product-integral}
\end{equation}
Here and below the two quantities separated by a comma each obey the
displayed bound.

Let \(Q_n^u\) denote the normalized nonlinear pressure in row \(n\),
so that \(\nabla Q_n^u=-\mathbb Q B_n\).  Its scalar reconstruction
is a double Riesz transform of the binomial tensor sum.  The multiplier
has norm at most one from the Frobenius tensor norm to the scalar
norm.  The product estimate consequently gives
\begin{equation}
 \begin{aligned}
 \norm{Q_n^u}_{L^1H^M}
 +\norm{\nabla Q_n^u}_{L^1H^{M-1}}
 &\le 2C_M^{\rm alg}
 \sum_{a=0}^n\binom na X_{a,M+1}X_{n-a,M+1}\\
 &\le 2C_M^{\rm alg}2^nR_{N,M}.
 \end{aligned}
 \label{r:eq:nonlinear-pressure-product-integral}
\end{equation}
The total pressure gradient includes \(\mathbb QF_n\) as well.

Using the next temporal coordinate gives a stronger time norm for the
projected transport.  The momentum row itself reads
\[
 \Pp B_n-\Pp F_n=\nu\Delta V_n-V_{n+1}.
\]
Its two terms on the right belong to the adjacent rectangle, so
\[
 \norm{\Pp B_n-\Pp F_n}_{L^2H^m}
 \le\nu X_{n,m+2}+X_{n+1,m}.
\]
Scalar Cauchy--Schwarz with coefficients \((\nu,1)\), followed by
summation over \(n\), gives the explicit complete-row bound
\begin{equation}
 \left(\sum_{n=0}^N
 \norm{\Pp B_n-\Pp F_n}_{L^2H^m}^2\right)^{1/2}
 \le\sqrt{\nu^2+1}\,R_{N,m+1}^{1/2}.
 \label{r:eq:projected-transport-adjacent}
\end{equation}
The right side retains the next temporal row.  This is the direct
bound of the nonlinear field participating in the equation along the
selected history.

The spatial recovery in that equation is also explicit.  The Fourier
inequality \((1+|\xi|^2)^2\le2(1+|\xi|^4)\) implies
\begin{equation}
 \norm{V_n}_{H^{m+2}}
 \le\sqrt2\left[
 \norm{V_n}_{H^m}
 +\nu^{-1}\norm{V_{n+1}+\Pp B_n-\Pp F_n}_{H^m}\right].
 \label{r:eq:adjacent-spatial-recovery}
\end{equation}
Both the temporal derivative and the complete projected nonlinear
response remain in this recovery of two spatial derivatives.

\subsection{Higher pressure rows and their time norms}

The pressure gain in \eqref{r:eq:pressure-critical-gain} extends to every
temporal row.  For two solenoidal fields, expanding the divergence gives
\[
 \operatorname{div}((v\cdot\nabla)w)
 =\partial_i v_j\,\partial_j w_i.
\]
This expression is symmetric under exchanging \(v,w\).  The gradient
projection has the exact norm relation
\[
 \norm{\mathbb Q((v\cdot\nabla)w)}_{\dot H^{1/2}}
 =\norm{\partial_i v_j\,\partial_j w_i}_{\dot H^{-1/2}}.
\]
With \(c_*=(\sqrt2\pi)^{-1}\), the Sobolev inequality and its dual
bound the right side by
\(c_*\norm v_{\dot H^{3/2}}\norm w_{\dot H^{3/2}}\).
Apply this to every binomial summand:
\begin{equation}
 \norm{\nabla Q_n^u}_{\dot H^{1/2}}
 \le c_*\sum_{a=0}^n\binom na
 \norm{V_a}_{\dot H^{3/2}}\norm{V_{n-a}}_{\dot H^{3/2}}.
 \label{r:eq:pressure-all-temporal-endpoint}
\end{equation}
The symmetry also gives
\begin{equation}
 \begin{aligned}
 \norm{\mathbb Q((v\cdot\nabla)w)}_{\dot H^{-1/2}}
 \le c_*\min\bigl\{&
 \norm v_{\dot H^{1/2}}\norm w_{\dot H^{3/2}},\\
 &\norm w_{\dot H^{1/2}}\norm v_{\dot H^{3/2}}\bigr\}.
 \end{aligned}
 \label{r:eq:pressure-factor-exchange}
\end{equation}
Either assignment of the two heights is available for the pressure
component because its divergence retains a derivative on both factors.

For the zeroth row, put
\[
 h(t)=\norm{u(t)}_{\dot H^{1/2}},\qquad
 d(t)=\norm{u(t)}_{\dot H^{3/2}},\qquad
 H_*:=\norm h_{L^\infty},\quad D_*:=\norm d_{L^2}.
\]
The bounds \(\norm{B_0}_{\dot H^{-1/2}}\le c_*hd\) and
\(\norm{\mathbb QB_0}_{\dot H^{1/2}}\le c_*d^2\) give
\begin{equation}
 \begin{aligned}
 \norm{\Pp B_0}_{L^2\dot H^{-1/2}},\quad
 \norm{\mathbb QB_0}_{L^2\dot H^{-1/2}}
 &\le c_*H_*D_*,\\
 \norm{\mathbb QB_0}_{L^1\dot H^{1/2}}
 &\le c_*D_*^2.
 \end{aligned}
 \label{r:eq:pressure-time-endpoints}
\end{equation}
Interpolating the two spatial norms at each time gives
\[
 \norm{\mathbb QB_0(t)}_2
 \le c_*h(t)^{1/2}d(t)^{3/2}.
\]
Raising this to the power \(4/3\) and integrating yields
\begin{equation}
 \norm{\mathbb QB_0}_{L^{4/3}L^2}
 \le c_*H_*^{1/2}D_*^{3/2}.
 \label{r:eq:pressure-intermediate-time}
\end{equation}
The critical rectangle supplies \(H_*\le\sqrt{D+R}\) and
\(D_*\le\sqrt R\), so the last right side is
\(c_*(D+R)^{1/4}R^{3/4}\).  A Sobolev and H\"older estimate also
gives the \(L^{4/3}L^2\) conclusion for full transport.  The additional
pressure regularity is its \(L^1\dot H^{1/2}\) endpoint; total forced
pressure still includes the prescribed gradient source.

\subsection{Material acceleration from the first rectangle}

The orthogonal acceleration identity can be evaluated in \(L^2_x\)
using \(R_{0,1}\).  Set
\[
 \mathsf A:=\norm u_{L^2H^2},\qquad
 \mathsf B:=\norm{u_t}_{L^2L^2},\qquad
 \mathsf Z:=\left(\norm{u_0}_{H^1}^2
                   +2\mathsf A\mathsf B\right)^{1/2}.
\]
The adjacent trace gives \(\norm u_{L^\infty H^1}\le\mathsf Z\).
Use the homogeneous Sobolev constants
\[
 S_6=\frac{2^{2/3}}{\sqrt3\pi^{2/3}},\qquad
 S_3=2^{-1/6}\pi^{-1/3}.
\]
H\"older in space and then in time gives
\begin{equation}
 \begin{aligned}
 \norm{B_0}_{L^2L^2}
 &\le\norm u_{L^\infty L^6}\norm{\nabla u}_{L^2L^3}\\
 &\le S_6S_3\mathsf Z\mathsf A
 =\frac1\pi\sqrt{\frac23}\,\mathsf Z\mathsf A.
 \end{aligned}
 \label{r:eq:first-rectangle-transport-L2}
\end{equation}
In unforced flow, viscosity supplies the solenoidal part of acceleration
and the pressure response supplies its gradient part.  Their squared
norms add by \eqref{r:eq:acceleration-orthogonal}.  We obtain
\begin{equation}
 \norm{D_tu}_{L^2L^2}
 \le\mathsf A\sqrt{\nu^2+\frac{2\mathsf Z^2}{3\pi^2}},
 \qquad f=0.
 \label{r:eq:first-rectangle-acceleration-L2}
\end{equation}
For the forced equation, retain the prescribed solenoidal source in
\eqref{r:eq:acceleration-orthogonal}. The resulting bound at this first
rectangle is
\begin{equation}
 \norm{D_tu}_{L^2L^2}^2
 \le\left(\nu\mathsf A+\norm{\Pp f}_{L^2L^2}\right)^2
 +\frac{2}{3\pi^2}\mathsf Z^2\mathsf A^2.
 \label{r:eq:first-rectangle-forced-acceleration-L2}
\end{equation}
The triangle inequality bounds the solenoidal component, while the
preceding transport estimate bounds the orthogonal gradient component.
The adjacent row also gives
\(\norm{\Pp B_0}_{L^2L^2}\le\nu\mathsf A+\mathsf B\) when
\(f=0\).  These bounds use the first rectangle, while the earlier
pointwise acceleration bound selects the stronger \(R_{1,2}\).

At critical negative-half order, the orthogonal identity gives a
particularly precise unforced coefficient:
\begin{equation}
 \nu d(t)\le\norm{D_tu(t)}_{\dot H^{-1/2}}
 \le\sqrt{\nu^2+\frac{h(t)^2}{2\pi^2}}\,d(t),
 \qquad f=0.
 \label{r:eq:critical-acceleration-two-sided}
\end{equation}
The lower bound retains the solenoidal viscous component, and the upper
bound uses \(\norm{\mathbb QB_0}_{\dot H^{-1/2}}\le c_*hd\).
Integration gives
\[
 \norm{D_tu}_{L^2\dot H^{-1/2}}
 \le\sqrt{\nu^2+\frac{H_*^2}{2\pi^2}}\,D_*,\qquad f=0.
\]
For general forcing, the solenoidal component is
\(\nu\Delta u+\Pp f\); retaining it gives
\begin{equation}
 \norm{D_tu}_{L^2\dot H^{-1/2}}^2
 \le\left(\nu D_*+\norm{\Pp f}_{L^2\dot H^{-1/2}}\right)^2
       +c_*^2H_*^2D_*^2.
 \label{r:eq:critical-acceleration-general-coordinates}
\end{equation}
Its substitution from \(R_{0,1}\) is
\eqref{r:eq:critical-material-bounds}.

\subsection{Acceleration at a higher adjacent height}

To place the full acceleration in \(L^2H^{1/2}\), select the stronger
adjacent coordinates
\begin{equation}
 u\in L^2H^{5/2},\qquad u_t\in L^2H^{1/2}.
 \label{r:eq:stronger-acceleration-coordinates}
\end{equation}
Higher integer rectangles supply both by Sobolev inclusion.  The
inhomogeneous pairing
\(2\ip{(1-\Delta)^{5/4}u}{(1-\Delta)^{1/4}u_t}\)
gives the continuous \(H^{3/2}\) trace and
\begin{equation}
 \sup_{0\le t\le T}\norm{u(t)}_{H^{3/2}}^2
 \le\norm{u_0}_{H^{3/2}}^2
 +2\norm u_{L^2H^{5/2}}\norm{u_t}_{L^2H^{1/2}}.
 \label{r:eq:stronger-acceleration-trace}
\end{equation}
For the pressure component, the negative-half estimate controls
frequencies below one and the positive-half estimate controls those
above one.  Together they give a continuous quadratic map from
\(H^{3/2}\) into \(H^{1/2}\).  Thus
\(\mathbb QB_0\in C([0,T];H^{1/2})\), with
\begin{equation}
 \norm{\mathbb QB_0}_{L^2H^{1/2}}
 \le C\norm u_{L^\infty H^{3/2}}
         \norm u_{L^2H^{3/2}}.
 \label{r:eq:stronger-pressure-time}
\end{equation}
If \(\Pp f\in L^2H^{1/2}\), all components of
\(a=\nu\Delta u+\Pp f+\mathbb QB_0\) now belong to
\(L^2H^{1/2}\).  The projected momentum row separately yields
\begin{equation}
 \norm{\Pp B_0}_{L^2H^{1/2}}
 \le\nu\norm u_{L^2H^{5/2}}
 +\norm{u_t}_{L^2H^{1/2}}
 +\norm{\Pp f}_{L^2H^{1/2}}.
 \label{r:eq:stronger-projected-transport}
\end{equation}
The \(L^2H^{5/2}\) bound controls the viscous term in
\(L^2H^{1/2}\), while the adjacent temporal bound yields
the continuous \(H^{3/2}\) trace.

\subsection{Unforced action and spatial recovery}

For \(f=0\), squaring the projected equation in
\(\dot H^{-1/2}\) retains an exact relation between the temporal
and viscous actions.  Since
\[
 \Pp B_0=\nu\Delta u-u_t,\qquad
 \ip{\Delta u}{u_t}_{\dot H^{-1/2}}
 =-\frac12\frac{\dd}{\dd t}h^2,
\]
we have
\begin{equation}
 \nu\frac{\dd}{\dd t}h^2+\nu^2d^2
 +\norm{u_t}_{\dot H^{-1/2}}^2
 =\norm{\Pp B_0}_{\dot H^{-1/2}}^2,
 \qquad f=0.
 \label{r:eq:unforced-critical-action}
\end{equation}
Add \(\norm{\mathbb QB_0}_{\dot H^{-1/2}}^2\) to both sides.
Orthogonality of the two transport projections and the acceleration
identity give
\begin{equation}
 \nu\frac{\dd}{\dd t}h^2
 +\norm{u_t}_{\dot H^{-1/2}}^2
 +\norm{D_tu}_{\dot H^{-1/2}}^2
 =\norm{B_0}_{\dot H^{-1/2}}^2,
 \qquad f=0.
 \label{r:eq:unforced-material-critical-action}
\end{equation}
These identities retain the complete nonlinear action on the right.
Adding the pressure square also adds its contribution there; the
operation supplies no additional sign for critical production.

The solenoidal acceleration determines the viscous spatial derivative
when \(\nu>0\).  In unforced flow,
\(\Pp a=\nu\Delta u\), so the Fourier estimate used in
\eqref{r:eq:adjacent-spatial-recovery} gives
\begin{equation}
 \norm u_{H^{m+2}}
 \le\sqrt2\left(\norm u_{H^m}
                  +\nu^{-1}\norm a_{H^m}\right),
 \qquad f=0.
 \label{r:eq:unforced-spatial-recovery-from-acceleration}
\end{equation}
This last relation recovers the spatial coordinate from the full
acceleration at that height, with the lower velocity norm retained.

\section{Viscous dissipation and material deformation}\label{r:joined-sec:metric}

The deformation of a material parcel is controlled in an integral
sense by viscous dissipation. Let $X(a,t)$ follow the parcel initially labeled $a$, let $F=\nabla_aX$, and let $C=F^TF$. The positive matrix $C$ records the lengths and angles of infinitesimal material vectors. Incompressibility preserves volume, while the symmetric velocity gradient $S=(\nabla u+\nabla u^T)/2$ changes this metric:
\begin{equation}\label{r:joined-eq:metric}
\dot F=(\nabla u)(X,t)F,\qquad
\det F=1,\qquad
\dot C=2F^TS(X,t)F.
\end{equation}
The full metric speed has the exact squared norm
\begin{equation}\label{r:joined-eq:metricspeed}
\tr(C^{-1}\dot C\,C^{-1}\dot C)=4|S(X,t)|^2.
\end{equation}
Indeed, $C^{-1}\dot C=2F^{-1}SF$, and cyclicity of the trace removes the two deformation matrices. After changing variables by the incompressible flow and integrating by parts in space,
\begin{equation}\label{r:joined-eq:metricaction}
\frac\nu2\int_0^t\!\int_{\R^3}
\tr(C^{-1}\dot C\,C^{-1}\dot C)\,da\,ds
=\nu\int_0^t\norm{\nabla u(s)}_2^2ds.
\end{equation}
Here $\norm S_2^2=\frac12\norm{\nabla u}_2^2$. The dissipation integral consequently controls the rate of change
of the material metric, including shape changes that preserve volume.

The energy budget makes this quantitative through the full time interval. Define
\begin{equation*}
F_T=\int_0^T\norm f_2dt,\quad
M_T=\norm{u_0}_2+F_T,\quad
K_T=\norm{u_0}_2^2+2M_TF_T.
\end{equation*}
The energy identity gives $\sup_t\norm u_2\le M_T$ and
$2\nu\int_0^T\norm{\nabla u}_2^2\le K_T$.
If $\sigma_i$ are the singular values of $F$, then the derivatives of $\log\sigma_i$ are diagonal entries of $S$ in a suitable orthonormal frame. Their squared sum is bounded by $|S|^2$. Integrating from $F(a,0)=I$ gives
\begin{equation*}
\norm{\log C(a,t)}_F\le2\int_0^t|S(X(a,s),s)|\,ds.
\end{equation*}
Repeated singular values are handled by diagonalizing the symmetric variation within each eigenspace. Cauchy--Schwarz and volume preservation now yield the maximal estimate
\begin{equation}\label{r:joined-eq:maxdistortion}
\int_{\R^3}\sup_{t<T}\norm{\log C(a,t)}_F^2\,da
\le\frac{TK_T}{\nu}.
\end{equation}
This is a bound over original material labels with the time supremum inside the integral. It rules out infinite distortion on a set of labels of positive measure. It permits concentration onto label sets whose measure tends to zero. A uniform bound over all labels would require further control beyond \eqref{r:joined-eq:maxdistortion}.


\section{Angular momentum and the role of symmetry}
\label{r:sec:angular}

The distinction between the leading vortex and the completed velocity
matters to the comparison. The background in OpenAI's construction is
axisymmetric. Its wave potentials contain factors
\(e^{ikm\Phi}\), with \(\Phi=p\theta+\varphi\) and nonzero integer
angular frequency \(kmp\)
\cite[(6.27), (7.3), (7.38)]{openai-paper}. Taking their curls retains
these harmonics. They enter the realized and localized velocity through
\cite[(9.21), (10.4)]{openai-paper}; the final spatial cutoff does not
average them over the angle. The complete field is nonaxisymmetric.

One can see exactly why that distinction affects a force bound. Define
angular momentum per unit mass about the axis by
\(\Gamma=ru_\theta\). The angular component of \eqref{r:eq:ns}, including
the derivatives of the cylindrical basis vectors, gives for \(r>0\)
\begin{equation}
 \begin{aligned}
 \left(\partial_t+u_r\partial_r
       +\frac{u_\theta}{r}\partial_\theta+u_z\partial_z\right)\Gamma
 &=\nu\left(\partial_{rr}-\frac1r\partial_r
       +\frac1{r^2}\partial_{\theta\theta}+\partial_{zz}\right)\Gamma\\
 &\quad-\partial_\theta p
       +\frac{2\nu}{r}\partial_\theta u_r+rf_\theta.
 \end{aligned}
 \label{r:eq:angular-full}
\end{equation}
Indeed, multiplying the angular momentum equation for \(u_\theta\) by
\(r\) combines its geometric term \(u_ru_\theta\) with the material
derivative of \(r\). The vector Laplacian contributes
\(2\nu r^{-1}\partial_\theta u_r\). Pressure supplies the torque
\(-\partial_\theta p\). These terms describe internal momentum transfer;
they are present even when the prescribed force is small.

\begin{proposition}[A force bound for axisymmetric angular momentum]
\label{r:prop:angular-max}
Suppose \(u,p,f\) are axisymmetric, solve \eqref{r:eq:ns} from rest, and
\(u\) is bounded on every \(\R^3\times[0,S]\), \(S<T\).
If \(G(t):=\int_0^t\norm{rf_\theta(s)}_\infty\dd s\) is finite through
\(T\), then
\begin{equation}
 \norm{ru_\theta(t)}_\infty\le G(t),\qquad 0\le t<T.
 \label{r:eq:angular-max}
\end{equation}
There is no source-amplitude restriction in this bound.
\end{proposition}

\begin{proof}
Under these symmetry assumptions the angular derivatives in
\eqref{r:eq:angular-full} vanish. The resulting scalar operator is
\[
 L=\partial_t+u_r\partial_r+u_z\partial_z
       -\nu(\partial_{rr}-r^{-1}\partial_r+\partial_{zz}).
\]
Fix \(S<T\) and let \(M_S=\sup_{[0,S]\times\R^3}|u|\).
For \(\Psi=1+r^2+z^2\) and \(c=M_S+2\nu+1\), direct differentiation
gives
\[
 L(e^{ct}\Psi)
 =e^{ct}\bigl(c\Psi+2ru_r+2zu_z-2\nu\bigr)>0.
\]
Here \(2M_S\sqrt{r^2+z^2}\le M_S\Psi\).
Compare \(\pm\Gamma\) with \(G(t)+\epsilon e^{ct}\Psi\).
The comparison functions dominate initially. At the axis,
\(\Gamma=0\) by Cartesian smoothness; at spatial infinity their
quadratic growth dominates \(|\Gamma|\le M_Sr\).
Apply the scalar maximum principle on bounded cylinders away from the
axis, using the integral of the maximum of \(|rf_\theta|\) on each
cylinder as the temporal barrier. That integral is at most \(G\).
For sufficiently large outer radius the quadratic barrier dominates the
outer boundary values; at a sufficiently small excluded axial radius,
\(|\Gamma|\le M_Sr\) is dominated by the positive barrier.
Exhausting the domain gives
\(\pm\Gamma\le G+\epsilon e^{ct}\Psi\) on \([0,S]\).
Let \(\epsilon\downarrow0\) and then use the arbitrariness of \(S<T\).
\end{proof}

Smooth rapidly decreasing forcing satisfies the hypothesis on \(G\).
This estimate would exclude an axisymmetric realization of the exterior
growth calculated in Section~\ref{r:sec:exact-cancellation}, since there \(ru_\theta\) grows like
\((1-t)^{-h}\). It does not assert global regularity for every
axisymmetric velocity: a bound on \(ru_\theta\) is not a uniform bound
on \(u_\theta\) near the axis.

Nor does angular averaging turn the completed nonaxisymmetric field into
a solution of the axisymmetric equation. Write a bar for angular average
and a prime for the fluctuation about that average. Averaging the
conservative form of \eqref{r:eq:angular-full} gives exactly
\begin{equation}
 \begin{aligned}
 \partial_t\bar\Gamma+\bar u_r\partial_r\bar\Gamma
       +\bar u_z\partial_z\bar\Gamma
 &=\nu(\partial_{rr}-r^{-1}\partial_r+\partial_{zz})\bar\Gamma
       +r\bar f_\theta\\
 &\quad-\frac1r\partial_r
       \left(r\overline{u_r'\Gamma'}\right)
       -\partial_z\overline{u_z'\Gamma'}.
 \end{aligned}
 \label{r:eq:angular-mean}
\end{equation}
The pressure torque and angular viscous derivative average to zero, but
the products in the transport flux do not. Two zero-mean harmonics can
have a nonzero averaged product. OpenAI's covariance construction uses
precisely these products to supply momentum flux
\cite[(7.23), (9.11)]{openai-paper}. A force-only estimate for
\(\bar\Gamma\) would have to control the two fluctuation fluxes in
\eqref{r:eq:angular-mean}. Dropping them would change the equation being
compared.



\section{Endpoint jets and material derivatives}\label{r:app:endpoint}

The intersection also gives \(V_n\in H^1(0,T;H^m)\) at every fixed
\(n,m\). Its continuous representative has a limit at \(T\).
Product continuity in the Sobolev spaces used above gives limits for \(B_n\);
\eqref{r:eq:pressure} gives limits for \(\nabla Q_n\). Viscosity and
spatial differentiation are continuous between the stated spaces. Each
momentum term thus extends to the endpoint in every finite spatial norm,
and \eqref{r:eq:embedding} gives uniform convergence of each fixed spatial
derivative.

\begin{samepage}
These limits have quantitative rates. Since
\(\partial_t B_n=B_{n+1}\), integrating the next row yields, for
\(|\alpha|=k\),
\begin{align}
 \norm{\partial_x^\alpha[V_{n+1}(t)-V_{n+1}(T)]}_\infty
 &\le c_\infty(T-t)A_{n+2,k+2},\label{r:eq:rate-time}\\
 \norm{\partial_x^\alpha[B_n(t)-B_n(T)]}_\infty
 &\le c_\infty(T-t)S_{n+1,k+2},\label{r:eq:rate-transport}\\
 \norm{\partial_x^\alpha[\nabla Q_n(t)-\nabla Q_n(T)]}_\infty
 &\le c_\infty(T-t)(L_{n+1,k+2}+S_{n+1,k+2}),\label{r:eq:rate-pressure}\\
 \norm{\nu\partial_x^\alpha\Delta[V_n(t)-V_n(T)]}_\infty
 &\le\nu c_\infty(T-t)A_{n+1,k+4}.\label{r:eq:rate-viscosity}
\end{align}
\end{samepage}
To justify the integration, first apply the fundamental theorem on
\([t,S]\), then let \(S\uparrow T\) in the continuous Sobolev
representatives. The pressure rate follows by applying \(I-\Pp\) to
\(F_{n+1}-B_{n+1}\). The rectangle \(R_{n+2,k+4}\), together with the
displayed source norm, contains all coordinates used by these four rates.

The terminal fields also obey the differentiated equation:
\[
 V_{n+1}(T)+B_n(T)+\nabla Q_n(T)-\nu\Delta V_n(T)=F_n(T).
\]
All its spatial derivatives agree with the limits of the corresponding
spatially differentiated rows. The endpoint is a compatible velocity and
pressure-gradient jet, rather than a collection of unrelated terminal fields.



\subsection{Material acceleration and its higher derivatives}

The fluid-particle acceleration uses both the first temporal row and
transport:
\begin{equation}
 D_tu=V_1+B_0=-\nabla p+\nu\Delta u+f.
\label{r:eq:material}
\end{equation}
Its uniform bound follows from \eqref{r:eq:material-one-rectangle}.
The next parcel derivative differentiates both contributions:
\begin{equation}
 \begin{aligned}
 D_t^2u={}&V_2+2(V_0\cdot\nabla)V_1+(V_1\cdot\nabla)V_0\\
 &+(V_0\cdot\nabla)\bigl[(V_0\cdot\nabla)V_0\bigr].
 \end{aligned}
 \label{r:eq:jerk-mixed-jet}
\end{equation}
The factor two collects the temporal derivative of transport and the
transport of \(V_1\). Jerk is already a composite of the mixed jet.

Write \(D=D_t\), \(a=Du\), and \(U_{ij}=\partial_j u_i\).
Ordinary derivatives commute, so
\(\nabla\cdot V_n=\partial_t^n(\nabla\cdot u)=0\).
Moving with the fluid introduces the velocity coefficient:
\[
 [D,\partial_i]w=-(\partial_i u_j)\partial_jw.
\]
For a smooth scalar or vector field \(w\), expanding the commutators gives
\begin{equation}
 \begin{aligned}
 [D,\nabla]p&=-U^{\mathsf T}\nabla p,\\
 [D,\Delta]w&=-(\Delta u)\cdot\nabla w
 -2\sum_i(\partial_i u)\cdot\nabla\partial_i w.
 \end{aligned}
 \label{r:eq:material-commutators}
\end{equation}
The jerk of the parcel is consequently
\begin{equation}
 \begin{aligned}
 Da={}&\nu\Delta a-\nabla(Dp)+U^{\mathsf T}\nabla p+Df\\
 &-\nu\left[(\Delta u)\cdot\nabla u
 +2\sum_i(\partial_i u)\cdot\nabla\partial_i u\right].
 \end{aligned}
 \label{r:eq:material-jerk}
\end{equation}
Every commutator coefficient is a spatial derivative of this velocity.
The pressure and viscous terms remain attached to the mixed jet.
In particular,
\begin{equation}
 \begin{aligned}
 D(\nabla\cdot u)&=\nabla\cdot a-\operatorname{tr}(U^2)=0,\\
 \nabla\cdot a&=\operatorname{tr}(U^2)
              =-\Delta p+\nabla\cdot f .
 \end{aligned}
 \label{r:eq:acceleration-divergence}
\end{equation}
Acceleration's divergence is determined by the velocity gradient and
the pressure equation. Its sign need not be positive:
\(\operatorname{tr}(U^2)\) differs from \(|U|_{\mathrm F}^2\).
The pressure pairing in a material row cannot be deleted using the
divergence-free property of \(V_1\).

Here spatial divergence and unbounded magnitude are distinct questions:
\(\nabla\cdot a\ne0\) does not assert that \(|a|\) is infinite.
Clay's incompressibility condition remains \(\nabla\cdot u=0\).
Equation~\eqref{r:eq:acceleration-divergence} follows by differentiating
that condition and is already encoded in its pressure Poisson equation.
Retaining this identity introduces no additional hypothesis on a smooth
solution of the original system.

A material parcel still preserves its volume. Let \(X(\ell,t)\) be the
flow from the particle label \(\ell\), and let
\(j(\ell,t)=\det\partial_\ell X(\ell,t)\), with \(j(\ell,0)=1\).
Differentiating the deformation matrix gives
\(\partial_t(\partial_\ell X)=U(X,t)\partial_\ell X\), so
\[
 \partial_tj=(\nabla\cdot u)(X,t)j=0,\qquad j(\ell,t)=1.
\]
For a bounded initial parcel \(E\), its volume is
\(|X(E,t)|=\int_E j(\ell,t)\dd\ell=|E|\).
Every time derivative of that volume vanishes. At second order the
Jacobian formula displays the cancellation explicitly:
\begin{equation}
 \partial_t^2j
 =\bigl[\nabla\cdot a-\operatorname{tr}(U^2)
                 +(\nabla\cdot u)^2\bigr](X,t)j=0.
 \label{r:eq:parcel-volume}
\end{equation}
The nonzero divergence of acceleration is precisely accompanied by
the changing parcel geometry. It does not create an additional degree
of freedom or violate incompressibility.
With the rectangle and constant \(A\) used in
\eqref{r:eq:material-one-rectangle},
\[
 \norm{\nabla\cdot a}_{L^\infty_{t,x}}
 \le\norm{\nabla u}_{L^\infty_{t,x}}^2
 \le\frac{A^2}{8\pi}.
\]

The constraint continues through every material order. Put
\(w_q=D^qu\) and \(\delta_q=\nabla\cdot w_q\). Expanding the
divergence of \(Dw_q\) gives the exact recurrence
\begin{equation}
 \delta_0=0,\qquad
 \delta_{q+1}=D\delta_q+\operatorname{tr}(U\nabla w_q).
 \label{r:eq:material-incompressibility}
\end{equation}
For example, \(DU=\nabla a-U^2\), so the jerk satisfies
\[
 \nabla\cdot(Da)
 =3\operatorname{tr}(U\nabla a)-2\operatorname{tr}(U^3).
\]
One further differentiation gives, with \(G=U\),
\(B_1=\nabla a\), and \(B_2=\nabla(Da)\),
\[
 \nabla\cdot(D^3u)
 =4\operatorname{tr}(GB_2)+3\operatorname{tr}(B_1^2)
 -12\operatorname{tr}(G^2B_1)+6\operatorname{tr}(G^4).
\]
This follows from \(DB_k=B_{k+1}-B_kG\) and cyclicity of the trace.

The higher spatial divergences are fixed by the lower material and
spatial derivatives. Along the flow map, these relations express the
compatibility of acceleration, jerk, and every subsequent derivative
with \(j\equiv1\).

For the full recurrence, put
\(\mathcal L_0=\Delta\), \(\mathcal G_0=\nabla\), and define
\(\mathcal L_{q+1}=[D,\mathcal L_q]\),
\(\mathcal G_{q+1}=[D,\mathcal G_q]\). Then
\begin{equation}
 \begin{aligned}
 D^{n+1}u={}&\nu\sum_{q=0}^n\binom nq\mathcal L_qD^{n-q}u\\
 &-\sum_{q=0}^n\binom nq\mathcal G_qD^{n-q}p+D^nf.
 \end{aligned}
 \label{r:eq:material-binomial}
\end{equation}
Induction uses \(D(Lw)=L(Dw)+[D,L]w\) and Pascal's identity.
Every operator coefficient is generated by the velocity in \eqref{r:eq:ns}.
The first two material source derivatives are
\[
 \begin{aligned}
 Df&=f_t+u\cdot\nabla f,\\
 D^2f&=f_{tt}+2u\cdot\nabla f_t
 +a\cdot\nabla f+\sum_{i,j}u_i u_j\partial_{ij}f.
 \end{aligned}
\]
These derivatives measure the force encountered by the moving parcel.
Their velocity factors are part of the coupled response, while the
Eulerian derivatives of the prescribed field remain source data.
Further material derivatives differentiate the velocity coefficients inside
\(D_t\). Put \(J_{r,\alpha}=\partial_t^r\partial_x^\alpha u\).
The exact reconstruction rule is
\begin{equation}
 D_tJ_{r,\alpha}
 =J_{r+1,\alpha}+\sum_{\ell=1}^3
      u_\ell J_{r,\alpha+e_\ell}.
 \label{r:eq:material-jet-reconstruction}
\end{equation}
Start with \(D_t^0u=J_{0,0}\), apply this rule to each factor, and use
Leibniz's rule for every product. After \(q\) steps, every component of
\(D_t^qu\) is a polynomial in mixed jet coordinates. Each monomial
contains at most \(q+1\) velocity factors, and the sum of their
derivative orders \(r+|\alpha|\) is exactly \(q\). This proves
the reconstruction at every finite order without introducing
independent material rows.

The reverse construction is equally explicit. Define the material--spatial
coordinates \(M_{q,\alpha}=\partial_x^\alpha D_t^qu\). Ordinary
spatial differentiation increments \(\alpha\). Expanding
\(\partial_x^\alpha(D_tw)\) for \(w=D_t^qu\) gives
\begin{equation}
 \partial_tM_{q,\alpha}
 =M_{q+1,\alpha}
 -\sum_{\beta\le\alpha}\binom\alpha\beta
       (M_{0,\beta}\cdot\nabla)M_{q,\alpha-\beta}.
 \label{r:eq:eulerian-jet-reconstruction}
\end{equation}
Leibniz's rule now generates every ordinary mixed derivative from the
material--spatial coordinates. These conversion rules relate the Eulerian and material derivative
algebras. Their different operator orders are accounted for by
the displayed commutator terms.

\begin{proposition}[Equality of the complete derivative intersections]
\label{r:prop:material-intersection}
Let \(0<T<\infty\), and let \(u\) be a smooth vector field on
\([0,T)\times\R^3\). Define \(D_t=\partial_t+u\cdot\nabla\) from this field.
Then
\begin{equation}
 \begin{aligned}
 &\partial_t^ru\in L^2(0,T;H^m)
       &&\text{for every }r,m\ge0\\
 &\hspace{20mm}\Longleftrightarrow\\[-2mm]
 &D_t^qu\in L^2(0,T;H^m)
       &&\text{for every }q,m\ge0.
 \end{aligned}
 \label{r:eq:intersection-equivalence}
\end{equation}
Here \(r,q,m\) are integers. No commutation of \(D_t\) with spatial
derivatives is assumed.
\end{proposition}
\begin{proof}
The Eulerian intersection gives
\(u\in H^1(0,T;H^m)\) at every spatial height, and gives this
property for each temporal derivative as well. Hilbert-valued time
embedding supplies their uniform \(H^m\) bounds. The polynomial
reconstruction \eqref{r:eq:material-jet-reconstruction}, spatial Sobolev
multiplication, and \(T<\infty\) give every material coordinate in
\(L^2_tH^m_x\).
For the reverse implication put \(w_q=D_t^qu\) and
\(b_{q,m}=\norm{w_q}_{L^2_tH^m_x}\). The exact identity
\[
 \partial_tw_q=w_{q+1}-(u\cdot\nabla)w_q
\]
first gives a time-integrable derivative. For integer \(m\ge0\),
Leibniz's rule and \(H^2\hookrightarrow L^\infty\) give a finite
constant \(C_m\) such that
\[
 \norm{v\cdot\nabla w}_{H^m}
 \le C_m\norm{v}_{H^{m+2}}\norm{w}_{H^{m+2}}.
\]
Indeed, in each derivative of the product, the differentiated \(v\)
factor is bounded in \(L^\infty\) using two additional derivatives,
and the differentiated \(\nabla w\) factor is bounded in \(L^2\).
Cauchy--Schwarz in time now yields
\[
 \norm{\partial_tw_q}_{L^1_tH^m_x}
 \le\sqrt T\,b_{q+1,m}+C_m b_{0,m+2}b_{q,m+2}.
\]
Thus \(w_q\in W^{1,1}(0,T;H^m)\). Select a time where its
\(H^m\) norm is at most \(T^{-1/2}b_{q,m}\), or take a limiting
sequence of such times, and integrate its derivative. This gives the
finite-coordinate bound
\begin{equation}
 \begin{aligned}
 \sup_{t<T}\norm{w_q(t)}_{H^m}
 &\le L_{q,m},\\
 L_{q,m}
 &:=T^{-1/2}b_{q,m}+\sqrt T\,b_{q+1,m}\\
 &\hspace{7mm}+C_m b_{0,m+2}b_{q,m+2}.
 \end{aligned}
 \label{r:eq:material-reverse-trace}
\end{equation}
Every material row now has uniform spatial bounds at every finite
height. The product reconstruction
\eqref{r:eq:eulerian-jet-reconstruction} expresses each Eulerian mixed
derivative using finitely many of those rows and spatial derivatives.
Sobolev multiplication bounds each product in \(L^\infty_tH^m_x\);
integration on the finite interval gives its \(L^2_tH^m_x\) norm.
This proves the reverse implication.
\end{proof}

The proposition identifies the complete intersection independently
of construction order. Its reverse trace estimate also shows how a finite
collection of material coordinates supplies an Eulerian bound.
For an explicit bound in the other direction, set
\[
 K_q=\max_{r+|\alpha|\le q}
 \norm{\partial_t^r\partial_x^\alpha u}_{L^\infty_{t,x}}<\infty.
\]
At the next material differentiation, a product of at most \(q+1\)
factors produces at most \(4(q+1)\) terms: one time derivative and
three convective derivatives per factor. Starting with one term at \(q=0\)
gives the sufficient componentwise coefficient bound \(4^q q!\). Hence
\begin{equation}
 \norm{D_t^q u}_{L^\infty_{t,x}}
 \le\sqrt3\,4^q q!\max(1,K_q)^{q+1}<\infty.
 \label{r:eq:material-all}
\end{equation}
The bound is finite for every fixed order \(q\). More generally,
applying \(D_t^q\) to a vector field \(w\)
produces products with one mixed derivative of \(w\) and at most \(q\)
velocity factors. The identical counting argument bounds it by
\(\sqrt3\,4^q q!\max(1,K_q)^qK_q^w\), where \(K_q^w\) is the
maximum of its mixed derivative norms through order \(q\).
For the prescribed source, set
\(K_q^f=\max_{r+|\alpha|\le q}
\norm{\partial_t^r\partial_x^\alpha f}_{L^\infty_{t,x}}\). This gives
\[
 \norm{D_t^q f}_{L^\infty_{t,x}}
 \le\sqrt3\,4^q q!\max(1,K_q)^qK_q^f.
\]
These formulas preserve the dependence of a material force derivative on
the fluid moving through its prescribed spatial field.

The bounded velocity and its bounded spatial gradient give a unique fluid
trajectory through each initial point on the closed interval. Along that
trajectory, ordinary time derivatives of \(u(X(t),t)\) are the material
derivatives just bounded. This connects the Eulerian rows to acceleration
and its higher material derivatives without identifying
\(\partial_t^n u\) with \(D_t^n u\).



\subsection{Explicit rectangles for finite material orders}
\label{r:sec:material-target-rectangle}

The finite footprint of the material jet gives an explicit answer when
the desired derivative orders are specified. Fix integers \(N,K,m\ge0\)
and put
\[
 J_*=N+K,\qquad q=\max(m,2),\qquad Q=q+K,
 \qquad R=R_{J_*,Q}.
\]
Define
\[
 d_0=\max_{0\le j\le J_*}\|V_j(0)\|_{H^Q},\qquad
 X=\sqrt R,\qquad Z=\sqrt{d_0^2+R}.
\]
The initial recurrence \eqref{r:eq:initial} includes the initial source
derivatives and makes \(d_0\) finite. Each temporal row in this selection
has a continuous \(H^Q\) representative, with
\begin{equation}
\begin{aligned}
 \sup_{0\le t\le T}\|V_j(t)\|_{H^Q}^2
 &\le\|V_j(0)\|_{H^Q}^2
 +2\|V_j\|_{L^2H^{Q+1}}\|V_{j+1}\|_{L^2H^{Q-1}}\\
 &\le d_0^2+R=Z^2.
\end{aligned}
\label{r:eq:material-target-trace}
\end{equation}
This follows by pairing the two adjacent Sobolev coordinates when
differentiating the squared \(H^Q\) norm. The Hilbert triple
\(H^{Q+1}\subset H^Q\subset H^{Q-1}\) also supplies continuity.

Every scalar component of \(D_t^ku\) is a sum of monomials
\[
 c\prod_{\ell=1}^{r}\partial_x^{\beta_\ell}V_{j_\ell},
 \quad 1\le r\le k+1,\quad
 \sum_\ell j_\ell=k+1-r,\quad
 \sum_\ell|\beta_\ell|=r-1.
\]
A temporal derivative increments one temporal index. Transport adds
one velocity factor and one spatial derivative. This proves the
displayed footprint by induction. After \(n\) ordinary temporal
derivatives, its temporal total is \(n+k+1-r\). The unique linear term
is \(V_{n+k}\); every nonlinear term has strictly lower temporal indices.

For integer \(q\ge2\), a scalar Sobolev algebra constant is
\[
 K_q=2^q(2\pi)^{-3/2}
 \left(\pi^{3/2}\frac{\Gamma(q-3/2)}{\Gamma(q)}\right)^{1/2}.
\]
Indeed, the weight inequality
\(\langle\xi\rangle^q\le2^{q-1}
(\langle\eta\rangle^q+\langle\xi-\eta\rangle^q)\), Fourier
convolution, and Cauchy--Schwarz with
\(\int\langle\xi\rangle^{-2q}d\xi
=\pi^{3/2}\Gamma(q-3/2)/\Gamma(q)\) give
\(\|vw\|_{H^q}\le K_q\|v\|_{H^q}\|w\|_{H^q}\).

For each component, the total absolute monomial coefficient in
\(D_t^ku\) is at most \(4^k k!\). A monomial with \(r\) factors
produces at most \(r\) temporal terms and \(3r\) transport terms
at the next material derivative. Applying \(n\) further temporal
derivatives costs at most \((k+1)^n\). Set
\[
 L_{n,k,m}=\sqrt3\,4^k k!(k+1)^n\max(1,K_q)^k.
\]
The factor \(\sqrt3\) combines the three scalar components. Since each
spatial factor has order at most \(K\), its \(H^q\) norm is bounded
by its parent's \(H^Q\) norm. Place one factor in \(L^2_tH^q_x\)
and all others in \(L^\infty_tH^q_x\). For every \(n\le N\),
\(k\le K\), this proves
\begin{equation}
\begin{aligned}
 \|\partial_t^nD_t^ku\|_{L^2(0,T;H^m)}
 &\le L_{n,k,m}X\max(1,Z)^k,\\
 \sup_{0\le t\le T}\|\partial_t^nD_t^ku(t)\|_{H^m}
 &\le L_{n,k,m}Z\max(1,Z)^k.
\end{aligned}
\label{r:eq:material-target-bounds}
\end{equation}
Continuity of the factors and of multiplication in \(H^q\) gives a
continuous \(H^m\) terminal trace. Replacing \(m\) by \(m+b\)
also controls spatial derivatives through order \(b\).
Every desired finite material order thus has a specified finite
rectangle within the full intersection. The constants may depend on
those requested orders.


\section{Conservative forcing and pressure adjustment}

A conservative force can be absorbed entirely into pressure.
Part~I supplies the global unforced history, so this reduction gives forced
continuation at arbitrary amplitude for conservative sources. Part~II treats
arbitrary prescribed smooth forcing directly on its own forced history.

\begin{proposition}[Conservative forcing]\label{r:prop:potential}
Suppose \(\Pp f=0\). Any smooth unforced history with datum \(u_0\)
gives a forced history with identical velocity and full mixed
pressure-gradient regularity on \([0,T]\).
\end{proposition}
\begin{proof}
Define the potential by
\[
 \widehat\phi(\xi,t)=
 -\frac{i\xi\cdot\widehat f(\xi,t)}{|\xi|^2}.
\]
Its gradient is \(f\). The inverse derivative at zero frequency is
legitimate in three dimensions for the stated source class. In fact,
for every integer \(m\ge0\),
\[
 \norm\phi_{H^m}^2
 \le \frac{2^{m-1}}{\pi^2}\norm f_1^2+\norm f_{H^m}^2.
\]
Split the Fourier integral at radius one, use
\(|\widehat f|\le(2\pi)^{-3/2}\|f\|_1\), and evaluate
\(\int_{|\xi|\le1}|\xi|^{-2}\dd\xi=4\pi\).
The high-frequency part is bounded by the displayed \(H^m\) norm.
Every temporal derivative satisfies this estimate as well.

If \(q\) is the unforced pressure, set \(p=q+\phi\).
The forced momentum equation follows directly, and all velocity
coordinates are those of the unforced history. The source potential
supplies the additional pressure jet.
\end{proof}

The restriction on the force is also necessary if the velocity is to be
retained. To see the precise obstruction, fix a smooth time and allow the
pressure to vary. For any pressure \(q\) with an admissible Sobolev gradient,
its unforced momentum residual is
\[
 \mathcal R_q:=u_t+(u\cdot\nabla)u-\nu\Delta u+\nabla q
 =f+\nabla(q-p).
\]
The solenoidal part is independent of that choice:
\[
 \Pp\mathcal R_q=\Pp f.
\]
Orthogonality of gradient and solenoidal fields in \(H^m\) gives
\begin{equation}
 \begin{aligned}
 \norm{\mathcal R_q}_{H^m}^2
 &=\norm{\Pp f}_{H^m}^2
   +\norm{(I-\Pp)f+\nabla(q-p)}_{H^m}^2,\\
 \inf_q\norm{\mathcal R_q}_{H^m}&=\norm{\Pp f}_{H^m}.
 \end{aligned}
 \label{r:eq:pressure-residual-distance}
\end{equation}
The infimum is attained at \(q=p-\phi\), where
\(\nabla\phi=(I-\Pp)f\). The Fourier formula above, applied to the
original \(f\), produces this potential. Thus pressure can convert the forced velocity
into an unforced solution exactly when \(\Pp f=0\). Subtracting a source
coordinate from a completed row does not change this residual of the
velocity itself.

Pressure normalization also matters in comparisons with periodic forcing
results. Tao's Proposition~1.7 equates unforced and forced periodic
regularity when nonperiodic pressure is permitted
\cite[pp.~28--29]{Tao2013}. Its accelerated Galilean transformation
does not preserve finite energy or \(H^1(\R^3)\), as stated on page~43.
That transformation is not the route used here: Part~II proves the compact
smoothly forced problem directly without importing an unforced velocity.

Proposition~\ref{r:prop:potential} also applies on a final interval if \(\Pp f\) vanishes
there: use the smooth velocity at its initial time as the datum for
the unforced continuation, and join by local uniqueness.
Vanishing of force derivatives at a single spatial point does not
meet this hypothesis. The next calculation distinguishes those two
properties in the actual localized source.

\section{The terminal rotational force in OpenAI's construction}

The localized velocity and pressure in \cite[Section~10]{openai-paper}
define the candidate source through the momentum equation. On the
regular exterior and cutoffs, these formulas determine its terminal
rotational component directly.

First set viscosity to one. In the paper's heat exterior, write
\(\tau=1-t\), \(A=1/2+h\), \(0<h<1/100\). Where the compactly
supported vector-potential correction is absent, its localized field is
\begin{equation}
 u=\chi(r,z)K(r,t)e_\theta,\qquad
 p=\chi(r,z)p_{\rm ext}(r,t),\qquad
 K=2^A\kappa r^{-2A}H(4\tau/r^2),
 \label{r:eq:exterior}
\end{equation}
with \(\kappa>0\). This amplitude is unrelated to our Sobolev constant
\(c_\infty\). The time cutoff equals one near \(t=1\).
The heat profile obeys
\[
 \partial_tK=(\partial_{rr}+r^{-1}\partial_r-r^{-2})K.
\]
The radial heat expansion in \cite[Appendix~A, (A.32)--(A.35)]{openai-paper}
gives, with \((b)_j\) denoting the rising factorial,
\begin{equation}
 K_j(r):=\partial_t^jK(r,1)
 =2^A\kappa\,4^j(h)_j(1+h)_j\,r^{-1-2h-2j}>0.
 \label{r:eq:heat-jets}
\end{equation}

The spatial cutoff satisfies \(0\le\chi\le1\), equals one for
\(r\le r_0/2,\ |z|\le z_0/2\), and has support compactly contained
in \(r<r_0,\ |z|<z_0\)
\cite[Section~10, (10.3)--(10.5)]{openai-paper}.
There is a connected rectangle
\[
 \Omega=(r_0/4,2r_0)\times(-\varepsilon,\varepsilon)
\]
lying in the heat exterior for all \(t\) sufficiently close to one.
To verify this region, use the source bound
\(q\le C_0(\tau+|z|^{1/D})\), \(D=1/2-h>0\), and
the exterior condition \(X=r^2/(2q)\ge X_{\rm ext}\).
Here \(q\) is the positive scale function used by that construction;
\(C_0,q_*,X_{\rm ext}\) are its fixed comparison and exterior constants.
Choose
\[
 q_b=\min\{q_*/2,r_0^2/(64X_{\rm ext})\},\qquad
 C_0(\eta+\varepsilon^{1/D})<q_b,\quad
 \varepsilon<z_0/2.
\]
Then \(X>2X_{\rm ext}\) on \(\Omega\) for \(1-\eta<t<1\),
as required by the exact exterior construction
\cite[Section~10, Step~4, p.~116]{openai-paper}.

The uncut swirl satisfies its azimuthal heat equation. Multiplication
by \(\chi\) creates spatial gradients at the edge of the cutoff.
Pressure has no azimuthal component, and axisymmetric pure swirl has
radial transport acceleration. The entire azimuthal residual is thus
\begin{equation}
 f_\theta=-K(\chi_{rr}+r^{-1}\chi_r+\chi_{zz})-2K_r\chi_r.
 \label{r:eq:cutoff-force}
\end{equation}
This is a direct momentum calculation. It identifies the force needed
to maintain the cutoff of a diffusing swirl.

\begin{proposition}[Nonzero projected terminal jets]\label{r:prop:terminal}
For the localized source defined above, each finite terminal temporal jet
\(F_j=\partial_t^jf(\cdot,1)\) has
\[
 \Pp F_j\not\equiv0,\qquad j=0,1,2,\ldots.
\]
This persists under the viscosity rescaling in the source paper.
\end{proposition}
\begin{proof}
Differentiate \eqref{r:eq:cutoff-force} and use
\(K_j'/K_j=-(1+2h+2j)/r\). On \(\Omega\),
\begin{equation}
 (F_j)_\theta=-K_j
 \left(\chi_{rr}+\chi_{zz}
             -\frac{1+4h+4j}{r}\chi_r\right).
 \label{r:eq:terminal-residual}
\end{equation}
If this component vanished throughout \(\Omega\), the cutoff would
solve the uniformly elliptic equation
\[
 \left(\partial_{rr}+\partial_{zz}
             -\frac{1+4h+4j}{r}\partial_r\right)\chi=0
\]
on that connected rectangle. Its coefficients are bounded there,
since \(r\ge r_0/4\). The cutoff attains its maximum one at the
interior point \((3r_0/8,0)\), and vanishes at the interior point
\((3r_0/2,0)\). The strong maximum principle forbids both properties.
Consequently \((F_j)_\theta\) is nonzero somewhere in \(\Omega\).

Choose a smooth compactly supported \(\psi(r,z)\) in a region where
this component has strict sign. The three-dimensional test field
\(w=\psi(r,z)e_\theta\) is smooth, supported away from the axis, and
divergence free. With its sign chosen appropriately,
\[
 \ip{\Pp F_j}{w}=\ip{F_j}{w}
 =2\pi\int (F_j)_\theta\psi(r,z)\,r\dd r\dd z>0.
\]
This proves the projected nonvanishing. The rescaling
\(f_\nu(x,t)=\sqrt\nu\,f(x/\sqrt\nu,t)\)
\cite[Section~10, (10.22)]{openai-paper}
commutes with the Leray projection and preserves that conclusion.
\end{proof}

The source paper makes its residual flat at the spatial origin at the
terminal time. Proposition~\ref{r:prop:terminal} locates nonzero terminal
force elsewhere. Smooth flatness at the proposed singular point is
compatible with a nonzero field in the surrounding cutoff region.
The exterior test also gives a quantitative obstruction to that pressure
conversion, without requiring convergence of the force on the whole space.
Use the field \(w\) from the proof at \(j=0\), and set
\[
 d:=2\pi\int (F_0)_\theta\psi(r,z)\,r\dd r\dd z>0.
\]
On the fixed support of \(w\), the heat formula gives
\(\ip{f(t)}w\to d\). Choose \(\eta_1>0\) so that this pairing is
at least \(d/2\) for \(1-\eta_1<t<1\). Since every gradient pairs to
zero with \(w\), Sobolev duality and
\eqref{r:eq:pressure-residual-distance} give, for each integer \(m\ge0\),
\begin{equation}
 \inf_q\norm{\mathcal R_q(t)}_{H^m}
 =\norm{\Pp f(t)}_{H^m}
 \ge\frac{d}{2\norm w_{H^{-m}}}>0,
 \qquad 1-\eta_1<t<1.
 \label{r:eq:actual-pressure-obstruction}
\end{equation}
The positive lower bound prevents any pressure adjustment from
turning this velocity into an unforced solution on a final interval.
It is an obstruction to the conservative-force reduction, not a sign
condition on the total work of the source.

\section{A direct bound from the critical source norm}
\label{r:app:critical-source}

The source can also be tested before velocity-intersection membership is
known. Begin with a classical forced history on its interval of existence
and retain the nonlinear current in the critical energy balance. Put
\[
 y=\norm u_{\dot H^{1/2}},\qquad
 z=\norm u_{\dot H^{3/2}},\qquad
 \beta=\norm{\Pp f}_{\dot H^{-1/2}},\qquad
 C_*=\frac1{\sqrt2\pi}.
\]
The product estimate used in \eqref{r:eq:critical-production}, applied
before substituting any rectangle bound, gives
\begin{equation}
 y y'+(\nu-C_*y)z^2\le\beta z.
 \label{r:eq:critical-source-energy}
\end{equation}
This calculation is on a strict classical interval. It assumes no bound
on a velocity coordinate at the terminal time.

As long as \(y<\nu/C_*\), completing the square in \(z\) yields
\[
 y y'\le\frac{\beta^2}{4(\nu-C_*y)}.
\]
Multiplication by the positive denominator makes the left side an exact
derivative. Define
\[
 \Psi(s)=2\nu s^2-\frac43C_*s^3,\qquad
 B(t)=\int_0^t\beta(s)^2\dd s.
\]
Then
\begin{equation}
 \Psi(y(t))\le\Psi(y(0))+B(t),\qquad
 \Psi'(s)=4(\nu-C_*s)s.
 \label{r:eq:critical-source-primitive}
\end{equation}
The function \(\Psi\) is strictly increasing between zero and
\(\nu/C_*\). At the upper endpoint it takes the value
\[
 \Psi(\nu/C_*)=\frac{2\nu^3}{3C_*^2}
 =\frac{4\pi^2}{3}\nu^3.
\]

\begin{theorem}[Critical source criterion]\label{r:thm:critical-source}
Under the datum and source assumptions of this paper, suppose
\(y_0=\norm{u_0}_{\dot H^{1/2}}<\nu/C_*\) and
\begin{equation}
 B(T)<\frac{4\pi^2}{3}\nu^3-\Psi(y_0).
 \label{r:eq:critical-source-condition}
\end{equation}
The forced history extends smoothly through \(T\) and belongs to
\(\mathscr I_T\). In particular, for rest data the sufficient condition is
\(B(T)<4\pi^2\nu^3/3\).
\end{theorem}

\begin{proof}
Let \(r<\nu/C_*\) be the unique nonnegative number satisfying
\(\Psi(r)=\Psi(y_0)+B(T)\). If \(y\) first exceeded \(r\) before
reaching \(\nu/C_*\), \eqref{r:eq:critical-source-primitive} would fail.
Continuity thus gives \(y\le r\) throughout classical existence up
to \(T\). At a zero of \(y\), the energy identity for \(y^2\) and
absolute continuity justify the integrated inequality without dividing
by \(y\).

Write \(a=\nu-C_*r>0\). Young's inequality in
\eqref{r:eq:critical-source-energy} now gives
\[
 y y'+\frac a2 z^2\le\frac{\beta^2}{2a},\qquad
 \int_0^S z^2\dd t\le\frac{y_0^2}{a}+\frac{B(T)}{a^2}
 \quad(S<T).
\]
These are bounds obtained from the source and initial datum.
To see directly how they prevent an endpoint loss, put
\(Y=\norm{\nabla u}_2\) and \(D=\norm{\Delta u}_2\).
Differentiation of the momentum equation, integration by parts, and
incompressibility give
\[
 \frac12(Y^2)'+\nu D^2
 \le\norm{\nabla u}_3\norm{\nabla u}_2
       \norm{\nabla u}_6+\norm{\Pp f}_2D.
\]
Let \(C_S\) be a finite constant in
\(\norm{\nabla u}_6\le C_S\norm{\Delta u}_2\), and let
\(C_3=2^{-1/6}\pi^{-1/3}\) as above. Since
\(\norm{\nabla u}_3\le C_3z\), two applications of Young's inequality
give
\[
 (Y^2)'+\nu D^2
 \le\frac{2(C_SC_3)^2}{\nu}z^2Y^2
       +\frac2\nu\norm{\Pp f}_2^2.
\]
Gronwall's inequality bounds \(Y\) uniformly. The ordinary energy
identity also gives
\(\norm{u(t)}_2\le\norm{u_0}_2+
\int_0^t\norm{\Pp f(s)}_2\dd s\).
Thus the full \(H^1\) norm remains bounded.

For completeness, this bound supplies a uniform local continuation time.
The estimate
\(\norm{(v\cdot\nabla)w}_{3/2}
\le C\norm v_{H^1}\norm w_{H^1}\), the boundedness of \(\Pp\)
on \(L^{3/2}\), and the heat estimates from \(L^{3/2}\) to \(H^1\)
bound the bilinear mild map on a time interval of length \(h\) by
\[
 C_\nu(h^{1/4}+h^{3/4})
 \norm v_{C_tH^1}\norm w_{C_tH^1}.
\]
The two powers integrate the heat factors \(t^{-3/4}\) and
\(t^{-1/4}\). A contraction ball of fixed radius now has a fixed
positive existence time, using also the bounded \(H^1\) norm of the
source. Restarting near a proposed finite endpoint continues the
history beyond it. Heat smoothing at positive elapsed time and the
smooth source preserve every finite spatial order; differentiation
of the equation then preserves the full temporal jet. Together with
the smooth initial datum and source, this proves \(\mathscr I_T\)
membership and the endpoint conclusions.
\end{proof}

The threshold is unchanged by the Navier--Stokes rescaling
\(u_\lambda(x,t)=\lambda u(\lambda x,\lambda^2t)\),
\(f_\lambda(x,t)=\lambda^3f(\lambda x,\lambda^2t)\).
The squared \(\dot H^{-1/2}\) source norm gains a factor
\(\lambda^2\), which its time integral cancels.

\subsection{A necessary \texorpdfstring{\(H^3\)}{H3} source budget}

For a candidate that starts from rest, write
\[
 Y(t)=\norm{u(t)}_{H^3},\qquad
 D(t)=\norm{\nabla u(t)}_{H^3},\qquad
 g(t)=\norm{\Pp f(t)}_{H^3},\qquad
 G_T=\int_0^Tg(t)\dd t.
\]
On each strict classical interval, the momentum equation and
\eqref{r:eq:product} give
\[
 YY'+\nu D^2\le8c_\infty Y^2D+gY,
 \qquad c_\infty=(8\pi)^{-1/2}.
\]
Completing the entire viscous square yields
\[
 8c_\infty Y^2D-\nu D^2
 \le\frac{16c_\infty^2}{\nu}Y^4
 =\frac2{\pi\nu}Y^4.
\]
Thus \(Y'\le2Y^3/(\pi\nu)+g\) wherever \(Y>0\).
At zeros, use \(Y_\varepsilon=(Y^2+\varepsilon^2)^{1/2}\);
its derivative obeys
\(Y_\varepsilon'\le2Y_\varepsilon^3/(\pi\nu)+g\).

If \(G_T=\infty\), the finite source condition already fails. Otherwise
put
\[
 Z_\varepsilon(t)=Y_\varepsilon(t)+\int_t^Tg(s)\dd s.
\]
Since \(Y_\varepsilon\le Z_\varepsilon\),
\[
 Z_\varepsilon'\le\frac2{\pi\nu}Z_\varepsilon^3,
 \qquad Z_\varepsilon(0)=G_T+\varepsilon.
\]
Integrating this scalar comparison and taking \(\varepsilon\downarrow0\)
gives
\begin{equation}
 Y(t)\le
 \frac{G_T}{\sqrt{1-4G_T^2t/(\pi\nu)}}
 \quad\text{if}\quad
 G_T<\sqrt{\frac{\pi\nu}{4T}}.
 \label{r:eq:direct-H3-budget}
\end{equation}
This bound uses the finite source budget and strict-interval equations;
it assumes neither a terminal velocity norm nor smooth continuation of
force through \(T\).

The calculated exterior has an unbounded \(L^3\) norm as \(t\uparrow1\).
Since \(H^3\) controls \(\dot H^{1/2}\), which embeds into \(L^3\),
its \(H^3\) norm is unbounded as well. Equation~\eqref{r:eq:direct-H3-budget}
thus proves the following lower bound for its actual residual:
\begin{equation}
 G_{\rm OA}\ge\frac{\sqrt{\pi\nu}}2.
 \label{r:eq:actual-H3-budget-failure}
\end{equation}

\subsection{Strict failure of the critical source condition}

For the prescribed exterior, cylindrical integration on
\(\mathcal E_\tau\) gives an explicit critical lower bound:
\[
 \int_{\mathcal E_\tau}|u|^3\dd x
 =C_L\tau^{-4h},\qquad
 C_L=4\pi d\int_c^{2c}k(y)^3y\dd y>0.
\]
The Sobolev inequality implies
\[
 \norm u_{\dot H^{1/2}}
 \ge C_3^{-1}C_L^{1/3}\tau^{-4h/3}\longrightarrow\infty.
\]
The viscosity rescaling preserves this divergence. Continuity gives a first
time \(S<1\) at which
\(y(S)=\nu/C_*\), with \(y(t)<\nu/C_*\) for \(0\le t<S\).
Integrate \eqref{r:eq:critical-source-primitive} before \(S\), then pass to
\(S\) by continuity. The rest datum gives
\begin{equation}
 B(S)\ge\Psi(\nu/C_*)=\frac{4\pi^2\nu^3}{3}.
 \label{r:eq:necessary-source-budget}
\end{equation}
The bound follows from the explicit exterior and momentum equation on
strict subintervals of \([0,1)\).

The residual continues to contribute a positive source norm after this
crossing. Use the divergence-free test field \(w\) from the fixed
exterior calculation leading to \eqref{r:eq:actual-pressure-obstruction},
and write \(\gamma_w:=\lim_{t\uparrow1}\ip{f(t)}w>0\).
For some \(\eta_1>0\), local convergence on its support gives
\(\ip{f(t)}w\ge\gamma_w/2\) for \(1-\eta_1<t<1\), independently
of terminal convergence elsewhere. Homogeneous Sobolev duality gives
\[
 \frac{\gamma_w}{2}\le\ip{\Pp f(t)}w
 \le\norm{\Pp f(t)}_{\dot H^{-1/2}}\norm w_{\dot H^{1/2}}.
\]
The test norm is finite and positive because \(w\) is smooth, compactly
supported, and nonzero. Let \(a=\max\{S,1-\eta_1\}<1\).
Integrating over \((a,1)\), which is disjoint from \((0,S)\), yields
\begin{equation}
 \begin{aligned}
 B(1)&\ge\frac{4\pi^2\nu^3}{3}
       +\frac{(1-a)\gamma_w^2}{4\norm w_{\dot H^{1/2}}^2}\\
     &>\frac{4\pi^2\nu^3}{3}.
 \end{aligned}
 \label{r:eq:strict-critical-budget-failure}
\end{equation}
For general viscosity, use the rescaled exterior and test field; the
source budget and positive correction both scale by \(\nu^3\).
The calculation also permits \(B(1)=\infty\). Thus no global terminal
source regularity is needed to establish failure of this sufficient
criterion. The source may be smooth and have finite norms, but these
particular quantitative estimates do not cover its magnitude. This failure
does not restrict the global theorem of Part~II.


\clearpage
\section*{Final conclusion}
\addcontentsline{toc}{section}{Final conclusion}

The datum-generated compatible intersection proves global smoothness and
uniqueness first for the unforced equations and then, on a separate history
carrying its complete source jet, for the forced equations.  The localized
residual of the field proposed in OpenAI's \emph{Finite Time Blowup for
Navier--Stokes} belongs to the forced theorem's source class.  The theorem
produces a global smooth solution for those exact prescribed data, while
strict-subinterval uniqueness identifies any classical candidate with that
solution on \([0,1)\).  Its finite velocity and mixed-rectangle bounds are
incompatible with the proposed field's unbounded velocity and with the
divergent positive exterior contribution to \(R_{0,1}\) calculated above.  The
proposed field does not
solve the stated forced initial-value problem.
